\documentclass[10pt,a4paper]{article}

% ----------------- Paquetes -------------------%
\usepackage[T1]{fontenc}
\usepackage[utf8]{inputenc}
\usepackage[a4paper,margin=2.5cm]{geometry}
\usepackage{mathptmx}             
\usepackage{changepage} 
\usepackage{setspace}
\usepackage[spanish]{babel}
\usepackage[labelfont=bf,labelsep=space]{caption}
\usepackage{amssymb}
\usepackage{amsmath}
\usepackage{ebgaramond}
\usepackage{placeins}
\usepackage{etoolbox}
\usepackage{setspace}
\usepackage{parskip} 
\usepackage{tcolorbox}
\usepackage{needspace}
\usepackage{xparse}
\usepackage{ocgx2}
\usepackage{hyperref}
\usepackage{hypcap}


%%%%%%%%%%%%%%%%%%%%%%%%%%%%%%%%%%%%%%%%%%%%%%%%%%%%%%%%%%%%%%%%
% --- Fija primero tus valores globales "buenos" ---
\setstretch{1.2}                 % Interlineado global
\setlength{\parskip}{0.7\baselineskip}
\setlength{\parindent}{0pt}
\newlength{\GlobalParskip}
\setlength{\GlobalParskip}{\parskip}

\newcounter{eqnum}[subsection]
\renewcommand{\theeqnum}{\arabic{eqnum}}
\newcounter{alphaitem}
\newcounter{numitem}

%% ------------------- Recuadros----------------------%%
\newcommand{\puntosclave}[1]{%
  \begin{tcolorbox}[
    colframe=black,
    title={Puntos clave},
    before upper={\setlength{\parindent}{0.5cm}\setlength{\parskip}{0.5\baselineskip}}
  ]
  #1
  \end{tcolorbox}
}

\newcommand{\modificaciones}[1]{
  \begin{tcolorbox}[
    colback=white,
    colframe=red,
    title={Modificaciones a realizar},
    before upper={\setlength{\parindent}{0.5cm}\setlength{\parskip}{0.5\baselineskip}}
  ]
  #1
  \end{tcolorbox}
}

%% ------------------- Listas----------------------%%
    % --- Lista alfabética ---
    \newenvironment{la}
      {\begin{list}{\alph{alphaitem})}{%
          \usecounter{alphaitem}%
          \setlength{\leftmargin}{1cm}%
          \setlength{\parskip}{3pt}% 
          \setlength{\itemsep}{0pt}% ← espacio entre ítems
          \setlength{\parsep}{0pt}%  ← párrafos dentro de un ítem
      }}
      {\end{list}}
      
    % --- Lista numérica ---
    \newenvironment{lnum}
      {\begin{list}{\arabic{numitem}.}{%
          \usecounter{numitem}%
          \setlength{\leftmargin}{1cm}%
          \setlength{\parskip}{3pt}% 
          \setlength{\itemsep}{0pt}% ← espacio entre ítems
          \setlength{\parsep}{0pt}%  ← párrafos dentro de un ítem
      }}
      {\end{list}}
      
% ------------------- Imágenes----------------------%
    \graphicspath{{Img/}}% Rutas por defecto 
    % --- Imagen adaptativa ---
    % Registros de dimensión definidos una sola vez
    \newdimen\imgcap
    \newdimen\imgmin
    \newdimen\imgmax
    \newdimen\imgremain
    \newdimen\imgh
    \newdimen\imgneed
    
    \newcommand{\imagenfit}[3]{%
      \addvspace{10pt}% separación antes
      \par\begingroup
        % Parámetros de composición (asignaciones, no \newdimen)
        \imgcap=1\baselineskip            % alto estimado de título+fuente
        \imgmin=0.15\textheight
        \imgmax=0.15\textheight
        \imgremain=\pagegoal
        \ifdim\pagegoal=\maxdimen \imgremain=0pt\fi
        \advance\imgremain by -\pagetotal
        \advance\imgremain by -\imgcap
        % Decisión de altura destino
        \ifdim\imgremain<\imgmin
          \newpage
          \imgh=0.20\textheight
        \else
          \imgh=\imgremain
          \ifdim\imgh>\imgmax \imgh=\imgmax \fi
        \fi
        % Reservar espacio para evitar cortes del bloque completo
        \imgneed=\imgh \advance\imgneed by \imgcap
        \Needspace{\imgneed}
        % Cuerpo
        \noindent\begin{minipage}{\textwidth}
          \centering
          \captionof{figure}{\textemdash\ #2}\par\vspace{0.3em}% título arriba
          \includegraphics[width=\linewidth,height=\imgh,keepaspectratio]{\detokenize{#1}}\par
          \vspace{0.3em}\small\textit{Fuente: }#3% fuente abajo
        \end{minipage}%
      \endgroup
      \par\addvspace{10pt}%
    }

%------------ Bloque de RECUADRO ESPECIAL -------------%
    \newenvironment{specialblock}[1]{%
      \begin{quote} % crea sangría a izquierda y derecha
      \small % texto más pequeño
      \noindent\textbf{#1}\\[-0.3em] % título en negrita
      \rule{\linewidth}{0.4pt}\\[0.6em]}{
      \end{quote}}
      

%-------------------------- Author Font -----------------------%
    \newcommand{\authorfont}[1]{{\fontfamily{EBGaramond-TLF}\selectfont #1}}

%-------------------------- Anexos -----------------------%
    \makeatletter
    \newcounter{anexo}
    \renewcommand{\theanexo}{\Roman{anexo}}
    \newcommand{\anexo}[1]{%
      \clearpage
      \phantomsection
      \refstepcounter{anexo}%
      \noindent\textbf{Anexo \theanexo.- }#1\par\medskip
      \addcontentsline{stc}{section}{Anexo \theanexo.- #1}%
    }
    \makeatother

    %-------------------------- Lista de figuras (personal) -----------------------%
\makeatletter
\newcommand{\MyListOfFigures}{%
  \clearpage
  \phantomsection
  {\centering\bfseries\Large Índice de figuras\par}%
  \medskip
  % Entrada en nuestro índice de contenido (stc)
  \addcontentsline{stc}{section}{Índice de figuras}%
  % Recuperación del listado (sin cabecera propia)
  \begingroup
    \parindent=0pt\parskip=2pt
    \@starttoc{lof}%
  \endgroup
}
\makeatother

%-------------------------- Bibliografía -----------------------%
\makeatletter
% Contador para las referencias
\newcounter{myref}
% Cabecera de la bibliografía, entra en el índice 'stc'
\newcommand{\MyBibliography}{%
  \clearpage
  \phantomsection
  {\centering\bfseries\Large Bibliografía y referencias\par}%
  \medskip
  \addcontentsline{stc}{section}{Bibliografía y referencias}%
  \setcounter{myref}{0}%
}
\newcommand{\MyBibItem}[4]{%
  \refstepcounter{myref}%
  \noindent[\themyref]\ #1.\ \emph{#2}. #3. #4\par
}
\makeatother

%--------- Establecer cuatro niveles de índice --------%
    \setcounter{tocdepth}{4}   
    \setcounter{secnumdepth}{4}% numera hasta \paragraph

%%%%%%%%%%%%%%%%%%%%%%%%%%%%%%%%

\makeatletter

    %--------- Interlineado Global --------%
    \edef\GlobalBaseLineStretch{\baselinestretch}
    
    %--------- Espaciado de seccinones --------%
    \renewcommand\section{\@startsection{section}
      {1}{\z@}
      {2.5ex \@plus 1ex \@minus .2ex} % espacio antes
      {1.5ex \@plus .2ex}             % espacio después
      {\normalfont\Large\bfseries}    % formato del título
    }
    \renewcommand\subsection{\@startsection{subsection}{2}{\z@}
      {3ex \@plus 0.8ex \@minus .2ex}
      {1.5ex \@plus .2ex}
      {\normalfont\large\bfseries}
    }
    \renewcommand\subsubsection{\@startsection{subsubsection}{3}{\z@}
      {3ex \@plus 0.5ex \@minus .2ex}
      {1ex \@plus .2ex}
      {\normalfont\normalsize\bfseries}
    }
    \renewcommand\paragraph{\@startsection{paragraph}{4}{\z@}%
    {-2ex \@plus -0.5ex \@minus -.2ex}% espacio antes
    {0.8ex \@plus .2ex}% espacio después
    {\normalfont\normalsize\itshape}}% ← cursiva en lugar de negrita

    %--------- Ecuaciones --------%   
    % Variables globales para almacenar títulos y notas
    \gdef\leq@current@section{}%
    \gdef\leq@current@subsection{}%
    \gdef\leq@last@printed@section{}%
    \gdef\leq@last@printed@subsection{}%
    
    % Comando para añadir nota (invisible en el cuerpo)
    \newcommand{\eqnote}[1]{%
      \addcontentsline{leq}{leqnote}{#1}%
    }
    
    % Comando para ecuaciones
    \newcommand{\eqblock}[2]{%
      % Verificar si necesitamos imprimir el título de sección
      \ifx\leq@current@section\leq@last@printed@section\else
        \addcontentsline{leq}{leqsection}{\leq@current@section}%
        \global\let\leq@last@printed@section\leq@current@section
        \gdef\leq@last@printed@subsection{}% Reset subsección al cambiar sección
      \fi
      % Verificar si necesitamos imprimir el título de subsección
      \ifx\leq@current@subsection\leq@last@printed@subsection\else
        \ifx\leq@current@subsection\@empty\else
          \addcontentsline{leq}{leqsubsection}{\leq@current@subsection}%
          \global\let\leq@last@printed@subsection\leq@current@subsection
        \fi
      \fi
      % Ecuación
      \refstepcounter{eqnum}%
      \begin{equation*}
        #1\tag{E.\theeqnum}
      \end{equation*}%
      \addcontentsline{leq}{leqt}{[E.\theeqnum] -- #2}%
      \addcontentsline{leq}{leqm}{#1}%
    }
    
    %%% ----- Índice de ecuaciones ----------%%%
    % Tamaño para []
    \newcommand{\leqIndexFont}{\normalsize}
    \newcommand{\leqTitleFont}{\tiny}
    \newcommand{\leqNoteFont}{\tiny}
    % Cabecera de sección 
    \newcommand*\l@leqsection[2]{%
      \addvspace{25pt}% Mayor separación antes de nueva sección
      \noindent\leqIndexFont
      \makebox[\linewidth]{\rule{.38\linewidth}{0.4pt}\ \ \textbf{  \MakeUppercase{#1}}\ \ \rule{.38\linewidth}{0.4pt}}%
      \par\vspace{-10pt}
    }
    % Cabecera de subsección 
    \newcommand*\l@leqsubsection[2]{%
      \par\addvspace{15pt}%
      \noindent\leqIndexFont #1
      \par\vspace{-7pt}
    }
    % Título de ecuación 
    \newcommand*\l@leqt[2]{%
    \par\addvspace{10pt}%
      \par\noindent\leqTitleFont #1\par
    }
    % Miniatura de ecuación
    \newcommand*\l@leqm[2]{%
      \vspace{0.3\baselineskip}%
      \par\scriptsize\noindent\hspace*{1.5em}\(#1\)\par%
    }
    % Nota de ecuación 
    \newcommand*\l@leqnote[2]{%
      \vspace{0.2\baselineskip}%
      \par\noindent\leqNoteFont\hspace*{1.5em}\textbf{Nota.--} #1\par\vspace{0.5\baselineskip}%
    }
    % Generación del índice
    \newcommand{\mylistofequations}{%
      \clearpage
      \phantomsection
      % Título visual de la lista (ajústalo a tu gusto):
      {\centering\bfseries\Large Índice de ecuaciones\par}
      % Entrada en nuestro índice 'stc':
      \addcontentsline{stc}{section}{Índice de ecuaciones}%
      % Llamada a tu lista real (usa el nombre exacto de tu macro):
      \@starttoc{leq}%
    }
    
    %--------- Captura de títulos de sección --------%
    \let\leq@original@section\section
    \let\leq@original@subsection\subsection
    % Flag para saber si estamos en una sección asterisco
    \newif\ifLeqStarSection
    \LeqStarSectionfalse
    
    % Redefinir \section CON CONTROL DE ASTERISCO
    \renewcommand{\section}{%
      \@ifstar{\leq@section@star}{\@dblarg\leq@section@nostar}%
    }
    \def\leq@section@star#1{%
      \global\LeqStarSectiontrue
      \gdef\leq@current@section{#1}%
      \gdef\leq@current@subsection{}%
      \leq@original@section*{#1}%
      \global\LeqStarSectionfalse
    }
    \def\leq@section@nostar[#1]#2{%
      \global\LeqStarSectionfalse
      \gdef\leq@current@section{#2}%
      \gdef\leq@current@subsection{}%
      \leq@original@section[#1]{#2}%
    }
    % Redefinir \subsection
    \renewcommand{\subsection}{%
      \@ifstar{\leq@subsection@star}{\@dblarg\leq@subsection@nostar}%
    }
    \def\leq@subsection@star#1{%
      \global\LeqStarSectiontrue
      \gdef\leq@current@subsection{#1}%
      \leq@original@subsection*{#1}%
      \global\LeqStarSectionfalse
    }
    \def\leq@subsection@nostar[#1]#2{%
      \global\LeqStarSectionfalse
      \gdef\leq@current@subsection{#2}%
      \leq@original@subsection[#1]{#2}%
    }

    % ----- Restauración de valores Globales de estilo ----------%
    % Flags
    \newif\ifInFootnote
    \InFootnotefalse
    \pretocmd{\@footnotetext}{\InFootnotetrue}{}{}
    \apptocmd{\@footnotetext}{\InFootnotefalse}{}{}
    % Profundidad de listas (soporta anidadas)
    \newcount\MyListDepth \MyListDepth=0
    \AtBeginEnvironment{la}{\advance\MyListDepth by 1}
    \AfterEndEnvironment{la}{\advance\MyListDepth by -1}
    \AtBeginEnvironment{lnum}{\advance\MyListDepth by 1}
    \AfterEndEnvironment{lnum}{\advance\MyListDepth by -1}
    \AtBeginEnvironment{itemize}{\advance\MyListDepth by 1}
    \AfterEndEnvironment{itemize}{\advance\MyListDepth by -1}
    \AtBeginEnvironment{enumerate}{\advance\MyListDepth by 1}
    \AfterEndEnvironment{enumerate}{\advance\MyListDepth by -1}
    % Restauración diferida: hazlo al INICIO del próximo párrafo real
    \newif\if@pendingRestore
    \@pendingRestorefalse
    \newcommand{\RequestRestore}{\global\@pendingRestoretrue}
    \everypar{%
      \if@pendingRestore
        \global\@pendingRestorefalse
        \ifInFootnote\else
          \ifnum\MyListDepth=0
            \setlength{\parskip}{\GlobalParskip}%
            \setstretch{\GlobalBaseLineStretch}%
          \fi
        \fi
      \fi
    }
    % Al final de bloques:
    \AfterEndEnvironment{equation}{\RequestRestore}
    \AfterEndEnvironment{equation*}{\RequestRestore}
    \AfterEndEnvironment{la}{\RequestRestore}
    \AfterEndEnvironment{lnum}{\RequestRestore}
    % Al final de macros:
    \apptocmd{\eqblock}{\RequestRestore}{}{}
    \apptocmd{\imagenfit}{\RequestRestore}{}{}
    
\makeatother

% ===================== ToC propio con 4 niveles (único bloque) =====================
\makeatletter

% --- Escritores: añadimos una línea al .stc en cada cabecera numerada ---
% Guardamos las versiones actuales para llamar al comportamiento normal
\let\MyTOC@orig@section\section
\let\MyTOC@orig@subsection\subsection
\let\MyTOC@orig@subsubsection\subsubsection
\let\MyTOC@orig@paragraph\paragraph

% SECTION
\renewcommand{\section}{%
  \@ifstar{\MyTOC@section@star}{\@dblarg\MyTOC@section@nostar}%
}
\def\MyTOC@section@star#1{%
  \MyTOC@orig@section*{#1}% sin entrada al .stc
}
\def\MyTOC@section@nostar[#1]#2{%
  \MyTOC@orig@section[{#1}]{#2}% comportamiento normal (escribe al .toc)
  % Escribimos también nuestra línea al .stc (texto con \numberline)
  \addcontentsline{stc}{section}{\protect\numberline{\thesection}#2}%
}

% SUBSECTION
\renewcommand{\subsection}{%
  \@ifstar{\MyTOC@subsection@star}{\@dblarg\MyTOC@subsection@nostar}%
}
\def\MyTOC@subsection@star#1{%
  \MyTOC@orig@subsection*{#1}%
}
\def\MyTOC@subsection@nostar[#1]#2{%
  \MyTOC@orig@subsection[{#1}]{#2}%
  \addcontentsline{stc}{subsection}{\protect\numberline{\thesubsection}#2}%
}

% SUBSUBSECTION
\renewcommand{\subsubsection}{%
  \@ifstar{\MyTOC@subsubsection@star}{\@dblarg\MyTOC@subsubsection@nostar}%
}
\def\MyTOC@subsubsection@star#1{%
  \MyTOC@orig@subsubsection*{#1}%
}
\def\MyTOC@subsubsection@nostar[#1]#2{%
  \MyTOC@orig@subsubsection[{#1}]{#2}%
  \addcontentsline{stc}{subsubsection}{\protect\numberline{\thesubsubsection}#2}%
}

% PARAGRAPH
\renewcommand{\paragraph}{%
  \@ifstar{\MyTOC@paragraph@star}{\@dblarg\MyTOC@paragraph@nostar}%
}
\def\MyTOC@paragraph@star#1{%
  \MyTOC@orig@paragraph*{#1}%
}
\def\MyTOC@paragraph@nostar[#1]#2{%
  \MyTOC@orig@paragraph[{#1}]{#2}%
  % Si no numeras los \paragraph, cambia \theparagraph por texto plano sin \numberline
  \addcontentsline{stc}{paragraph}{\protect\numberline{\theparagraph}#2}%
}

% --- Impresor del índice propio (lee el .stc) ---
\newcommand{\MyTOC}{%
  \par\bigskip
  {\centering\bfseries\Large Índice de contenido\par}%
  \medskip
  \begingroup
    \parindent=0pt\parskip=2pt
    \@starttoc{stc}%
  \endgroup
}

\makeatother
% ===================== fin ToC propio =====================


%%%%%%%%%%%%%%


% --- Datos del tema ---
\title{Sistema Fiscal en España (II): Impuesto de Sociedades\\
\large }
\author{Marcelo Ortega}
\date{Septiembre 2026}

\begin{document}
\maketitle
\newpage

% --- Página de resumen y modificaciones ---
\puntosclave{ %% Escribir puntos clave Apuntes CECO
     
    }
\vspace{1cm}

\modificaciones{
    
    }

\newpage
\MyTOC
\newpage

\mylistofequations
\clearpage

\MyListOfFigures
\clearpage


\pagenumbering{arabic}
\setcounter{page}{1}


\section{Introducción}



\subsection{Contextualización}


\subsubsection{Relevancia}




\subsubsection{Problemática}



\subsection{Estructura de la exposición}




\clearpage
\section{IS: Marco conceptual}
\puntosclave{
    El Impuesto sobre Sociedades se puede definir corno un tributo ele carácter directo y naturaleza personal que grava la renta de las sociedades y demás entidades jurídicas.
    
    Respecto del ámbito de aplicación el impuesto se exige en todo el territorio español, sin perjuicio de los regímenes tributarios forales de concierto y convenio económico en vigor, respectivamente. en los territorios históricos de la Comunidad del País Vasco y en la Comunidad Foral de Navarra.
    
    Las exenciones existentes son o bien objetivas o subjetivas. En el caso de las segundas, o plenas o parciales.
}

El IS. constituye la tercera fuente de ingresos impositiva estatal (aproximadamente 28.066 mili. de E' en 2020) y vemos como su recaudación es muy sensible al ciclo económico en el que se encuentre el País. A continuación, vemos la evolución de la recaudación de este Impuesto a lo largo de los últimos ai'ios:

\textbf{TABLA DE RECAUDACIÓN IS}

Como podernos observar. durante los años de recesión económica, 2008 en adelante, se produjo una caída notable de la recaudación del IS. A partir del 2012 la recaudación comienza a recuperarse debido a las medidas de política fiscal introducidas por el Gobierno para incrementar la recaudación, así corno, a partir de 201 3 . a la tímida mejora de la situación económica espai'iola que más tarde se va consolidando.

Los impuestos directos, como prácticamente único elemento de política fiscal sobre el que los Gobiernos de nuestro entorno tienen capacidad real de actuación al no encontrarse armonizados a nivel europeo. cobran especial importancia a la hora de marcar la política económica de un país.

En este sentido. en los últimos años se han producido en España cambios normativos. no sólo en el IRPF sino también en el IS. para tratar de intervenir en el ciclo económico y corregir o modular el déficit público.

La crisis económica que sufrió España a partir de 2007 elevó exponencialmente el déficit público y derivó en una desviación respecto del objetivo marcado por Europa de más de 2,5 puntos del Pll3 en el año 20 l 1 . Como consecuencia de dicha situación. una primera fase de la política fiscal de los últimos aí'íos consistió en incrementar la recaudación a través de una serie de medidas que afectó fundamentalmente al Impuesto sobre Sociedades de las grandes empresas a lo largo del año 201 2. tal y como veremos más adelante.

Una vez superada la crisis económica y con un déficit público más acomodado a los estándares europeos. se inició una reforma fiscal en 2014, aplicable a partir de 2015, que consistió fundamentalmente. en el ámbito del IS, en una reducción de la presión fiscal sobre las personas jurídicas a través de una bajada del tipo nominal del IS del 30 al 25% y una ampliación de bases imponibles que buscaban un Impuesto más generalista y más competitivo desde el punto de vista del derecho comparado.

Una vez situado en contexto tanto el peso cuantitativo respecto al total de ingresos tributarios como las funciones de cada uno de estos impuestos, IS e IRNR, p rocedemos a analizar cada uno de ellos por separado. estableciendo. en primer lugar. su definición y normas que los amparan. a continuación, una referencia a su evolución histórica y modificaciones recientes, para luego profundizar en el esquema del impuesto en sí y sus principales características.

El Impuesto sobre Sociedades se puede definir como un tributo de carácter directo y naturaleza personal que grava la renta ele las sociedades y demás entidades jurídicas. Su regulación básica la encontramos en la Ley 27/2014, de 27 de noviembre. del Impuesto sobre Sociedades.

\subsection{Historia: ¿?}
La primera Ley reguladora del Impuestos sobre Sociedades, fue la Ley 61/1978 como consecuencia de la Reforma Ordoñez. Dicha Ley establecía una lista tasada de gastos deducibles. como las adquisiciones corrientes de bienes y servicios. las amortizaciones, intereses. aquileres, · etc ... Fuera de esa lista, los gastos no eran deducibles. A esta primera ley, le sustituyó la Ley 43/1995, de 27 de diciembre, del Impuesto sobre Sociedades que establecilas reglas esenciales de la actual estructura del 1mpuesto sobre Sociedades, inspirada en los principios de neutralidad, transparencia. sistematización. coordinación internacional y competitividad. No osbstante, la Ley 43/1995, de 27 de diciembre. estableció como una de sus principales novedades la determinación de la base imponible del Impuesto ele manera sintética. a partir del resultado contable, corregido por las excepciones legalmente tipificadas. Es decir, la Ley 43/1978 cambiaba el concepto al pasar de una una lista cerrada de gastos decibles (la de la Ley 61/1978) a los admitidos contablemente salvo algunas excepciones.

Posterionnente, con la finalidad de incrementar la claridad del sistema tributario y mejorar la seguridad jurídica, se aprobó el Texto Refundido de la Ley del Impuesto sobre Sociedades, a través del Real Decreto Legislativo 4/2004, de 5 de marzo, que t11vo como objetivo fundamental integrar en un único cuerpo nomiativo todas las disposiciones que afectaban a este Impuesto. salvo casos excepcionales.

\subsubsection{Reciente: Modificaciones}
No obstante. desde la aprobación este Texto Refundido, éste ha sido o��jeto de constantes modificaciones de carácter parcial. Fundamentalmente a pa11ir de '.201 1 . el Gobierno ha sido muy activo en relación a su política fiscal como consecuencia de la situación macroeconómica que atravesaba Espaíia. por lo que se introducen diferentes modifícaciones en el IS. Vamos a destacar, a cominuación. las principales medidas introducidas en el ámbito del Impuesto. distinguiendo tres etapas·:

\textbf{1. Etapa de crisis ecouámica con elevado nil'el de d��ficit público: }

Esta primera etapa está marcada por la imperante necesidad de incremento de los ingresos públicos. por lo que se tomaron medidas fundamentalmente de índole recaudatoria, entre la s que podemos destacar las siguientes según el orden cronológico en el que fueron tomadas:

A) Marzo 2012: En este momento se optó por modificar el Impuesto sobre Sociedades con medidas que tuvieron por objeto aumentar el tipo efectivo del mismo. sin modificar el tipo de gravamen fo1mal. Estas medidas afectan fundamentalmente a quien dispone de una mayor capacidad de pago e incrementa más el tipo efectivo de las grandes empresas que el de las restantes. Entre ellas, podemos destacar las siguientes: se tijan límites para fa deducción de la libertad de amortización se establece un pago fraccionado mínimo general del 8 por ciento del resultado contable para entidades con importe neto de cifra de negocios superior a 20 millones de euros se limita la deducibíl idad de los gastos financieros en los grupos mercamiles

B) Julio 2012: Dado que los ingresos p úblicos no se recuperaban y siguiendo las recomendaciones de la Unión Europea. se introdujeron en el Impuesto sobre Sociedades: se limita la compensación de bases negativas se aumenta la base imponible de los pagos fraccionados con el 25 por ciento del importe de los dividendos y las rentas a los que resulte de aplicación la exención por doble imposición internacional se incrementan los tipos impositivos de los pagos fraccionados se eleva el tipo general del pago fraccionado mínimo del 8 al 12 por ciento
C) Diciembre 2012: En el ámbito del IS se adoptan las siguientes medidas: se establece un límite a la deducibilidad fiscal de las amortizaciones del inmovilizado material cifrado en un 70% del coeficiente máximo de amortización previsto en tablas, para las grandes empresas

\textbf{2. Etapa de reversión de la situación macroeconómica con leve crecimiento:}

En esta etapa. con carácter general. se establece un marco fiscal más favorable para el autónomo que inicia una a ctividad emprendedora con el objetivo de incentivar la creación de empresas y reducir la carga impositiva durante los primeros años de ejercicio de dicha actividad. Así, en marzo de 2013, en el ámbito del lS, con la finalidad de fomentar el inicio de la actividad emprendedora, se establece, para las entidades de nueva creación que realicen actividades económicas, un tipo de gravamen del 15% para los primeros 300.000 euros de base imponible. y del 20% para el exceso sobre dicho importe. aplicable el primer período impositivo en que la base i mponible de las entidades resulte positiva y en el período impositivo siguiente a este. En septiembre de 2013, se aprueba la Ley de Emprendedores (Ley 14/2013. de 2 7 de septiembre, de apoyo a los emprendedores y su internacionalización). que recoge. en lo relativo al IS, entre otras. las siguientes medidas: 

Deducción por inversión de beneficios: Con el objetivo de incentivar la capitalización empresarial y la inversión en activos nuevos tangibles y afectos a actividades económicas, y mejorar la situación patrimonial de las entidades de reducida dimensión frente a terceros por el incremento que experimentan los fondos propios. se aprueba esta medida. que supone que las empresas de reducida dimensión (importe neto de la cifra de negocios inferior a 1 O millones de euros) pueden aplicarse una deducción con carácter general del 10% en la cuota íntegra del I rn puesto sobre Sociedades de los beneficios obtenidos en el período impositivo, que se invirtieran en elememos nuevos del inmovilizado material o inversiones inmobiliarias afectos a la actividad económica 

I+D+i: Las deducciones por gastos e inversiones en l+D+i podrán opcionalmente ser objeto de aplicación sin quedar sometidas a ningún límite en la cuota y resultar abonadas, con un descuento del 20% de su importe, cuando no hayan podido aplicarse por insuficiencia de cuota. El importe de la deducción aplicada o abonada por las actividades de· investigación y desarrollo e innovación tecnológica no podrá superar conjuntamente. y por todos los conceptos, los 3 millones de euros anuales. Con l a reforma. las empresas ganan ·'seguridad" y '·certeza'· en cuanto que, en un momento u otro. se van a poder beneficiar de esta deducción y van a poder recuperar sus inversiones en l+D+i. 

Patent box (reducción de rentas procedentes de determinados activos intangibles): Con el objetivo de situar a España al nivel de los países de nuestro entorno que han implantado regímenes similares, como son Bélgica. Holanda, Luxemburgo. Reino Unido, Franc ia o Hungría, se aprueba esta medida, consistente en la integración en la base imponible del IS de las rentas procedentes de la cesión del uso o explotación de los activos en un 40% (es decir. se establece una reducción del 60% de las rentas). El incentivo también es aplicable a las rentas derivadas de la transmisión de los activos intangi bles. De esta fomrn, se establece un incentivo espaíiol competitivo con los incentivos análogos de otros países de la Unión Europea.

\textbf{3. Consolidación de la recuperación econámica:}

Ya en fase económica expansiva. se considera el momento adecuado para aprobar una reforma fiscal que afecte a los principales impuestos, fundamentalmente IR.PF e IS. Esta reforma tiene como final i dad devolver a los contribuyentes el esfuerzo contributivo realizado durante los años de crisis bajando los tipos de gravamen del JRPF y del IS. En el ámbito del IS se aprueba la nueva Ley 27/2014, de 27 de noviembre. del Impuesto sohre Sociedades (BOE 28 noviembre), que se analiza en el epígrafe siguiente.

ll.l.lll. ley 2 712014, del Impuesto sohre Sociedades
Todas estas modilicaciones parciales del Texto Refundido del IS, siendo todas ellas parcialmente consistentes. no han ido acompaiiadas de una revisión de carácter global hasta la Refomrn Fiscal de 2014 con la aprobación de la Ley 27/2014 del Impuesto sobre Sociedades. El objetivo fundamental de esta refonna es ampliar bases imponibles, simplificando y haciendo más equitativo el im puesto a cambio de rebajar el tipo impositivo nominal al 25%. Este objetivo global se puede desglosar de la siguiente fom1a: 
Busca alcanzar los principios de neutralidad, igualdad y justicia en el seno del IS 
Acercamiento entre tipo nominal y tipo efectivo. Una crítica tradicional del sistema anterior era la existencia de un tipo nominal muy alto en relación a la existencia de un tipo efectivo muy inferior. 
Am pliación de la base im ponible. Se introducen medidas que incrementarán la base del Impuesto. perm itiendo que la reducción del tipo de gravamen tenga un m enor impacto recaudatorio. 
Reducción de la diferencia e11tre el tratamiento fiscal de la financiación ajena y financiación propia. Se adoptan medidas para incentivar la capitalización de empresas y un menor sesgo al endeudamiento. 
Sim plificación del Impuesto. Se adoptan medidas que si mplifican la estructura del Impuesto.

Las medidas concretas en vigor a partir de esta refonna son, entre otras. las siguientes: Reducción de los tipos de gravamen:

Situación anterior:

Tras la refomrn:
► Tipo general: 30%
► Entidades de Reducida Dimensión (Cifra de negocios igual o inferior a 1 0mill€): 25% por los primeros 300.000€ y 30% del resto
► Microernpresas (Cifra de negocios igual o inferior a 5m ill€) con creación y mantenimiento de empleo: 20% por los primeros 300.000E y 25% del resto
► Entidades de nueva creación: 1 5% por lo�� primeros 300.000€ y 20% del resto (durante 2 años)
► M inoración del tipo de gravamen general en dos etapas. El tipo impositivo actual es del 25%.
► Mantenimiento del tipo del 30% para entidades financieras.
► M antenimiento del Lipo de Entidades de Reducida Dimensión al 25%
► Se suprime el tipo aplicable a microempresas.
► Entidades de nueva creación: 1 5% para toda la base imponible ( durnnte 2 ai'íos) 

Establecimiento de una reserva de capitalización em presarial y una reserva de nivelación para Pymes, como reducciones para la determinación de la base imponible. Se analizarán en profundidad más ade_lante. 
Establecimiento de la exención para evitar la doble im posición de dividendos y rentas derivadas de la transmisión de valores representativos de los fondos propios de entidades residentes y no residentes en territorio espai'íol, de tal manera que se iguala el tratamiento fiscal a las inversiones realizadas en entidades nacionales y extranjeras. 
Se establece como gasto no deducible el deterioro de elementos del inmovilizado material, inversiones inmobiliarias e inmovilizado intangible, de acciones y participaciones e n el capital y de valores representativos de deuda.

Situación Actual

En el ejercicio 2022 se introdujeron una serie de m odificaciones importantes: Tributación m ínima. Se introduce un nuevo artículo 30 bis. según para contribuyentes cuyo impo11e neto de la cifra de negocios sea al menos de 20 millones. la cuota líquida no podrá ser inferior al resultado de aplicar el 1 5 por ciento a la base imponible, Régimen de entidades en arrendamiento de viviendas: Se estableció una bonificación del 40 % para aquellas entidades que se dedicaran al arrendamiento de viviendas.

Derogación del apartado 6 del artículo 1 2 1 como consecuencia de la sentencia del Tribunal de .Justicia de la Unión Europea de 27 de enero de 2022 en el asunto C-788/19 que determinó que determinados aspectos del régimen jurídico asociado a la obligación de declaración de bienes y derechos en el extranjero (modelo 720) incunen en incumplimiento de la nonnativa europea.

La modificación normativa. supone, como luego se verá la derogación de la imprescriptibilidad (declarada contraria al derecho de la unión), para aquellos bienes que no fueron declarados en la referida declaración informativa.

El apartado 6 del artículo 1 2 1 de la Ley establecía: ·'6. En todo caso. se entenderá que han sido adquiridos con cargo a renta no declarada que se imputará al periodo impositivo más antiguo de entre los no prescritos susceptible de regularización, los bienes y derechos respecto de los que el contribuyente no hubiera cumplido en el plazo establecido al efecto la obligación de información a que se refiere la Disposición adicional decimoctava de la Ley General Tributaria.

No obstante. n o resultará de aplicación lo previsto en este apartado cuando el contribuyente acredite que los bienes y derechos cuya titularidad le corresponde han sido adquiridos con cargo a rentas declaradas o bien con cargo a rentas obtenidas en periodos impositivos respecto de los cuales no tuviese la condición de contribuyente de este Impuesto.''

El problema afíadido. era que dicha ganancia de patrimonio nojustificad.i imputable al último ejercicio no prescrito, llevaba aparejada una sanción específica del 1 50%. de forma tal que ante un bien que tuviera un valor. por ejemplo de I millón de euros no declarado, a un tipo marginal por ej. del 50 %. la cuota a ingresar sería de 500.000 euros, más la sanción del 150%. esto es 750.000 €. es decir un total de 1.250.000 C. Libertad de amortización para elementos que utilicen energías renovables. 


Principales novedades para 2023: .Tipo de gramen para Entidades con un I mporte Neto de la Cifra de Negocios inferior a 1 millón de euros al 23 %. Régimen Fiscal Especial de las Islas Baleares. Similar al REF Canario. introduciéndose la Reserva para inversiones en Baleares (RIB). Como nota más destable se establece una reducción a que se aplicará a las dotaciones que en cada período i mpositivo se hagan a la reserva para inversiones en las Islas Baleares. hasta el límite del 90 por ciemo de la parte de beneficio obtenido en el mismo período que no sea objeto de distribución. La inversión deberá realizarse en: La dquisición de elementos patrimoniales del inmovilizado material o intangible, de elementos patrimoniales que contribuyan a la mejora y protección del medio am biente en el territorio de las llles Balears. así como los gastos de investigación y desarrollo derivados de actividades de investigación, desarrollo e innovación tecnológica a que se refiere el artículo 3 5 . 1 y 2 de la Ley 27/20 1 4. de 27 de noviembre, del I m puesto sobre  Sociedades. Dichos elementos deberán estar situados o ser recibidos en el archipiélago balear, utilizados en el 1úism o, afectos y necesarios para el desarrollo de actividades económicas del contribuyente. salvo en el caso de los que contribuyan a la mejora y protección del medio ambiente en el territorio balear.

Entendiéndose situados en el territorio balear: Las aeronaves que, por su destino, contribuyan a mejorar las conexiones de las llles Balears. en los términos que reglamentariamente se determinen; Los buques con pabellón español y con puerto base en las Jlles Balears; Las redes de transporte y de comunicaciones que conecten el archipiélago balear con el exterior. por el tramo de la misma que se encuentre dentro del territorio de las Illes Balears y a la parte situada fuera del mismo que se utilice para conectar entre sí las distintas islas del archipiélago; Las aplicaciones inform áticas y los derechos de propiedad industrial, que vayan a aplicarse exclusivamente en procesos productivos o actividades comerciales que se desarrollen en el ámbito territorial balear, así como los derechos de propiedad intelectual que sean objeto de reproducción y distribución exclusivamente en el archipiélago balear; las concesiones administrativas de uso de bienes de dominio público radicados en las llles Balears; las concesiones administrativas de p restación de servicios públicos que se desarrollen exclusivamente en el archipiélago: Las concesiones administrativas de obra pública para la ejecución o explotación de infraestructuras públicas radicadas en las llles Balears.

La creación de puestos de trabajo relacionada de forma directa con las inversiones previstas en la letra A. que se produzca dentro de un período de seis meses a contar desde la fecha de entrada en funcionamiento de dicha inversión.

La suscripción de acciones o pai1icipaciones en el capital emitidas por sociedades como consecuencia de su constitución o ampliación de capital que desarrollen en el archipiélago su actividad, siempre que se cumplan una serie de requisitos.\footnote{
1 .- Que realicen las inversiones en la adquisición de elementos patrimoniales en las Islas Baleares o creación de puestos de trabajo.
2.- El plazo para la rcallizaciún de las in versiones serán 3 mios desdo: desde la f.:cha del devengo del impuesto correspondiente al ejercicio en el que el contribuyente que adquiere las acciones o las participaciones en su capital hubiera dorado la reserva regulada en esre apartado.
3.- Los elemento�� patrimoniales así adquiridos deberán mantenerst" en funcionamiento en las 1 lles Balears en los términos pre"isto, en este apartado.
4.- El importe del , alor de adquisición de las inversiones realizadas por la sociedad participada deberá alcan,ar. como mínimo, el importe desembolsado de las accione�� o participaciones adquiridas por el contribuyente.
}

Las inversiones realizadas por la sociedad participada no darán lugar a la aplicación de ningún otro beneficio fiscal.

Se crea un Régimen es pecial para empresas industriales. agrícolas, ganaderas ypesqueras:

Los contribuyentes del Impuesto sobre Sociedades aplicarán una bonificación del I O por ciento de la cuota íntegra correspondiente a los rendimientos derivados de la venta de bient:s corporales producidos en las llles Balears por ellos mismos, propios de actividades agrícolas, ganaderas. industriales y pesqueras, en este último caso en relación con las capturas efectuadas en su zona pesquera y acuícola. Se podrán beneficiar de esta bonificación las personas o entidades domiciliadas en las llles Balears o en otros territorios que se dediquen a la producción de tales bienes en el archipiélago, mediantesucursal o establecimiento pe1manente.
La bonificación se aplicará sobre la parte de la cuota íntegra que proporcionalmente corresponda a los rendimientos derivados de las actividades de producción sciialadas.
La aplicación de la bonificación en cada período impositivo requerirá que la plantilla media de la entidad en dicho período no sea inferior a la plantilla media correspondiente a los doce meses anteriores al inicio del primer período impositivo en que tenga efectos el régimen previsto en este apartado.
Cuando la entidad se haya constituido dentro del sefialado plazo anterior de doce meses se tendrá en cuenta la pla11tilla media que resulte de ese período.
La bonificación · se incrementará hasta el 25 por ciento en aquellos períodos impositivos en los que. además de cumplirse el requisito previsto en el número anterior, se haya producido un incremento de plantilla media no inferior a la unidad respecto de la plantilla media del período impositivo anterior y dicho incremento se mantenga durante. al menos. un plazo de tres años a partir de la fecha de finalización del período impositivo en el que se aplique esta bonificación incrementada.
Se modifica el Régimen Fiscal Especial de Canarias, estableciéndose la posibilidad de aplicar el tipo impositivo de la Zona Especial Canaria ( 4%) a las operaciones triangulares.
Libe11ad de amortización para los vehículos eléctricos.


\subsection{Ámbito de aplicación}
Respecto del ámbito de arlicación el impuesto se exige en todo el territorio espafiol. sin perjuicio de los regímenes tributarios forales de concierto y convenio económico en vigor. respectivamente. en los territorios históricos de la Comunidad del País Vasco y en la Comunidad Foral de Navarra. y de los tratados y convenios internacionales además del Régimen especial de Canarias por la insularidad. denominado Régimen Económico y Fiscal (REF).

Destacar así mismo, que en la Ley de Presupuestos Generales del Estado para 2023, como ya se ha puesto de manifiesto anterionnente, ha aprobado el Régimen Fiscal Especial Balear parar los ejercicios económicos entre el I de enero de 2023 y el 31 de diciembre de 2028. Similar en algunos aspectos al REF canario a excepción de la Zona Especial Canaria que es exclusiva de Canarias.

Se trata de un impuesto estatal cuya recaudación no se encuentra cedida a las Comunidades Autónomas. por lo que el 100% de su recaudación permanece en las arcas del Estado central a diferencia del IRPF o IVA cuya recaudación se encuentra cedida a las Comunidades Autónomas en un 50% o la mayor parte de los Impuestos Especiales en un 58%.

\subsection{Hecho impobible}
Constituye el hecho imponible la obtención de renta, cualquiera que fuere su fuente u origen. por el sujeto pasivo.

IJ.Jll.J. Rentas exentas

A la hora de estudiar las exenciones ha de diferenciarse entre exenciones objetivas y exenciones subjetivas.

Exenciones objetivas

- Exención sobre dividendos y rentas derivadas de la transmisión de valores representativos de los fondos propios de entidades residentes y no residentes en territorio espaliol. Se trata de evitar la doble imposición económica sobre dividendos y plusvalías ya procedan de fuente extranjera o de fuente doméstica.

- Exención de las rentas obtenidas en el extranjero a través de un establecimiento permanente cuando el mismo haya estado sujeto y no exento a un impuesto de naturaleza idéntica o análoga al IS con un tipo nominal de al menos un 10%. Estarán exentas igualmente las rentas derivadas de la transmisión de un EP cuando se cumpla el requisito anteriormente selialado.

- Exención del 50 % de la renta positiv a derivada de la transmisión de bienes inmuebles urbanos que tengan la condición de activo no corriente y que hubieran sido adquiridos a título oneroso a partir de la entrada en vigor del RDL l 8i2012, 12 de mayo de 2012, hasla el 3 1 de diciembre de 2012.

Exenciones subjetivas
Dentro de esta categoría la Ley considera exentos del I S a dos tipos de sujetos en función del alcance de la exención.


Exenciones plenas
El Estado, las Comunidades Autónomas. y las Entidades Locales,
Los Organismos autónomos del Estado y entidades autónomas de análogo carácter de las Comunidades Autónomas y de las Entidades locales,
El Banco de Espaiia. los Fondos de Garantía de Depósitos, y los Fondos de G arantía de Inversiones,
Las entidades públicas encargadas de la gestión de la Seguridad Social,
El Instituto de Espafia y las Reales Academias Oficiales integradas en el mismo y las instituciones de las Comunidades Autónomas con lengua oficial propia que tengan fines análogos a los de la Real Academia Espafiola,
Los restantes organismos públicos mencionados en las disposiciones adicionales 9 y l O de la Ley de Organización y Funcionamiento de la Administración General del Estado. asi como los entes públicos de análogo carácter de las Comunidades Autónomas y de las
Entidades Locales,
Las Agencias Estatales.
El Consejo Internacional de Supervisión Pública. 

Exenciones parciales

Las entidades sin ánimo de lucro a las que es de aplicación el régimen fiscal de las entidades sin fines lucrativos y de los incentivos fiscales al mecenazgo. Es aplicable a las fundaciones, asociaciones declaradas de utilidad pública, ONG, delegaciones de fundaciones extranjeras, federaciones deportivas y federaciones de entidades sin fines lucrativos que cumplan una serie de requisitos.

Estas entidades gozan de exención en el IS por los resultados obtenidos en el ejercicio de las actividades que constituyen su o��jeto social. Por lo que la base imponible está formada sólo por las rentas de explotaciones económicas no exentas. La base i mponible detern1inada de esta fonna es gravada al tipo de 10%.

Una serie de entidades recogidas dentro del régimen especial de entidades parcialmente exentas. Estas entidades son. entre otras. las entidades e instituciones sin ánimo de lucro no acogidas a la Ley 49/2002; las uniones. federaciones y confederaciones de cooperativas; los colegios profesionales, las asociaciones empresariales. las cámaras oficiales, los sindicatos de trabajadores; o los fondos de promoción de empleo.

Estas entidades gozan de exención respecto de las rentas que procedan de la realización de actividades que constituyan su objeto social o finalidad específica. las derivadas de adquisiciones y de transmisiones a título lucrativo, siempre que se realicen en cumplimiento de su objeto social o finalidad específica y las que se pongan de manifiesto en la transmisión onerosa de bienes afectos a la realización del objeto social cuando el total producto obtenido se dedique a nuevas inversiones relacionadas con el objeto social.

Por último, señalar un régimen tributario particular de los partidos políticos. Según este régimen. los partidos políticos gozan de exención en el IS por las rentas obtenidas para la financiación de las actividades que constituyen su objeto o finalidad específica. La exención comprende. entre otras. las cuotas satisfechas por sus afiliados. las subvenciones y las donaciones privadas. Por las rentas no exentas tributan al 25%.

\subsection{Sujeto pasivo}
En cuanto a la definición de sujeto pasivo. la Ley procede a la enumeración de la�� entidades que tributan por este impuesto sefialando las ��iguientes. siempre que tengan su residencia en territorio español:
las personas jurídicas. excepto las sociedades civiles que no tengan objeto men:antil
las sociedades agrarias de transformación
los fondos de inversión
las Uniones Temporales de Empresas
los fondos de capital riesgo
los fondos de pensiones.
los fondos de regulación del mercado hipotecario
los fondos de titulización
los fondos de garantía de inversiones,
las comunidades titulares de montes vecinales en mano común determinados Fondos de Activos Bancarios

Siendo gravados por la totalidad de la renta que obtengan, con independencia del lugar donde se hubiere producido y cualquiera que sea la residencia del pagador.

Tienen su residencia en España las entidades en las que concurran alguno de los siguientes requisitos:
que se constituyan conforme a las leyes españolas.
que tengan su domicilio social en territorio español,
o que tengan su sede de dirección efectiva en territorio español. entendiendo a estos efectos por tal cuando radique en él la dirección y control del conjunto de sus actividades.

Se presume que son residentes en España las entidades radicadas en territorios de nula o baja tributación o en un territorio calificado como paraíso fiscal cuando sus activos principales, directa o indirectamente, consistan en bienes situados o derechos que se cumplan o ejerciten en territorio español, o cuando su actividad principal se desarrolle en éste, salvo que dicha entidad acredite que su dirección y efectiva gestión tienen lugar en aquel país o territorio. así como que la constitución y operativa de la entidad responde a motivos económicos válidos y razones empresariales sustantivas distintas de la simple gestión de valores u otros activos.


\subsection{Periodo impositivo}
EL IS es un impuesto periódico en el que la realización del hecho im ponible tiene lugar en un período indefinido de tiempo, lo que ha obligado al legislador a dividir ese hecho en varios períodos determinando el último día de cada uno de ellos el nacimiento de la obligación tributaria.

El período impositivo coincide con el ejercicio económico de la entidad y no e>,cederá de 12 m eses.

No obstante, el período impositivo se entiende concluido en todo caso en los siguientes supuestos:
l . cuando la entidad se extinga, siendo las causas de extinción las previstas en la regulación sustantiva de la entidad,
2. cuando tenga lugar un cambio de residencia de la entidad residente, en principio, en territorio espafiol al extranjero;
3. cuando se produzca la transformación de la forma jurídica de la entidad y ello determine la no sujeción al IS de la entidad resultante;

4. cuando se produzca la transformación de la forma jurídica de la entidad y ello determine la modificación del tipo de gravamen en el IS o la aplicación de un régimen tributario especial a la entidad resultante;

En cuanto al devengo, éste se producirá el último día del período im positivo y determina, a efectos del IS. la legislación aplicable y el momento respecto al cual ha de fijarse l a base im ponible. 

El periodo impositivo no tiene por que coincidir con el afio natural. ya que el ejercicio económico puede empezar en cualquier fecha del afio. no superando los 12 meses.


\clearpage
\section{IS: Esquema de liquidación}
El esquema de liquidación del Impuesto es el conjunto de elementos técnicos que. entrelazados entre sí, determinan la cuota a ingresar o a devolver. La estructura del esquema de liquidación. que iremos desgranando a continuación a lo largo del tema. es la siguiente:

\subsection{Renta sujeta a gravamen: ¿Cómo se calcula la Base Imponible?}
La Base Imponible estará constituida por el importe de la renta obtenida en el período impositivo minorada por la compensación de bases imponibles negativas de ejercicios anteriores (BINs).

La base se determina, con carácter general. en régimen de estimación directa: por el de estimación objetiva en algunos regímenes especiales; y, subsidiariamente, por el de estimación indirecta de acuerdo con la Ley General Tributaria.

En el régimen de estimación directa. el cálculo de la Base Imponible se reduce a corregir el Resultado contable. calculado de acuerdo con l a nonnativa mercantil. con aquellas correcciones que, en su caso, establezca la normativa fiscal y que veremos a continuación. En el caso de estimación objetiva se calcula la base en función de signos. índices o módulos.

Los ingresos y gastos derivados de las transacciones o hechos económicos se imputarán al período impositivo e n que se produzca su devengo, con arreglo a la n01mativa contable, con independencia de la fecha de su pago o de su cobro, respetando la debida correlación entre unos y otros.

\subsection{Ajustes}
Como se refleja en el esquema de liquidación, el resultado contable es objeto de una serie de ajustes positivos o negativos para obtener la renta del pe1iodo, de la que, una vez compensadas las BINs, se obtiene la base imponible. Vamos a ver los principales ajustes.

A11wrtizaciones
Por amortización hemos de entender aquel proceso sistemático que tiene como finalidad la imputación al resultado del ejercicio del deterioro de elementos patrimoniales cuya vida útil supera el ejercicio económico ordinario.

Las amortizaciones serán deducibles cuando respondan a l a depreciación efectiva que sufran los elementos del inmovil izado material, intangible y las inversiones inmobiliarias por funcionamiento, uso, disfrute u obsolescencia.

Entendiéndose efectivamente producida la depreciación cuando la amortización se calcule de acuerdo con alguno de los siguientes métodos recogidos en la Ley del l mpuesto:
a) en primer l ugar, que sea el resultado de aplicar los coeficientes de amortización de las tablas de amortización aprobados oficialmente. Siendo necesario destacar a este respecto que con la Reforma Fiscal de 2 0 1 4 se simplificaron las tablas de amo11ización y se introdujeron en el cuerpo de la Ley del Impuesto.
b) en segundo l ugar, mediante la aplicación de un coeficiente constante. E l porcentaje constante se determinará ponderando el coeficiente de amortización lineal obtenido a partir del período de amortización según tablas de amortización oficialmente aprobadas, por los unos coeficientes.
e) otro método reconocido en la Ley es el de número�� dígitos en fu nción del período de
amo11ización establecido en tablas. y aplicable a todos los activos excluidos los edificios,
mobiliario, enseres y demás equipos de oficina.
d) o mediante un plan de amortización formulado por el sujeto y aceptado por la administración.
e) y por último, cuando el sujeto justifique su impone.

En lodo caso es amortizable el precio de adquisición o coste de producción. excluido. en su caso, el valor residual.

Como excepción a principio de inscripción contable. se puede destacar la figura de la libertad de amortización. que. aunque su ámbito de aplicación se ha visto restringido por la Reforma Fiscal de 2014, es aplicable a, entre otros elementos: 
Activos materiales, intangible e inversiones inmobiliarias adquiridos en los primeros 5 afios por las sociedades anónimas laborales y las explotaciones asociativas prioritarias agrarias. Los elementos del inmoviliz ado material e intangible, excluidos los edificios, afectos a las actividades de investigación y desarrollo. Los edificios podrán amortizarse de forma lineal durante un período de I O afios. en la pa11e que se hallen afectos a las actividades de investigación y desarrollo. Los gastos de investigación y desarrollo activados como inmovilizado intangible. excluidas las amortizaciones de los elementos que disfruten de libertad de amortización. Los elementos del inmovilizado material o inta ngible de Ins entidades que tengan la calificación de explotaciones asociativas prioritarias de acuerdo con lo dispuesto en la Ley 19/1995, de 4 de julio, de modernización de las explotaciones agrarias, adquiridos durante los cinco primeros afios a partir de la fecha de su reconocimiento como explotación prioritaria.

Los elementos del inmovilizado material nuevos, cuyo valor unitario no exceda de 300 euro��, hasta el límite de 25.000 euros referido al período impositivo. Si el período impositivo tuviera una duración inferior a un aiio, el límite señalado será el resultado de multiplicar 25.000 euros por la proporción existente entre la duración del período impositivo respecto del año.

\textbf{Pérdidas por deterioro de valor de eleme11tos patri111011iales}

Serán deducibles las pérdidas por deterioro de los créditos derivadas de las posibles insolvencias de los deudo1'es. cuando en el momento del devengo del Impuesto concun-a alguna de las siguientes circunstancias:
a) Que haya transcurrido el plazo de 6 meses desde el vencimiento de la obligación.
b) Que el deudor esté declarado en situación de concurso.
e) Que el deudor esté procesado por el delito de alzamiento de bienes.
d) Que las obligaciones hayan sido reclamadas judicialmente o sean obj eto de un litigio judicial o procedimiento arbitral de cuya solución dependa su cobro.

\textbf{Provisiones}
No son deducibles los siguientes gastos asociados a provisiones:
a) Los gastos por obligaciones implícitas o tácita<;. 
b) Los concemientes a los costes de cumplimiento de contratos que excedan a los beneficios económicos que se esperan recibir de los mismos.
c) Los derivados de reestructuraciones. excepto s1 se refieren a obligaciones legales o contractuales.
d) Los relativos al 1iesgo de devoluciones de ventas.
e) Los de personal que se correspondan con pagos basados en instrumentos de patrimonio, tanto si se satisface en efectivo o mediante la entTega de dichos elementos.
f) Los gastos por actuaciones medioambientales sólo son deducibles cuando respondan a un plan aprobado por la Administración.

Estos gastos serán deducibles cuando la provisión se aplique a su finalidad.

\textbf{Gastos no deducibles}

Por otro lado, a la hora de calcular la base imponible a partir del resultado contable existen una serie de gastos que la normativa fiscal considera en todo caso como no deducibles, como son. entre otros:
a) Los que representen una retribución de los fondos propios.
b) Los derivados de la contabilización del Impuesto sobre Sociedades. N o tendrán la consideración de ingresos los procedentes de dicha contabilización.
e) Las multas y sanciones penales y administrativas, el recargo de apremio y el recargo por presentación fuera de plazo de declaraciones liquidaciones y autoliquidaciones.
d) Las pérdidas del juego
e) Los donativos y liberalidades. No se entienden como tales los gastos por relaciones públicas con clientes o proveedores ni los que con aiTeglo a los usos y costumbres se efectúen con respecto al personal de la empresa ni los realizados para promocionar, directa o indirectamente, la venta de bienes y prestación de servicios, ni los que se hallen correlacionados con los mgresos.
f) Los gastos derivados de acciones contrarias al ordenamiento jurídico.
g) Los gastos de servicios corrcspondícnt.es a operaciones realizadas, directa o indirectamente, con personas o entidades residentes en países o territorios calificados reglamentariamente por su carácter de paraísos fiscales. o que se paguen a través de personas o entidades residentes en los mismos. excepto que el sujeto pasivo pruebe que el gasto devengado responde a una operación o transacción efectivamente realiz.ada.
h) Los gastos financieros derivados con deudas con empresas del grupo destinadas a la adquisición a otras entidades del grupo de participaciones en el capital o fondos propios de otras entidades o a la realización de apo11aciones al capital salvo que se acrediten que existen motivos económicos válidos. En caso de endeudamiento no vinculado, la deducibilidad de los gastos financieros se limita a un 30 % del beneficio operativo en caso de que los gastos superen el millón de euros.
i) Los gastos que excedan para cada perceptor del importe de I millón euros o el i mpo11e exento en caso de ser superior por extinción de la relación laboral.
j) Las pérdidas por deterioro de valores en el capital o fondos propios con carácter general.
k) Como consecuencia del considerando 9 de la Directiva (lJE) 201 7/952 del Consejo, de 29 de mayo de 20 1 7. seí'íala que las normas sobre las asimetrías híbridas «( . . . ) deben tratar situaciones de asimetría derivadas de dobles deducciones. de conflictos en la calificación de los instrumentos financieros, pag.os y entidades, o de la atribución de pagos. Dado que las asimetrías híbridas podrían dar lugar a una doble deducción o a una deducción sin inclusión, es necesario establecer nonnas en virtud de las cuales el Estado miembro implicado deniegue la deducción de un pago, unos gastos o unas pérdidas, u obligue al contribuyente a incluir el pago en su renta imponible. según sea adecuado ( . . . )». se crea un nuevo artículo 1 5 bis el cual establece que no serán fiscalmente deducibles los gastos correspondientes a operaciones realizadas con personas o entidades vinculadas residentes en otro país o territorio que. como consecuencia de una calificación fiscal diferente en estas del gasto o de la operación, no generen un i ngreso, generen  un ingreso exento o sujeto a una reducción del tipo impositivo o a cualquier deducción o devolución de impuestos distinta de una deducción para evitar la doble imposición jurídica.
l) Lim itación en la deducibilidad de gastos financieros. Los gastos financieros netos serán deduci bles con el límite del 30 por ciento del beneficio operativo del ejercicio. 


\textbf{Ajustes de valoracián: Operaciones a mlor de mercado}

La Ley del Impuesto señala que las variaciones de valor originadas por aplicación del criterio del valor razonable no tendrán efectos fiscales mientras no deban imputarse a la cuenta de pérdidas y ganancias. El importe de las revalorizaciones contables no se integrará en la base imponible, excepto cuando se lleven a cabo en virtud de nonnas legales o reglamentarias que obliguen a incluir su impo11e en el resultado contable.

Sin embargo, la Ley del Impuesto seí'íata una serie de elementos patrimoniales que se deberán valorar a valor no1mal de mercado, como son:
a) Los transmitidos o adquiridos a título lucrativo. A estos efectos no se entienden como adquisiciones a título lucrativo las subvenciones
b) Los aportados a entidades y los valores recibidos en contraprestación.c) Los transmitidos a los socios por causa de disolución, separación de los mismos, reducción del capital con devolución de aportaciones. reparto de la prima de emisión y distribución de beneficios.
d) Los transmitidos en vi11ud de fusión, absorción y escisión total o parcial. Debiendo destacar que en todas estas operaciones es la transmitente la que integrará en su base la diferencia entre el valor n01mal de mercado y su valor contable.
e) Y además los adquiridos por pennuta.
f) Los adquiridos por canje o conversión. Se entiende por valor nonnal del mercado el que hubiera sido acordado en condiciones nonnales de mercado entre partes independientes. Además de estos supuestos de valoración a precios de mercado antes mencionados, debe hacerse especial referencia a las operaciones vinculadas. La regulación del tratamiento fiscal de este tipo de operaciones ha sufrido una modificación sustancial como consecuencia de la aprobación de la Ley de Medidas para la prevención del fraude fiscal ya que, como sefiala su Exposición de Motivos. se ha procedido a adaptar la normativa española a las directrices de la OCDE en materia de precios de transferencia.

La actual regulación establece la obligación de los sujetos pasivos de valorar las operaciones vinculadas por su valor normal de mercado, definido como aquel que se habría acordado por personas o entidades independientes en condiciones de libre competencia.

Para la determinación del valor normal de mercado se debe aplicar uno de los siguientes métodos:
a) Método del precio libre comparable, que compara el precio del bien o servicio en una operación entre partes vinculadas con el precio de bienes o servicios idénticos o de características similares entre partes no vinculadas.
b) Método del coste incrementado, por el que se añade al valor de adquisición o coste de producción del bien o servicio el margen habitual en operaciones idénticas o similares entre partes no vinculadas.
c) Método del precio ele reventa. por el que se sustrae del precio de venta de un bien o servicio el margen que aplica el propio revendedor en operac10nes idénticas o similares con personas o entidades independientes.
d) Método de la distribución del resultado. por el que se asigna, a cada persona o entidad vinculada que realice de fonna conjunta una o varias operaciones. la parte del resultado común derivado de dicha operación en función de u n criterio adecuado.
e) Método del margen neto del conjunto de operaciones,_por el que se atribuye a las operaciones realizadas con una persona o entidad vinculada el resultado neto, calculado sobre costes, ventas o la magnitud que resulte más significativa, obtenido en un estudio de comparabilidad entre empresas no vinculadas.

\subsection{Reducciones: reserva de capitalización y nivelación}
Reserl'a de capitali:.ació11
,- Consiste en una minoración en la base imponible del 1 0% del incremento de los fondos propios. con el límite del 1 0% de la base imponible. Si no hay base imponible positiva en ese ejercicio se podrá aplicar en los tres siguientes.
:;... Requisitos: Creación de una reserva indisponible sin que sea necesario invertir en activos de la empresa y los fondos propios no deben verse minorados en un plazo ele 5 a,'íos, salvo por pérdidas.
;,, Objetivos: Incremento de la competiti vidad de la empresa española. reducción de l a diferencia en el tratam iento fiscal de la financiación ajena y l a financiación propia. así como fomentar la capitalización empresarial.
r 


Como ejemplo podemos destacar el siguiente:
3 1 / 1 2/20 14:
Capital: l 0.000
Bº del ejercicio: 3.000
3 1/ l 2i20 1 5
Capiial: 1 0.000
Reservas: 3.000
Bº del ejercicio: 3.500
La entidad decide llevar a reservas todo el beneficio del ejercicio.
RC: 1 0% 6 FP. Para su cálculo se prescinde de los beneficios del ejercicio y del ejercicio
anterior.
Límite: 1 0% de la BI del período impositivo. En caso de insuficiencia de base imponible,
posibil idad de aplicar en dos años siguientes.
6 FP= 3.000
Bl previa (no hay ajustes) = 3.500
RC = - 300 ( 1 0% 3.000)
B I = 3.200
CI (25%) = 800

Tipo efectivo en este ejemplo, 22,85%

Reserva de nfre/ac:ión
► Consiste en una m inoración del 1 0% ele l a base imponible, siempre que se constituya una reserva indisponible. con el lím ite máximo de I m i llón de euros.
► Esta cantidad se compensará con bases imponibles negativas (BINS) futuras, en un plazo de 5 años.
► Si en ese plazo l a empresa no genera BINS. la empresa adicionará a su base imponible la correspondiente reserva transcurrido ese plazo de 5 años.
► En definitiva, se dejan de pagar impuestos en el momento de constitución de la reserva por las BINS futuras. En el caso en que no haya BINs se produce un diferimiento en la tri butación.
► Objetivos: disminución de l a diferencia entre el tratamiento fiscal de la financiación ajena y la financiación propia y mayor estabilidad tributación PYME
► En el caso de aplicación conjunta de la reserva de capital ización y de nivelación por una PYME, su tipo efectivo puede reducirse hasta un 20%
► Como ejemplo podemos destacar el siguiente:
La entidad A obtiene en el afio O una base imponible positiva por valor de 200.000 euros.
Situación anterior:
Cuota integra: 200.000 x 25% = 50.000
Tras la refom1a:
Puede minorar su base ú11ponible en 20.000 euros si constituye una reserva por ese
i mporte.
Cuota íntegra: (200.000 - 20.000) x 25% = 45.000
( Esto supone una tributación efectiva en el año O al 22.5%)
Si en el año I a 5 obtiene una base imponible negativa por valor de -20.000 euros, se
'·compensa'· con los 20.000 euros del año O. En la práctica, se trata de una compensación
anticipada de BINS futuras.
S i en esos 5 arios no ha obtenido BINS, en la declaración del año 5 se integran 20.000 en
la base imponible tributando en ese momento.
20.000 X 25'1/ó = 5.000
(En este caso se produce un diferimiento de la tributación del 2.5% de la base imponible
del año 0).
► Ejemplo de aplicación conjunta de reserva de capitalización y de nivelación:
3 1/1 2/2014: 3 1/ 1 2/20 15
Capital: 10.000
8° del ejercicio: 3 .000
Capital: 10.000
Reservas: 3 .000
Bº del ejercicio: 3.500
La entidad decide llevar a reservas todo el beneficio del ejercicio.
RC: 10% /:,. FP. Para su cálculo se prescinde de los beneficios del ejercicio y de los del
ejercicio anterior.
Límite 10% de la BI del período impositivo. En caso ele insuficiencia de base imponible.
posibilidad de aplicar en dos años siguientes.
RN: 10% BI ele la entidad. Límite I millón de euros.
/:,. FP= 3 .000
B I previa (no hay ajustes) = 3.500
R C = - 300 ( 10% 3.000)
B 1 = 3.200 RN = - 320 Bl con RN = 2.880 Tipo efectivo en este ejemplo del 20,57%, Deberá constituir una reserva de nivelación de 3 20 y otra ele capitalización de 300

\subsection{Compensación de bases imponibles negativas}
Una vez realizados los ajustes al resultado contable, se obtiene una base previa que puede ser reducida en el importe de las bases imponibles negativas de ejercicios anteriores. en un intento de la Ley de paliar los efectos distorsionadores que sobre el cálculo del resultado global tiene la división de la vida económica de la entidad en unos períodos que no dejan de ser <lrtificiales

Las bases imponibles negativas que hayan sido objeto de liquidación o autoliquidación podrán ser compensadas con las rentas positivas de los períodos impositivos sigu ientes con el límite del 70 por ciento de la base imponíble previa a la aplicación de la reserva de capitalización y a su compensación. En todo ca<;o, se podrán compensar en el periodo impositivo bases imponibles negativas hasta el importe de I millón de euros.

\subsection{Cuota íntegra}
Para la determinación de la cuota integra. sobre la base imponible se aplicará el tipo de gravamen, que con carácter general cs. a partir del 2016, del 25%. excepto para las entidades cuyo impone neto de la cifra de negocios del período impositivo inmediato anterior sea inferior a I millón de euros que será el 23 por ciento.Eliminándose de esta fonna la ventaja con la que contaban previamente a la Reforma Fiscal las entidades de reducida dimensión. que tributaban a un tipo más bajo (el 25% por los primeros 300.000 euros frente al 30% de las entidades ordinarias). Como ya se ha comentado en las novedades del ejerecicio 1022, se introdujo un nuevo ai1ículo 30 bis que regula la tributación mínima para contribuyentes cuyo importe neto de la cifra de negocios sea al menos de 20 millones de euros durante los 1 2 meses anteriores a la fecha en que se inicie el período impositivo o que tributen en el régimen de consolidación fiscal. La cuota líquida no podrá ser inferior al resultado de aplicar el 1 5 por ciento a la base imponible, minorada o incrementada, en su caso y según corresponda, por las cantidades derivadas la Reserva por

Inversiones en Canarias. Dicha cuota tendrá el carácter de cuota líquida mínima. Será del I O por ciento en las entidades de nueva creación y el 18 por ciento si se trata de entidades crédito, así como las entidades que se dediquen a la exploración, investigación y explotación de yacimientos y almacenamientos subterráneos de hidrocarburos.

No obstante, no será de aplicación esta tributación mínima a los contribuyentes que tributen a los tipos de gravamen que t1ibuten al 10 %, 1 % y 0% ni a las Sociedades Anónimas Cotizadas de Inversión en el Mercado Inmobiliario.

En el caso de las cooperativas, la cuota líquida m inima no podrá ser inferior al resultado de aplicar el 60 por ciento a la cuota íntegra.

En las entidades de la Zona Especial Canaiia, la base imponible positiva no incluirá la pai1e de la misma correspondiente a las operaciones real izadas material y efectivamente en el ámbito geográfico de dicha Zona que tribute al tipo ele gravamen especial del Régimen Económico y Fiscal de Canarias.

Además, la regulación del I S prevé una serie de tipos específicos, entre otros, podemos destacar: 
a) Tributarán al tipo del 30% las entidades de crédito, así como las entidades que se dediquen a la exploración. investigación y explotación de yacimientos y almacenamientos subterráneos de hidrocarburos.
b) Tributarán al 20% las cooperativas fiscalmente protegidas.
c) Las entidades de nueva creación que realicen actividades económicas tributarán tipo de gravamen del 1 5% en el primer período impositivo con base impositiva positiva y el siguiente.
d) Tributarán al 10% las entidades de régimen tiscal de entidades sin fines lucrativos e incentivos fiscales al mecenazgo
e) Tributarán al 1 %: las sociedades de inversión de capital variabk siempre que el número de accionistas sea como mínimo de I OO. los fondos de inversión de carácter financiero siempre que el número de partícipes sea como mínimo de 1 00, las sociedades de inversión inmobiliaria y los fondos de inversión inmobiliaria siempre que el número de socios y pa11ícipes sea como mínimo de 100 y que tengan por objeto exclusivo la inversión en cualquier tipo de inmueble de naturaleza urbana para su arrendamiento, las sociedades de inversión inmobiliaria y los fondos de inversión inmobiliaria que desa1rollen la actividad de promoción exclusivamente de viviendas para destinarlas a su arrendamiento y cumplan condiciones adicionales, el fondo de regulación de carácter público del mercado hipotecario.
f) Tributarán al 0% los fondos de pensiones.

\subsection{Deducciones y bonificaciones}
Deducciones para evitar doble imposición

Se trata de mecanismos que tienen por objeto impedir que una renta sea sometida dos veces a gravamen en imposición directa.
Los métodos que tradicionalmente se han empleado para evitar el fenómeno de la doble imposición son dos:
el método de exención. en el que las rentas que ya han sido sometidas a un gravamen anterior no se i ntegran en base del preceptor (que ya ha sido analizado en el epígrafe de rentas exentas)
el método de imputación con deducción. las rentas se integran en base y son gravadas, pero se aplica una deducción que minora en todo o en parte la tributación en el impuesto personal del preceptor.

En el actual IS, conviven ambos sistemas de minoración del fenómeno de la doble imposición, si bien su regulación ha sido simplificada a través de la Reforma Fiscal de 2014, dándole el mismo tratamiento a la doble imposición de rentas por gravámenes previos interiores y a la doble imposición derivada de rentas ya gravadas previamente en el extranjero.

En concreto, nos centramos en este apa11ado, en el método de imputación con deducción por doble imposición internacional jurídica y la deducción por doble imposición internacional económica.

En primer lugar, en cuanto a la deducción por doble imposición internacional jurídica, trata de evitar la doble imposición existente al gravarse una renta obtenida por un mismo contribuyente en Estados diferentes. Sería el caso de un beneficio generado por una sociedad espafiola en un país extranjero por el que paga el impuesto de dicho país en base al criterio de territorial idacl ( pago en fuente) y también paga el I S español por dicho beneficio en base al criterio de residencia. En este caso. cuando en la base imponible del contribuyente se integren rentas positivas obtenidas y gravadas en el extranjero. se deducirá de la cuota íntegra la menor de las dos cantidades siguientes:
a) El importe efectivo de lo satisfecho en el extranjero por razón del gravamen de naturaleza idéntica o análoga a este Impuesto.
b) El importe de la cuota íntegra que en Espafia correspondería pagar por las mencionadas renta<; si se hubieran obtenido en territorio espa,iol.

En segundo lugar, en cuanto a la deducción por doble imposición internacional económica, trata de evitar l a doble imposición existente al gravarse una renta en entidades distintas en dos Estados d i ferentes. Sería el caso de los dividendos percibidos por la matriz de una filial situada en el extranjero, que tributan primero como beneficio de la filial por el impuesto extranjero y luego como dividendo al integrarse en la base imponible de la matriz por el IS nacional. En este caso, la Ley establece que se deducirá el impuesto efectivamente pagado en el extranjero respecto de los beneficios con cargo a los cuales se abonan los dividendos, en la cuantía correspondiente de tales dividendos, siempre que dicha cuantía se incluya en la base i mponible del contribuyente. Para l a aplicación de esta deducción

\textbf{Deducción para i11centivar delermi11adas actividades}

Estas deducciones para incentivar detenninadas actividades se rigen por unas no1mas comunes: una misma inversión no puede dar lugar a deducción en más de una entidad. que el conj unto de las deducciones no podrá superar el 25% de la cuota íntegra minorada en las deducciones por doble imposición y las bonificaciones. que las cantidades no deducidas por superar el tope podrán trasladarse a los 1 5 ejercicios siguientes, las deducciones por J+D+i pueden quedar excluidas del límite y aplicarse con un descuento del 20 % y en caso de insuficiencia de cuota se puede solicitar su abono a la Administración.

En cuanto a las principales deducciones para incentivar detern1inadas actividades son las siguientes:
a) Deducción por actividades de investigación científica e innovación tecnológica.
Se establece una deducción general del 25% de los gastos efectuados por investigación científica en el periodo impositivo. Para las actividades de innovación tecnológica se establece una deducción del 1 2% de una serie de gastos.
b) Deducción por inversión en producciones cinematográficas o audiovisuales espaí'íolas.Dan derecho a la deducción del 30% respecto d��I primer millón de base de la deducción y del 25% sobre el exceso de dicho importe. Por otro lado, los productores registrados en el Registro de Empresas Cinematográficas del Ministerio de Educación, Cultura y Deporte que se encarguen de la ��jecución de una producción extranjera de largometrajes cinematográficos o de obras audiovisuales. tendrán derecho a una deducción del 30% de los gastos realizados en territorio espafiol. siempre que los gastos realizados en territorio español sean, al menos. de I millón de euros y del 25% sobre el exceso de dicho importe.

c) Deducción por creación de empleo discapacitado. Se deducen 9000€ por cada persona/año de incremento del promedio de plantilla de trabajadores con discapacidad.

d) La deducción por creación de empleo. Las entidades que contraten a su primer trabajador a través de un contrato de trabajo por tiem po indefinido de apoyo a los emprendedores que sea menor de 30 aiios. podrán deducirde la cuota íntegra la cantidad de 3 .000 euros.

\textbf{Bonijicaciones}

Podemos destacar una serie de minoraciones a realizar sobre la cuota íntegra conocidas como bonificaciones:

l . Por las rentas obtenidas en Ceuta y Melilla para las entidades espafiolas domiciliadas fiscalmente en Ceuta y Melilla. y aquellas otras no domiciliadas allí. tanto españolas como extranjeras y que operen allí mediante establecimiento pemrnnentc. La bonificación será del 50% de la parte de la cuota íntegra que correspondiese a estas rentas.

2. Bonificación del 99% para las rentas obtenidas en la prestación de determinados servicios por parte de entidades dependientes del Estado, Comunidades Autónomas o de Corporaciones Locales.

\subsection{Gestión y liquidación del impuesto}
Pasaremos a continuación a examinar los aspectos relativos a la gestión del impuesto. comenzando por señalar que cada Delegación de la Agencia Estatal de la Administración Tributaria debe llevar un índice de entidades no exentas que tengan su domicilio fiscal en su ámbito territorial. Los titulares de los Registros Públicos y los Notarios deben remitir de forma mensual a la AEAT una relación de inscripciones o de documentos otorgados respectivamente y que se refieran a constitución, establecimiento. modificación. transformación o extinción de entidades. Lo que no exime a los sujetos de sus deberes censales.

\textbf{Obligación de declarar.}

En cuanto al modo de declarar. están obligados a presenta la autoliquidación del IS los sujetos pasivos en el lugar y forma detemúnados por el Ministro de Hacienda y Función Pública. Si bien el devengo del impuesto es el último día del periodo impositivo. la declaración debe presentarse en el plazo de ��5 días naturales siguientes a los 6 meses posteriores a la conclusión del período impositivo. como habitualmente el ejercicio económico coincide con el año natural, el límite para la presentación de la declaración del impuesto seria el 25 de julio . 

No estando obligados a declarar los sujetos que estén exentos. mientras que las entidades parcialmcme exentas tampoco deberán presentar declaración cuando no obtengan rentas sujetas al mismo o las únicas rentas sujetas que obtengan sean rendimientos sujetos a retención.

\textbf{Medio de pago.}

Los sujetos pasivos al tiempo de presentar su declaración deberán determinar la deuda correspondiente e ingresarla en el lugar y en la forma determinados por el MinistTo de Hacienda y Función Pública.

Pudiendo realizarse el pago de la misma mediante efectivo o la entrega de bienes integrantes del Patrimonio Histórico Español que estén inscritos en el Inventario General de Bienes Muebles o en el Registro General de Bienes de Interés Cultural, sin que las rentas obtenidas por la transmisión se integren en la base imponible. 

Cuando la cuota resultante de la autoliquidación o. en su caso, de la liquidación provisional sea inferior a la suma de las cantidades efectivamente retenidas, ingresos a cuenta y pagos fraccionados. la Administración tributaria procederá a devolver de oficio el exceso sobre la citada cuota.


\subsection{III.VUI . Pagos a cuenta }

Por otra parte, en el ámbito del IS, podemos hablar de obligación de ingresar cantidades en el Tesoro respecto de los pagos a cuenta, que tienen la consideración de deuda tributaria distinguiendo entre retenciones, ingresos a cuenta y pagos fraccionados.

/JI. VIII. /. Retenciones e ingresos a cuenta

Son cantidades que el pagador de detem1inados rendimientos ingresa directamente a la Hacienda Pública pudiendo el perceptor deducirlos de su cuota reducida, en tanto en cuanto suponen anticipos de la obligación que corre a su cargo.

Están obligados a efectuar estas operaciones las entidades. incluidas las comunidades de bienes y las de propietarios, que satisfagan o abonen rentas sujetas a este Impuesto. También ��stán obligados a retener e ingresar los empresarios individuales y los profesionales respecto de las rentas que satisfagan o abonen en el ejercicio de sus actividades empresariales o profesionales, así como las personas físicas, jurídicas y demás entidades no residentes en territorio español que operen en él mediante establecimiento permanente.

Debe practicarse retenc1on. en concepto de pago a cuenta del Impuesto sobre Sociedades con-espondiente al perceptor, respecto de, entre otras:
a) Las rentas derivadas de la pa1ticipación en fondos propios de entidades, de la cesión a terceros
de capitales propios,
b) Los premios derivados de la pa,ticipación en juegos:
c) Las contrapres1aciones ob1enidas como consecuencia de la atrihución de cargos de
administrador o consejero en otras sociedades,
d) Las rentas procedentes de la cesión del derecho a l<1 explotación de la imagen.
e) Las rentas procedentes del arrendamiento de inmuebles urbanos. 

Con carácter general, constituye la base para el cálculo de la obligación de retener la acon1raprestación íntegra exigible o satisfecha.

El porcentaje general al que se debe efectuar la retención o el ingreso a cuenta es del 19%.

\subsubsection{III. VIII.Jl. Pagosfracc,ionados}

Son cantidades que los contribuyentes del IS deben anticipar a cuenta del impuesto final, en los primeros vei nte días naturales de los meses de abril, octubre y diciembre, teniendo el carácter de deuda tributaria. Existen dos formas de cálculo del pago fraccionado:

En primer lugar, aplicar el porcentaje de 1 8%. a la cuota íntegra minorada en deducciones, bonificaciones y pagos a cuenta del último ejercicio con plazo voluntario de declaración vencido el día primero de los meses de abril. octubre y diciembre. Fij ándose u n plazo de 20 días a partir de esas fechas para el desembolso del pago.

Y, en segund9 lugar, aplicar el porcentaje que resulte de multiplicar por cinco séptimos el tipo de gravamen redondeado por defecto a la base imponible obtenida en lo que va de ejercicio al cabo de los tres. nueve y once meses de cada afio natural.

Estarán obligados a aplicar la modalidad a que se refiere este apa1tado los contribuyentes cuyo importe neto de la cifra de negocios haya superado la cantidad de 6 millones de euros durante los 1 2 meses anteriores a la fecha en que se inicie d período impositivo al que corresponda el pago fraccionado.










\clearpage
\section{Impuesto sobre Sociedades: Marco jurídico-conceptual}
\puntosclave{
    
}

El Impuesto sobre Sociedades comparte con el IRPF su fundamento último en el art. 31.1 CE y su condición de pilar de la imposición directa española, pero conviene no dar por supuesto, solo por esa coincidencia constitucional, que su desarrollo se explica igual. El IRPF nace en 1978 como ruptura frente a un sistema cedular fragmentado; el IS, en cambio, ya tenía una identidad propia bastante consolidada antes de la Constitución, y su historia reciente está marcada menos por preguntas de justicia distributiva entre personas físicas —el gran tema del IRPF— que por preguntas de neutralidad económica, de competitividad internacional y, sobre todo, de **relación entre la contabilidad mercantil y la fiscalidad**, una tensión que no tiene equivalente en el impuesto personal y que va a condicionar buena parte de este apartado. Empiezo por el desarrollo histórico, sigo con el marco jurídico vigente, y cierro con el esquema básico de liquidación, que aquí tiene una arquitectura —partiendo del resultado contable— sustancialmente distinta de la del IRPF.

\subsection{Desarrollo histórico}
Antes de que existiera un impuesto sobre sociedades identificable como tal, la renta de las personas jurídicas en España tributaba, igual que ocurría con la renta de las personas físicas antes de 1978, dentro de un sistema de figuras de producto: la tributación de los beneficios societarios se articulaba, ya desde el siglo XIX, a través de la Contribución de Utilidades de la Riqueza Mobiliaria, sin que existiera un tributo autónomo y unificado sobre la renta de las entidades jurídicas como tales. [Seguro] El nacimiento formal del Impuesto sobre Sociedades como figura tributaria diferenciada se sitúa en la reforma de 1957, promovida por el entonces ministro de Hacienda Mariano Navarro Rubio, tal y como lo describe Sanz Gadea (2003) en su historia del impuesto recogida en el *Manual del Impuesto sobre Sociedades* del Instituto de Estudios Fiscales, obra de referencia obligada para quien quiera profundizar en esta evolución. [Seguro] Una nueva reforma tributaria en 1964 —de nuevo con ánimo de simplificación del sistema entonces vigente— dio lugar, tres años después, al Texto Refundido del Impuesto sobre Sociedades de 1967, que sin embargo no resolvió en toda su amplitud los problemas estructurales del tributo, y que llega a la Transición arrastrando aún carencias importantes.

El siguiente hito no es todavía el impuesto moderno, sino su antesala: la **Ley 50/1977, de 14 de noviembre**, de medidas urgentes de reforma fiscal —la misma generación normativa que, meses después, traería la Ley 44/1978 del IRPF, aunque conviene no sobredimensionar el paralelismo: la Ley 50/1977 no crea todavía un nuevo IS, sino que aborda cuatro cuestiones instrumentales que preparan el terreno de toda la reforma tributaria de la Transición: la regularización voluntaria de situaciones fiscales pendientes, el levantamiento del secreto bancario frente a la Hacienda Pública, la creación del Impuesto sobre el Patrimonio, y el establecimiento del delito fiscal como figura penal autónoma. Es, en otras palabras, la ley que dota a la Hacienda española de las herramientas de control y de las nuevas figuras tributarias complementarias sin las cuales una reforma sustantiva del IS —o del IRPF— habría carecido de eficacia práctica.

[Seguro] La culminación de este proceso llega con la **Ley 61/1978, de 27 de diciembre**, del Impuesto sobre Sociedades, que la Audiencia Nacional describiría años después, en su sentencia de 7 de marzo de 2002, como una norma que representó una ruptura con la tradición legislativa tributaria anterior. Entre sus efectos menos conocidos pero jurídicamente significativos está la extensión de la sujeción al impuesto a entidades hasta entonces al margen del sistema, incluidas las confesiones religiosas, salvo régimen de exención específico —un dato de interés expositivo que ilustra bien hasta qué punto la reforma de 1978 amplió el ámbito subjetivo del impuesto respecto de su antecedente inmediato. [Suposición, nexo pedagógico sin desarrollar aquí] Si te resulta útil como ancla mnemotécnica: la Ley 44/1978 del IRPF es de 8 de septiembre; la Ley 61/1978 del IS es de 27 de diciembre — ambas hijas de la misma reforma Fuentes Quintana, pero tramitadas y aprobadas con casi cuatro meses de diferencia dentro del mismo año.

El impuesto de 1978 rige durante toda la década de los ochenta sin alteración estructural, pero hacia el final de esa década se acumulan tres factores de presión que acabarán forzando una reforma integral. [Seguro] Así lo explica la propia Exposición de Motivos de la **Ley 43/1995, de 27 de diciembre**, que enumera con precisión las causas de la reforma que se disponía a aprobar: primera, la reforma parcial de la legislación mercantil operada por la **Ley 19/1989, de 25 de julio** —que adaptó el Derecho español de sociedades y de cuentas anuales a las directivas comunitarias contables, introduciendo un régimen de contabilidad mercantil moderno y armonizado con el resto de la entonces Comunidad Europea—; segunda, la propia reforma del IRPF operada por la Ley 18/1991, que ya vimos con detalle en el Tema 16; tercera, la creciente apertura de la economía española a los flujos transfronterizos de capitales; cuarta, la dispersión normativa que para entonces afectaba ya al articulado del Impuesto sobre Sociedades tras años de modificaciones parciales; y quinta, la propia evolución de la teoría hacendística y jurídico-financiera sobre la imposición societaria en el entorno comparado. [Seguro] De estas cinco causas, la primera es, con diferencia, la que produce el cambio más profundo en la arquitectura técnica del impuesto: la Ley 43/1995 adopta, por primera vez de forma sistemática, el **principio de conexión** entre contabilidad y fiscalidad —el resultado contable de la entidad, calculado conforme a las normas mercantiles ya armonizadas por la Ley 19/1989, pasa a ser el punto de partida obligado para determinar la base imponible del impuesto, sujeto después a los ajustes que la propia norma fiscal establezca—. [Probable] Este es el elemento que, adelanto ya sin desarrollarlo, va a convertirse en la bisagra estructural de todo el apartado 2 de este tema: antes de 1995, la base imponible del IS se determinaba de forma autónoma según reglas fiscales propias, con relación solo indirecta con la contabilidad mercantil de la empresa; desde 1995, la contabilidad deja de ser un simple punto de referencia y pasa a ser, por mandato legal expreso, el punto de partida obligatorio de la liquidación —un giro metodológico sin equivalente en la historia del IRPF, que nunca ha tenido una fuente de determinación de la renta externa al propio impuesto.

La siguiente etapa no es una reforma sustantiva sino una consolidación técnica: el **Real Decreto Legislativo 4/2004, de 5 de marzo**, aprueba el Texto Refundido de la Ley del Impuesto sobre Sociedades (TRLIS), unificando en un único cuerpo normativo las numerosas modificaciones parciales que la Ley 43/1995 había ido acumulando desde su aprobación. [Seguro] Entre las novedades sustantivas incorporadas en ese proceso de refundición —no meramente técnicas, pese a la naturaleza formal del instrumento normativo empleado— destacan la introducción de medidas de corrección de la doble imposición interna e intersocietaria, la separación definitiva de la tributación de los no residentes en un texto normativo propio (dando lugar al actual Impuesto sobre la Renta de No Residentes como figura autónoma, ya no integrada dentro del IS), la desaparición del antiguo régimen de **transparencia fiscal** —un mecanismo, ya entonces problemático, para imputar directamente a los socios la renta de determinadas sociedades cerradas o de mera tenencia de bienes— y su sustitución por el régimen de **sociedades patrimoniales**.

[Seguro] Llegamos así a la norma vigente, la **Ley 27/2014, de 27 de noviembre**, del Impuesto sobre Sociedades, aprobada —como ya sabes por el Tema 16— dentro del mismo paquete de reforma tributaria integral derivado del Informe de la Comisión de Expertos de 2014 que dio lugar, simultáneamente, a la Ley 26/2014 del IRPF. La propia Exposición de Motivos de la Ley 27/2014 enmarca la reforma en la necesidad de modernizar el impuesto y de favorecer la competitividad empresarial, y entre sus rasgos definitorios destacan tres, que conviene fijar ya como anticipo del resto del tema: la **reducción del tipo general de gravamen, del 30% al 25%**, acompañada —siguiendo exactamente la misma doctrina de *ensanchamiento de bases y reducción de tipos* que ya analizamos a propósito del Informe Lagares en el Tema 16— de un endurecimiento simultáneo de varias reglas de determinación de la base imponible (limitación reforzada a la deducibilidad de gastos financieros, restricciones a la compensación de bases imponibles negativas); la introducción, ex novo, de los conceptos de **actividad económica** y de **entidad patrimonial**, alineados conceptualmente —aunque no idénticos— con las categorías ya empleadas en el IRPF; y la sustitución de la antigua deducción por reinversión de beneficios extraordinarios por dos figuras nuevas, la reserva de capitalización y la reserva de nivelación, que trataremos con detalle en el apartado 3.

[Seguro] La historia no se detiene en 2014. Dos hitos posteriores, ambos posteriores a la reforma estructural pero de calado suficiente como para no omitirlos aquí, cierran el relato hasta la actualidad. Por un lado, la **Ley 22/2021**, de Presupuestos Generales del Estado para 2022, introdujo una **tributación mínima doméstica del 15%** (art. 30 bis LIS) para grandes empresas —aquellas con importe neto de cifra de negocios igual o superior a 20 millones de euros— y para los grupos que tributan en régimen de consolidación fiscal, estableciendo que la cuota líquida no pueda situarse por debajo de ese 15% de la base imponible, con independencia de las deducciones que en otro caso resultarían aplicables. Por otro, la ya citada **Ley 7/2024** trasladó a España el **Pilar 2 de la OCDE/G20**, creando el Impuesto Complementario para garantizar un nivel mínimo global de imposición del 15% a los grandes grupos multinacionales y nacionales —una figura formalmente distinta del art. 30 bis LIS de 2021, con la que conviene no confundirla: una es una regla interna de tributación mínima sobre la base imponible española de cualquier gran empresa; la otra es un mecanismo internacional coordinado, dirigido específicamente a corregir la localización de beneficios de grupos multinacionales en jurisdicciones de baja tributación. [Probable] Esta doble capa de tributación mínima, aprobada con apenas tres años de diferencia y por razones distintas —consolidación fiscal doméstica en el primer caso, coordinación fiscal internacional en el segundo— es el mejor indicador de hacia dónde se dirige la reforma del impuesto en el momento actual, y será el punto de llegada natural del debate sobre competencia fiscal internacional que desarrollaremos en el apartado 2.


\subsection{Marco jurídico actual: Ley 24/2015}
El IS vigente se regula en la Ley 27/2014, de 27 de noviembre, desarrollada reglamentariamente por el Real Decreto 634/2015, de 10 de julio (Reglamento del Impuesto sobre Sociedades). Como en el IRPF, rige supletoriamente la Ley 58/2003, General Tributaria. [Seguro] Pero aquí aparece de inmediato una fuente normativa sin equivalente en el impuesto personal: el Código de Comercio y el Plan General de Contabilidad (aprobado por Real Decreto 1514/2007, con sus adaptaciones sectoriales), porque —como ya anticipamos en el bloque histórico a propósito de la Ley 43/1995— el punto de partida legal de la liquidación es el resultado contable de la entidad, determinado conforme a esas normas mercantiles. [Probable] Esto no convierte al Derecho contable en fuente del Derecho tributario en sentido estricto —la Ley 27/2014 sigue siendo la única norma que puede alterar la tributación—, pero sí lo convierte en una fuente de remisión obligada y constante, hasta el punto de que cualquier reforma del Plan General de Contabilidad tiene, de forma casi automática, repercusión en la base imponible del impuesto salvo que la propia LIS establezca un ajuste específico para neutralizarla. Es una dependencia que no tiene parangón en el esquema de liquidación del IRPF que construimos en el Tema 16.

2.2. Ámbito de aplicación

[Seguro] El art. 2 LIS extiende la aplicación del impuesto a todo el territorio español, incluidas las zonas adyacentes a las aguas territoriales sobre las que España ejerce derechos de explotación de recursos naturales del suelo y subsuelo marino, sin perjuicio —fórmula ya conocida del Tema 16— de los regímenes forales de Concierto y Convenio en el País Vasco y Navarra. [Probable] La diferencia relevante con el IRPF no está tanto en esta declaración territorial general como en el criterio que determina la sujeción: mientras en el IRPF la residencia habitual de una persona física se fija mediante criterios empíricos de presencia física (los 183 días, el núcleo de intereses económicos), en el IS la "residencia" de una persona jurídica es, por definición, una construcción normativa, porque una sociedad no "vive" en un lugar en el sentido en que lo hace una persona física —de ahí que el criterio decisivo, y el más disputado en la práctica internacional, sea el de la sede de dirección efectiva.

[Seguro] El art. 8 LIS considera residentes en territorio español a las entidades en las que concurra alguna de estas tres circunstancias —basta una sola, como en el IRPF—: que se hubieran constituido conforme a las leyes españolas; que tengan su domicilio social en territorio español; o que tengan en España su sede de dirección efectiva, entendida como el lugar donde radique la dirección y control del conjunto de sus actividades. [Probable] Este tercer criterio es el que genera mayor litigiosidad internacional, porque es el que persiguen combatir las estructuras de planificación fiscal agresiva: una sociedad puede constituirse formalmente en una jurisdicción de baja tributación y, sin embargo, ser residente fiscal en España si es aquí donde efectivamente se toman las decisiones estratégicas y de gestión —lo que exige a la Administración probar hechos de gestión real, no solo formalidades registrales, con la consiguiente dificultad probatoria que este tipo de comprobaciones plantea en la práctica. (OJO: este es el terreno donde se libra, en su vertiente más técnica, buena parte del debate sobre erosión de bases y traslado de beneficios —BEPS— que desarrollaremos en el apartado 2; la sede de dirección efectiva es, en el fondo, el mismo problema que el establecimiento permanente digital, aplicado a estructuras societarias tradicionales en lugar de a plataformas tecnológicas.)

2.3. El contribuyente

[Seguro] Son contribuyentes del IS las personas jurídicas, con la excepción expresa de las sociedades civiles sin objeto mercantil —que, como ya vimos en el Tema 16, tributan en régimen de atribución de rentas, imputando su resultado directamente a los socios en el IRPF o en el propio IS de estos, según su naturaleza—. Pero, igual que en el IRPF encontrábamos entidades sin personalidad jurídica sometidas a un régimen especial (las entidades en atribución de rentas), en el IS encontramos el fenómeno inverso: entidades carentes de personalidad jurídica propia que, sin embargo, sí tributan de forma autónoma por este impuesto, precisamente porque la ley opta por tratarlas como centros de imputación de renta independientes de sus partícipes. [Seguro] El art. 7 LIS enumera, entre otras, los fondos de inversión, las uniones temporales de empresas (UTEs), los fondos de capital-riesgo, los fondos de pensiones, los fondos de regulación del mercado hipotecario, y las comunidades titulares de montes vecinales en mano común —esta última una figura jurídica muy singular, de arraigo en el derecho consuetudinario gallego y asturiano, que ilustra bien hasta qué punto el legislador ha tenido que resolver, caso por caso, situaciones de titularidad colectiva atípica que no encajan ni en el molde de la persona jurídica ni en el de la atribución de rentas del IRPF.

2.4. Hecho imponible y periodo impositivo

[Seguro] El art. 4 LIS define el hecho imponible en términos deliberadamente paralelos a los del IRPF: la obtención de renta por el contribuyente, cualquiera que fuese su fuente u origen. [Probable] Pero aquí aparece una diferencia estructural de primer orden que conviene fijar ya, porque va a condicionar todo el apartado 2 del tema: el IRPF, como vimos con detalle en el Tema 16, clasifica la renta por fuentes (trabajo, capital, actividades económicas) precisamente porque esa clasificación determina un tratamiento fiscal diferenciado —la dualización renta general/renta del ahorro—. El IS, en cambio, no distingue la renta societaria por fuentes a efectos de su integración en la base imponible: un beneficio de explotación, un dividendo recibido de otra sociedad o una plusvalía por venta de un activo se integran, en principio, de forma unitaria en la misma magnitud, sin la arquitectura dual que estructuraba buena parte del Tema 16. [Probable] Esto no significa que el origen de la renta sea irrelevante —ya veremos que la exención de dividendos y plusvalías del art. 21 LIS depende crucialmente de su origen—, sino que esa relevancia opera por la vía de exenciones específicas y no mediante una clasificación general de fuentes como la del IRPF.

[Seguro] El periodo impositivo (art. 27 LIS) es el segundo punto de divergencia notable respecto del IRPF: coincide con el ejercicio económico de la entidad, que la sociedad fija libremente en sus estatutos y que no tiene por qué coincidir con el año natural —a diferencia del IRPF, donde el periodo impositivo es siempre el año natural salvo por fallecimiento del contribuyente—. El único límite es que el periodo impositivo no puede exceder de doce meses, y el impuesto se devenga el último día del periodo impositivo, con independencia de cuál sea esa fecha. [Suposición, dato de interés expositivo] Esto explica por qué, en la práctica profesional, es habitual encontrar sociedades —especialmente filiales de grupos multinacionales— con ejercicios fiscales que no van de enero a diciembre sino, por ejemplo, de octubre a septiembre, para hacerlos coincidir con el ejercicio de la matriz extranjera: una posibilidad sencillamente inexistente en el IRPF, donde el periodo impositivo no es nunca objeto de elección del contribuyente.

2.5. Gestión y aplicación del impuesto

[Seguro] La gestión corresponde, igual que en el IRPF, a la AEAT (salvo especialidad foral), mediante autoliquidación en el Modelo 200, que debe presentarse, conforme al art. 124.1 LIS, en el plazo de los 25 días naturales siguientes a los seis meses posteriores a la conclusión del periodo impositivo —para una entidad con ejercicio coincidente con el año natural, esto sitúa el plazo de declaración en torno al 25 de julio del año siguiente, sensiblemente más tarde que el plazo de la Renta del IRPF—. [Seguro] Como en el IRPF, existen pagos a cuenta, articulados aquí como pagos fraccionados (art. 40 LIS, Modelo 202): tres pagos, en los primeros veinte días naturales de abril, octubre y diciembre de cada año, calculados según dos modalidades alternativas —un porcentaje sobre la cuota del último periodo impositivo liquidado, o un porcentaje directo sobre el resultado contable positivo del periodo transcurrido en el propio ejercicio (art. 40.3 LIS), esta segunda modalidad obligatoria para las entidades cuyo importe neto de cifra de negocios supere los 6 millones de euros—. [Probable] Este umbral de 6 millones de euros es la primera aparición, en este marco jurídico general, de un criterio de tamaño empresarial que va a reaparecer de forma constante en el apartado 3: buena parte del régimen del IS no es uniforme para todas las entidades, sino que se segmenta según magnitudes de cifra de negocios (6 millones para la modalidad de pago fraccionado, 10 millones para el régimen de empresas de reducida dimensión, 20 millones para la tributación mínima del art. 30 bis que vimos en el bloque histórico) — una arquitectura por tramos de tamaño sin equivalente en el IRPF, donde la segmentación relevante era casi siempre por nivel de renta personal, no por magnitud de la actividad económica desarrollada.

[Probable] La competencia normativa autonómica, a diferencia del IRPF, es aquí prácticamente inexistente: el IS no es un tributo cedido en el sentido de la Ley 22/2009, y las Comunidades Autónomas de régimen común carecen, salvo bonificaciones muy tasadas en algún régimen especial, de capacidad para modular su tarifa o su base. [Seguro] La única dispersión normativa relevante es, de nuevo, la foral: País Vasco y Navarra mantienen tipos de gravamen y regímenes de deducción propios y sustancialmente distintos del territorio común —con el mismo esquema de tres Normas Forales dentro del País Vasco que ya vimos en el Tema 16 a propósito del IRPF—, con una diferencia de enfoque relevante que trataremos en el apartado 3: aquí la competencia fiscal entre territorios no se libra tanto por atraer contribuyentes-personas físicas como por atraer la sede social y la actividad productiva de las empresas, lo que convierte esta dispersión foral en un terreno de disputa política y económica de primer orden en el debate sobre la localización empresarial en España.

\subsection{Esquema básico de liquidación}

Cierro este apartado con el esquema que, igual que en el Tema 16, servirá de columna vertebral para el resto del tema —y aquí es donde el contraste con el IRPF se hace más visible:

Resultado contable, determinado conforme al Código de Comercio y al Plan General de Contabilidad, tal y como resulta de la cuenta de pérdidas y ganancias del ejercicio.
Ajustes extracontables (positivos y negativos), previstos por la propia LIS para corregir el resultado contable allí donde la norma fiscal se aparta deliberadamente del criterio contable —diferencias que la doctrina clasifica en permanentes (nunca revierten; por ejemplo, una multa contabilizada como gasto pero fiscalmente no deducible) y temporarias (revierten en ejercicios futuros; por ejemplo, una provisión contable no deducible en el momento de su dotación pero sí en el momento en que se aplica).
Base imponible previa, resultado de aplicar esos ajustes al resultado contable.
Compensación de bases imponibles negativas de ejercicios anteriores, con los límites cuantitativos que analizaremos en el apartado 2.
Base imponible, tras esa compensación.
Tipo de gravamen (general del 25%, con tipos especiales para determinadas entidades y regímenes que iremos viendo).
Cuota íntegra.
Deducciones para evitar la doble imposición (internacional e interna) y bonificaciones.
Cuota líquida, con el suelo mínimo del 15% ya visto para grandes empresas (art. 30 bis LIS).
Deducciones por incentivos (I+D+i, producciones cinematográficas, creación de empleo, entre otras que veremos en el apartado 3).
Retenciones, ingresos a cuenta y pagos fraccionados, minorando la cuota resultante.
Cuota diferencial, resultado final a ingresar o a devolver.

El punto 2 —el ajuste extracontable— es, para este tema, lo que el mínimo personal y familiar y la base del ahorro fueron para el Tema 16: el elemento que, presentado aquí de forma anticipada y sin desarrollo, va a estructurar por completo el apartado 2. Adelanto ya la idea que lo sustenta, sin desarrollarla: si en el IRPF el benchmark analítico era un ideal económico externo a la propia ley (SHS), en el IS el "benchmark" natural para medir cualquier divergencia fiscal es la propia norma contable — lo que cambia radicalmente la naturaleza del debate teórico que viene a continuación, y es la razón por la que te advertía, al empezar este tema, que la analogía con el Tema 16 se rompe justo en este punto.



\clearpage
\section{IS: Determinación de la base imponible}
\puntosclave{
    
}

Este apartado no puede arrancar,  situando un benchmark económico externo y midiendo distancias respecto de él. Aquí la primera pregunta es previa y más incómoda: ¿por qué gravar a una persona jurídica en absoluto? Una sociedad no tiene, en sentido estricto, capacidad económica propia distinta de la de las personas físicas que están detrás de ella; toda la renta que genera una sociedad termina, tarde o temprano, en el bolsillo de alguna persona física. 

Si el principio de capacidad económica del art. 31.1 CE se predica, en última instancia, de personas, ¿qué justifica interponer un impuesto autónomo sobre una ficción jurídica antes de que esa renta llegue a su destinatario final? Esta pregunta, es el punto de partida obligado y de su respuesta dependen directamente los dos grandes debates que cierran el bloque.

\subsection{De la Teoría al Método: ¿Por qué gravar la persona jurídica?}
La doctrina de hacienda pública ofrece, de forma clásica, dos grandes familias de respuesta a esta pregunta, sistematizadas con particular claridad por MUSGRAVE en su obra de referencia sobre teoría de la hacienda pública (\textit{Public Finance in Theory and Practice}).

La primera es la teoría del beneficio o de la contraprestación: la sociedad, como entidad jurídica separada, se beneficia directamente de bienes y servicios públicos —seguridad jurídica para hacer cumplir sus contratos, infraestructuras para transportar sus mercancías, un sistema judicial que protege su personalidad jurídica y su responsabilidad limitada frente a terceros— con independencia de quiénes sean, en cada momento, sus accionistas. Bajo esta lógica, el IS no es un anticipo de nada: es la contribución propia y autónoma que la sociedad debe al Estado por el privilegio, nada trivial, de operar bajo el escudo de la responsabilidad limitada, un privilegio que el Estado concede y que tiene un coste social —el riesgo trasladado a los acreedores— que conviene internalizar de algún modo.

La segunda es la teoría de la retención en la fuente o del anticipo del IRPF: el IS no tendría, bajo esta lectura, justificación autónoma alguna; existiría sencillamente porque es mucho más eficaz recaudar el impuesto en la fuente —en el momento en que la sociedad genera el beneficio— que esperar a perseguir esa misma renta cuando llegue, mucho más fragmentada y mucho más difícil de fiscalizar, a los cientos o miles de accionistas dispersos que la reciben en forma de dividendo. Esta segunda teoría explica, mejor que la primera, un fenómeno que citamos ya en el bloque histórico sin desarrollarlo: la propia Exposición de Motivos de la Ley 43/1995 describe el IS como cumpliendo, entre sus funciones, la de "función de retención en la fuente respecto de las rentas del capital obtenidas por los inversores extranjeros a través de sociedades de su propiedad residentes en territorio español" —una descripción que el legislador español hace de sí mismo, y que encaja con precisión casi literal en la teoría del anticipo: sin el IS, gran parte de la renta generada en España por inversores no residentes escaparía sencillamente al control fiscal español, porque el IRNR sobre no residentes sin establecimiento permanente tiene limitaciones de alcance que el IS, actuando en la fuente, no tiene.

Ninguna de las dos teorías, por separado, explica satisfactoriamente el diseño real del impuesto. Si el IS fuera solo un anticipo del IRPF de los socios, debería existir un mecanismo de devolución o de imputación plena cuando el socio es persona física residente y tributa efectivamente por ese dividendo; y si el IS respondiera únicamente a la teoría del beneficio, no se explicaría por qué su recaudación varía tan sensiblemente con la cifra de beneficios de la entidad y no, por ejemplo, con una medida más directa de su uso de bienes públicos. [Suposición] La lectura más defendible, y la que uso como marco de trabajo para el resto del apartado, es que el IS español —como la mayoría de sus homólogos comparados— es un híbrido pragmático entre ambas justificaciones, cuyo peso relativo varía según el tipo de socio: frente a un accionista persona física residente, el IS opera predominantemente como anticipo del IRPF (de ahí la preocupación por la doble imposición que veremos); frente a un accionista no residente o frente a la propia sociedad como agente económico autónomo, opera predominantemente como tributo con justificación propia de teoría del beneficio.

\subsubsection{Benchmark: ¿por qué aquí no hay un SHS societario?}
Conviene ahora zanjar explícitamente algo que ya adelanté al empezar el tema: no existe, en la literatura de hacienda pública, un equivalente consensuado del concepto Schanz-Haig-Simons aplicado a la renta societaria. La razón no es casual: SHS mide capacidad económica de una persona, y una sociedad no tiene, en el sentido relevante para la equidad tributaria, capacidad económica propia que "medir" con independencia de sus socios —de ahí precisamente la pregunta que acabamos de plantear en el epígrafe anterior—. [Seguro] Por eso el benchmark que vamos a usar en este tema, ya anticipado en el apartado 1, es de naturaleza distinta: no un ideal económico de renta, sino la norma contable, es decir, el resultado que arroja la aplicación correcta y de buena fe del Código de Comercio y del Plan General de Contabilidad. Esta elección de benchmark no es arbitraria ni mía: es la que adopta la propia Ley 27/2014 en su art. 10.3, al ordenar que la base imponible se calcule corrigiendo el resultado contable, no calculándose de forma autónoma según un ideal externo.

[Probable] Esta diferencia de benchmark tiene una consecuencia metodológica importante para todo lo que sigue, y es la que ya hemos usado para resolver tu pregunta anterior: mientras que en el Tema 16 el criterio para distinguir "estructural" de "gasto fiscal" era la distancia respecto de SHS, aquí el criterio equivalente es la distancia respecto de lo que exigiría una correcta medición del resultado económico real de la entidad a lo largo de su existencia —lo que en la literatura anglosajona se resume en la fórmula ya citada de Gutman: gravar la renta que la empresa ha obtenido "a lo largo de su existencia, sin atender a la imposición artificial, aunque necesaria, de un periodo de declaración anual". Un ajuste extracontable que corrige una diferencia meramente temporal entre contabilidad y fiscalidad —amortización, provisión, compensación de pérdidas— encaja en ese benchmark corregido y es, por tanto, estructural. Un ajuste que premia una conducta —invertir en I+D+i, capitalizar la empresa, rodar una película en España— no corrige nada: añade un incentivo nuevo sobre la base ya correctamente medida, y es, por tanto, gasto fiscal en el sentido estricto de Surrey.

1.3. Roadmap de las divergencias que se analizan

Fijado el porqué del impuesto y el benchmark de referencia, el resto del apartado se organiza en tres bloques: primero, la divergencia temporal entre el devengo contable y el devengo fiscal, con la compensación de bases imponibles negativas ya situada como su pieza central; segundo, la divergencia por atribución, es decir, los mecanismos que aseguran que una misma renta tribute una sola vez dentro de un grupo societario o en operaciones entre partes vinculadas; y tercero, los dos grandes debates abiertos que la propia existencia del impuesto deja sin cerrar del todo —quién soporta realmente su carga, y cómo se corrige, o no, la doble imposición del dividendo—, que son, en realidad, la culminación natural de la pregunta que hemos planteado en este primer subapartado y no una cuestión distinta de ella.

\subsection{Divergencias temporales: contabilidad y fiscalidad}
\textcolor{red}{Antes de desarrollar el catálogo, quiero anclarlo con un hallazgo que cambia el tono de todo el subapartado: los límites a la compensación de BINs no son solo un problema de diseño técnico —son, además, el objeto de una de las declaraciones de inconstitucionalidad más relevantes de la última década en materia tributaria, y por una razón que conecta directamente con el argumento estructural que defendíamos en el subapartado anterior. Lo desarrollo en su lugar.}

\begin{specialblock}{¿Por qué la compensación de rentas anteriores es una divergencia temporal y no un beneficio fiscal?}
    Confirmado con una fuente que no deja margen de duda: el Joint Committee on Taxation estadounidense —la autoridad de referencia en metodología de tax expenditures, la misma tradición que arranca con Surrey— excluye explícitamente las compensaciones de pérdidas (carryback/carryforward) de su lista de gastos fiscales, calificándolas como parte de la "normal income tax law", precisando además que incluso los límites temporales o cuantitativos que se les impongan se presumen, salvo prueba en contrario, parte de esa misma normalidad estructural, justificados por razones de conveniencia administrativa. La razón de fondo, bien resumida por Hank Gutman —antiguo Chief of Staff del propio JCT— es esta: el impuesto sobre la renta se basa en periodos anuales artificiales pero necesarios por razones de gestión; sin un mecanismo de compensación, una empresa que gana 100 un año y pierde 100 al siguiente —con una renta vital neta de cero— acabaría pagando impuesto sobre una renta que nunca tuvo, solo por el azar de en qué ejercicio concreto cayó el beneficio. La compensación no es una gracia del legislador: es la corrección necesaria de una distorsión que el propio impuesto introduce al trocear en años una actividad económica continua.

    Mantengo, por tanto, la compensación de BINs en el apartado 2, como manifestación —de hecho, la más pura posible— de la divergencia temporal: es exactamente el mismo fenómeno que el criterio de realización que vimos en el Tema 16, aplicado ahora no a un activo individual sino a la renta total de una entidad a lo largo de su vida. Sus límites cuantitativos (el porcentaje máximo de base imponible positiva que puede absorberse cada año, hoy en torno al 60-70% con un suelo de un millón de euros siempre compensable) no la convierten en beneficio fiscal: son, igual que el límite del 25% de compensación de pérdidas de capital que vimos en el IRPF, una restricción recaudatoria sobre un principio estructural, no una concesión.

    Dicho esto, tu intuición señala un problema real que sí debo corregir, y es el siguiente: en la literatura que acabo de revisar aparece, junto a la compensación de BINs, la figura de la reserva de capitalización (art. 25 LIS) —que ya tenía prevista para el apartado 3— descrita expresamente como "beneficio fiscal novedoso". Y tiene toda la razón esa calificación: a diferencia de la compensación de pérdidas, la reserva de capitalización no corrige ninguna distorsión temporal —no está devolviendo al contribuyente algo que el impuesto le quitó de más—, sino que premia una conducta que el legislador quiere fomentar (incrementar fondos propios en vez de financiarse con deuda). Misma ubicación formal —ambas son "reducciones en la base imponible"—, naturaleza radicalmente distinta. Aprovecho tu pregunta para blindar el criterio de corte con un caso más de la misma familia, que si no lo señalo ahora corre el riesgo de colárseme más adelante: la libertad de amortización (distintos supuestos del art. 12 LIS y regímenes especiales, típicamente para empresas de reducida dimensión con creación de empleo, o para determinados activos de I+D+i) tampoco es una corrección estructural —a diferencia de las tablas de amortización ordinarias, que sí aproximan la depreciación económica real del activo—, sino un incentivo deliberado a adelantar la deducción por razones de política de inversión. Ambas, reserva de capitalización y libertad de amortización, quedan confirmadas para el apartado 3; la compensación de BINs y las tablas ordinarias de amortización, para el apartado 2.
\end{specialblock}


Conviene fijar primero por qué existe esta divergencia, porque a diferencia del IRPF —donde el criterio de realización nace de una elección deliberada del legislador frente al devengo económico puro— aquí el problema es de otra naturaleza: la contabilidad ya opera, desde la Ley 43/1995, bajo un criterio de devengo económico razonablemente fiel a la realidad de la empresa. El problema no es que la fiscalidad "prefiera" un criterio distinto al contable por convicción propia, sino que la fiscalidad desconfía, en ciertos puntos muy concretos, de la fiabilidad o la neutralidad recaudatoria del juicio contable, y ajusta en consecuencia. Es una desconfianza selectiva, no sistemática: la inmensa mayoría del resultado contable se traslada sin corrección alguna a la base imponible; los ajustes se concentran en un número relativamente reducido de partidas donde el legislador ha decidido que el criterio contable, por su naturaleza estimativa o valorativa, no debe determinar directamente cuánto impuesto se paga.

2.2. Amortizaciones: coeficientes fiscales frente a vida útil económica

[Seguro] El art. 12 LIS exige que la amortización contabilizada como gasto sea, además, fiscalmente deducible, lo que en la práctica significa que debe respetar los coeficientes máximos y los periodos máximos de amortización recogidos en la tabla del propio artículo —una tabla notablemente simplificada por la Ley 27/2014 respecto de la, mucho más detallada, tabla reglamentaria anterior a 2015—. [Probable] El criterio contable de fondo (norma de valoración 2ª del PGC) exige amortizar el activo conforme a su vida útil económica real, estimada por la propia empresa según criterios técnicos; el criterio fiscal, en cambio, fija topes objetivos por categorías de activos, sin atender a las circunstancias singulares de cada empresa concreta. Cuando ambos criterios coinciden —lo habitual—, no hay ajuste. Cuando la empresa amortiza contablemente más deprisa de lo que permite la tabla fiscal —por ejemplo, porque técnicamente considera que un activo tiene una vida útil más corta que la que asume la Hacienda Pública de forma genérica para toda esa categoría de bienes—, surge un ajuste extracontable positivo en el ejercicio del exceso, que revertirá como ajuste negativo en ejercicios posteriores, cuando el activo ya esté contablemente amortizado pero fiscalmente aún queden cuotas pendientes de deducir dentro del margen que permite la tabla. Es una diferencia temporaria en el sentido técnico-contable del término: no cambia el importe total deducible a lo largo de la vida del activo, solo el ejercicio en que se deduce.

Distingo aquí, con toda intención, este mecanismo de la libertad de amortización que ya identificamos como gasto fiscal genuino en tu pregunta anterior: la tabla ordinaria del art. 12 LIS no premia ninguna conducta —solo fija, de forma objetiva y general, un tope de prudencia recaudatoria a lo que en otro caso sería una estimación subjetiva de la propia empresa—; los regímenes de libertad de amortización, en cambio, permiten deducir el coste íntegro del activo en el primer ejercicio, muy por delante de cualquier criterio de depreciación económica real, precisamente para incentivar la inversión en determinados activos o por determinadas empresas. La frontera entre ambos no está en su ubicación normativa —ambos aparecen formalmente como "amortización"— sino en si corrigen una estimación o si conceden una ventaja.

2.3. Provisiones, deterioros, y el caso singular del deterioro de participaciones

[Seguro] El art. 13 LIS somete a condiciones estrictas la deducibilidad fiscal de las pérdidas por deterioro de créditos derivadas de insolvencias de deudores —exige, entre otros requisitos alternativos, que hayan transcurrido al menos seis meses desde el vencimiento de la obligación, o que el deudor esté declarado en concurso, o que se hayan iniciado acciones judiciales de reclamación—, mientras que la contabilidad puede reconocer el deterioro antes, en el momento en que la dirección de la empresa estima, con criterios técnicos, que la recuperación del crédito es dudosa. [Probable] La lógica fiscal aquí no es de incentivo sino de verificabilidad: un deterioro contable es, en parte, un juicio de valor de la propia empresa sobre sus propios activos, y el legislador fiscal exige un hecho externo y objetivamente comprobable —el transcurso de un plazo, una declaración judicial— antes de permitir que ese juicio reduzca la recaudación. Es, de nuevo, una diferencia temporaria: el gasto acabará siendo deducible, solo que más tarde de lo que la contabilidad lo reconoce.

[Seguro] Distinto es el caso, más drástico, del deterioro de participaciones en el capital de otras entidades: desde la reforma operada por la Ley 16/2013, de 29 de octubre, en pleno contexto de consolidación fiscal post-crisis, este tipo de deterioro dejó de ser fiscalmente deducible en ningún caso, con independencia de que la participada hubiera perdido valor real y de forma verificable. [Probable] Aquí ya no hablamos de una diferencia temporaria sino de una diferencia permanente: el gasto contable nunca llega a ser fiscalmente deducible, ni antes ni después. Es un ejemplo interesante para el debate de fondo del apartado, porque roza la frontera entre lo estructural y lo recaudatorio puro: no corrige ninguna distorsión temporal —no hay "reversión" futura que compense la no deducción—, así que en rigor no encaja en la lógica de "corrección de una desviación introducida por el propio impuesto" que hemos usado para defender el carácter estructural de la compensación de BINs. [Suposición] Es, más bien, una medida de ensanchamiento de base por pura necesidad recaudatoria en tiempos de crisis —coherente con el patrón que ya vimos repetidas veces en el Tema 16 a propósito de la Ley 26/2014 del IRPF—, y aunque formalmente aparece aquí, entre las diferencias permanentes, conviene que la trates con matiz: no es un beneficio fiscal (no premia ninguna conducta), pero tampoco es una corrección estructural en el sentido estricto que hemos defendido para las BINs; es, si acaso, el reverso de un gasto fiscal —una restricción de base sin contrapartida de incentivo alguno.

2.4. Instrumentos financieros a valor razonable: el eco del criterio de realización

[Probable] El art. 15.k) LIS establece, con carácter general, que no son fiscalmente deducibles las disminuciones de valor originadas por aplicación del criterio del valor razonable correspondientes a valores representativos de participaciones en el capital o fondos propios de entidades, salvo que se integren en la base imponible a través de reglas específicas de excepción —fundamentalmente, cuando se trate de instrumentos financieros mantenidos para negociar, donde el propio ciclo de negocio del activo justifica seguir el criterio contable—. [Suposición] Es, en su núcleo conceptual, el mismo fenómeno que analizamos en el Tema 16 a propósito de las ganancias patrimoniales del IRPF: la contabilidad puede reconocer, año a año, la variación de valor de un activo financiero cotizado sin que se haya producido ninguna transmisión, mientras que la fiscalidad, con carácter general, prefiere esperar a que esa variación se materialice en un negocio jurídico real antes de darle trascendencia tributaria —el mismo criterio de realización, trasplantado del ámbito de la persona física al de la persona jurídica, con la particularidad de que aquí compite directamente con un criterio contable de devengo que sí reconoce esas variaciones no realizadas, generando por tanto un ajuste extracontable explícito que en el IRPF ni siquiera hace falta, porque allí no hay punto de partida contable del que discrepar.

2.5. La compensación de bases imponibles negativas: mecánica, evolución y un episodio de inconstitucionalidad

[Seguro] El art. 26 LIS permite compensar las bases imponibles negativas generadas en un periodo impositivo con las rentas positivas de los periodos siguientes, sin límite temporal alguno desde la reforma de 2014 —a diferencia del régimen anterior, que había ido ampliando progresivamente el plazo de compensación (10 años bajo la Ley 43/1995 original, 15 años desde 2001, 18 años en plena crisis) hasta que la Ley 27/2014 eliminó por completo el límite temporal—. [Probable] Esta eliminación del límite temporal parece, a primera vista, una liberalización del régimen; pero la misma reforma introdujo, por primera vez, un límite cuantitativo que no existía bajo el régimen de plazo limitado: la compensación no puede superar el 70% de la base imponible previa a la propia compensación (y a la reserva de capitalización), con un suelo de un millón de euros siempre compensable con independencia de ese porcentaje. [Seguro] Este 70% se redujo transitoriamente al 60% para el ejercicio 2016 —Ley de Presupuestos para 2015—, y el Real Decreto-ley 3/2016, de 2 de diciembre, introdujo además, solo para grandes empresas (cifra de negocios igual o superior a 20 millones de euros), límites todavía más estrictos: 50% para empresas entre 20 y 60 millones, y 25% para empresas por encima de 60 millones —límites que, junto con una restricción paralela a las deducciones por doble imposición y con la obligación de reintegrar en la base imponible, de forma automática y fraccionada en cinco años, deterioros de participaciones deducidos en ejercicios anteriores, formaban parte de un mismo paquete de medidas de consolidación fiscal urgente aprobado por decreto-ley.

[Seguro] Este último dato —aprobado por decreto-ley— es el que determina el desenlace del episodio. La STC 11/2024, de 18 de enero (BOE de 20 de febrero de 2024), resolviendo una cuestión de inconstitucionalidad planteada por la Audiencia Nacional, declaró inconstitucionales y nulos estos tres preceptos del Real Decreto-ley 3/2016 —los límites reforzados a la compensación de BINs para grandes empresas, el nuevo límite a las deducciones por doble imposición, y la reincorporación automática de deterioros—, por vulneración del art. 86.1 CE: un decreto-ley, instrumento normativo reservado por la Constitución a situaciones de extraordinaria y urgente necesidad, no puede afectar a los derechos, deberes y libertades de los ciudadanos regulados en el Título I, entre los que se encuentra expresamente el deber de contribuir al sostenimiento de los gastos públicos del art. 31.1 CE. [Seguro] El Tribunal razonó que estas medidas alteraban elementos esenciales del impuesto —base imponible y cuota— con un impacto recaudatorio notable sobre una figura básica del sistema tributario, lo que basta para situarlas fuera del ámbito material que un decreto-ley puede regular, con independencia de que el contenido concreto de la medida —limitar más una compensación de pérdidas para grandes empresas— pudiera ser, en sí mismo, una opción legítima de política fiscal si se hubiera tramitado como ley ordinaria. [Seguro] El Tribunal conecta expresamente esta sentencia con un precedente idéntico: la STC 78/2020, que ya había declarado inconstitucional, por el mismo motivo, el Real Decreto-ley 2/2016 —la norma "gemela" que, meses antes, había incrementado por la misma vía los pagos fraccionados del IS de las grandes empresas—.

[Probable] Conviene que extraigas de este episodio dos lecturas, no solo una. La primera, la más evidente, es formal: confirma un límite constitucional real y recurrente al uso del decreto-ley en materia tributaria, que cualquier Gobierno en apuros de consolidación fiscal debe respetar, y que ya habíamos visto operar, con otro objeto, a propósito de la tributación conjunta del IRPF en el Tema 16. La segunda, más sustantiva y más relevante para el hilo argumental de este apartado, es que el propio Tribunal Constitucional está reconociendo, con otras palabras, exactamente el argumento que hemos defendido para mantener la compensación de BINs en el terreno estructural: describe estos límites reforzados como una alteración de "elementos esenciales" que "inciden en la determinación de la carga tributaria" de una "pieza básica del sistema tributario" — es decir, el propio Tribunal trata la compensación de pérdidas como parte del núcleo estructural del impuesto, no como un beneficio fiscal periférico y prescindible que el Gobierno pudiera limitar sin mayor trámite. [Suposición, nexo pedagógico] Si te preguntan en el ejercicio oral por qué defiendes el carácter estructural de la compensación de BINs, este es el argumento de autoridad más sólido que puedes invocar: no es solo la opinión del JCT estadounidense, es la propia doctrina del Tribunal Constitucional español, llegada por una vía completamente distinta —el control del uso del decreto-ley— pero convergente en el fondo.

[Suposición, dato comparado de interés expositivo] Para cerrar con perspectiva comparada: Francia limita la compensación de pérdidas al 50\% de la base imponible que exceda de un millón de euros, sin límite temporal; Alemania aplica una lógica de "cláusula mínima" (Mindestbesteuerung) de estructura similar a la española. España no es, por tanto, una anomalía en el contexto comparado —el propio diseño de "porcentaje con suelo mínimo garantizado" es, de hecho, el estándar europeo habitual—; lo que sí resultó anómalo, y por eso fue anulado, fue el vehículo normativo elegido para endurecerlo, no el endurecimiento en sí.

2.6. La neutralidad fiscal en las reestructuraciones empresariales

[Seguro] El régimen especial del Capítulo VII LIS (arts. 76 a 89), aplicable a fusiones, escisiones, aportaciones de rama de actividad, canjes de valores y cesiones globales de activo y pasivo, resuelve exactamente el mismo problema que la compensación de BINs, solo que proyectado sobre plusvalías latentes en vez de sobre pérdidas: si una empresa se reorganiza jurídicamente —por ejemplo, si dos sociedades del mismo grupo se fusionan en una sola, o si una rama de actividad se segrega a una filial nueva—, la aplicación estricta de las reglas generales de transmisión patrimonial obligaría a tributar de inmediato por la diferencia entre el valor de mercado y el valor fiscal de todos los activos transmitidos en la operación, aunque no haya habido ningún cambio real en quién controla económicamente el negocio. [Probable] El régimen de neutralidad evita precisamente eso: la entidad adquirente se subroga en los valores y fechas de adquisición fiscales de la entidad transmitente —exactamente el mecanismo de subrogación que ya conoces, planteado en otros términos, del art. 36 LIRPF y los pactos sucesorios del Tema 16—, de modo que la plusvalía no desaparece, sino que se traslada intacta hacia el futuro, hasta que el activo se transmita de verdad frente a un tercero ajeno a la operación de reestructuración. Es diferimiento por realización, en estado puro, aplicado ahora no a un activo individual de un contribuyente persona física sino al patrimonio empresarial completo de una entidad jurídica.

[Seguro] El origen de este régimen es la Directiva 2009/133/CE (que codifica y sustituye a la originaria Directiva 90/434/CEE, de 1990), transpuesta por primera vez al ordenamiento español por la Ley 29/1991, e integrada después en el cuerpo general del IS por la Ley 43/1995. Como en el caso de las BINs, aquí también existe una salvaguarda recaudatoria frente al uso abusivo del régimen: el art. 89.2 LIS excluye la aplicación de la neutralidad cuando la operación se realice con la finalidad principal de fraude o evasión fiscal, exigiendo que exista un motivo económico válido distinto del mero ahorro fiscal —racionalización o reestructuración de actividades, por ejemplo— para justificar la reorganización. [Seguro] Esta cláusula, y la propia noción de "motivo económico válido" como estándar de control, no es una invención española: procede directamente del art. 15 de la propia Directiva de Fusiones, y su interpretación de referencia en toda la Unión Europea es la del Tribunal de Justicia en el asunto Leur-Bloem (Sentencia de 17 de julio de 1997, C-28/95), donde el Tribunal precisó que los Estados miembros no pueden presumir automáticamente la existencia de fraude por el mero hecho de que la operación no responda a una actividad "normal" u ordinaria, sino que la Administración debe examinar caso por caso el conjunto de la operación. [Probable] El paralelismo con lo que ya vimos en el subapartado 5 —la aparición reiterada de cláusulas anti-abuso allí donde el legislador introduce un mecanismo estructural de diferimiento— confirma que este patrón no es casual: cada vez que el impuesto renuncia a gravar de inmediato una plusvalía no realizada, aparece, casi como contrapeso automático, una cláusula de vigilancia que impide que ese diferimiento se convierta, por la vía de una reorganización artificial, en una exención encubierta.

[Suposición, nexo con la compensación de BINs] Un último punto que conecta directamente con el subapartado 5 y que conviene no perder: el art. 84 LIS regula la subrogación de la entidad adquirente no solo en los valores fiscales de los activos transmitidos, sino también en los derechos y obligaciones tributarias de la entidad transmitente, lo que incluye, expresamente, el derecho a compensar las bases imponibles negativas pendientes que esta tuviera generadas. Esta regla, aparentemente técnica, fue durante años objeto de un uso abusivo bien documentado en la práctica española —la llamada "compraventa de sociedades con bases imponibles negativas", donde se adquiría una sociedad inactiva pero con un importante crédito fiscal acumulado, únicamente para aprovechar ese crédito en una reestructuración posterior—, lo que llevó al legislador a introducir límites específicos a esta subrogación cuando la entidad transmitente se encuentra en una situación de inactividad prolongada o de cambio sustancial en su actividad.


\subsection{Divergencia por atribución: participaciones y operaciones vinculadas}

### 3.1. El hilo común: que la renta tribute una vez, y en el lugar correcto

Antes de separar estas dos figuras —que a primera vista pertenecen a mundos distintos, una a la fiscalidad de los grupos societarios puramente domésticos y otra a la planificación fiscal internacional— conviene fijar por qué las trato juntas. Ambas responden a la misma preocupación de fondo: evitar que la existencia de varias sociedades relacionadas entre sí —por participación de capital en un caso, por cualquier vínculo de control o parentesco en el otro— permita que una misma renta económica tribute más de una vez (doble imposición) o, en el extremo contrario, que no tribute en ningún sitio o tribute donde no debería (elusión mediante desplazamiento artificial de beneficios). Es exactamente el mismo problema, mirado desde sus dos caras: el art. 21 LIS resuelve el riesgo de "demasiado impuesto"; el art. 18 LIS y el marco BEPS resuelven el riesgo de "demasiado poco impuesto, en el sitio equivocado". No es casualidad que ambas figuras convivan en el mismo bloque temático de cualquier manual de la asignatura: son las dos mitades de una misma pregunta sobre neutralidad entre entidades vinculadas.

### 3.2. La exención de dividendos y plusvalías del art. 21 LIS: de la deducción a la exención, y de la exención plena a la exención del 95%

[Probable] Conviene situar primero el problema que resuelve, porque su relevancia económica es enorme y suele pasar desapercibida sin un ejemplo: imagina un grupo societario en cascada —una filial que reparte beneficios a su matriz, que a su vez es filial de una matriz superior, que a su vez reparte dividendo a sus accionistas—. Si cada eslabón de esa cadena tributase íntegramente por el dividendo recibido antes de repartirlo hacia arriba, la misma renta original quedaría sometida a tantas capas de IS como sociedades interpuestas existan en la estructura, con un efecto multiplicador brutal que penalizaría exclusivamente la organización empresarial en varias sociedades frente a la organización en una sola —una distorsión que ningún principio de neutralidad fiscal puede justificar, porque la forma jurídica de organización de un grupo empresarial no debería determinar, por sí sola, cuánto impuesto total paga ese grupo sobre el mismo beneficio económico.

[Seguro] España corrige este problema mediante un régimen de **exención**, regulado en el art. 21 LIS, aplicable tanto a los dividendos percibidos de otras entidades como a las plusvalías obtenidas en la transmisión de participaciones, siempre que se cumplan dos requisitos principales: una participación mínima del 5% en el capital de la entidad participada (mantenida, o con compromiso de mantenerla, durante al menos un año), y —cuando la participada es una entidad no residente— que esta haya estado sujeta a un impuesto de naturaleza análoga al IS con un tipo nominal mínimo del 10%, requisito que se entiende automáticamente cumplido si existe con ese país un convenio de doble imposición con cláusula de intercambio de información.

[Probable] Este diseño no es el que existía originalmente. Bajo la Ley 43/1995 y hasta la reforma de 2014, España corregía la doble imposición intersocietaria mediante dos regímenes separados y de naturaleza distinta: para participaciones domésticas, una **deducción en cuota** (método de imputación, deducción por doble imposición interna) que neutralizaba, mediante un crédito fiscal, el impuesto ya pagado por la filial; para participaciones internacionales, ya existía un régimen de exención similar al actual. La Ley 27/2014 unificó ambos regímenes bajo un único mecanismo de exención —abandonando definitivamente el método de imputación también para el ámbito doméstico—, en línea con la tendencia internacional mayoritaria: la exención es, administrativamente, mucho más sencilla que calcular y acreditar el impuesto exacto pagado en cada eslabón de la cadena societaria, un cálculo que se complica exponencialmente cuantas más sociedades interpuestas existan.

[Seguro] El régimen, sin embargo, dejó de ser una exención **plena** desde el 1 de enero de 2021: la **Ley 11/2020**, de Presupuestos Generales del Estado para 2021, introdujo un recorte del 5% —"en concepto de gastos de gestión no deducibles referidos a la participación"—, de modo que la exención efectiva pasó del 100% al **95%** del dividendo o de la plusvalía. [Probable] Esta medida no fue una ocurrencia española: la propia Directiva matriz-filial de la UE (Directiva 2011/96/UE) permite expresamente a los Estados miembros introducir este recorte a tanto alzado, con el límite del 5%, y países como Francia, Bélgica o Alemania ya lo habían incorporado con anterioridad a sus respectivas transposiciones. [Probable] La justificación oficial —una estimación objetiva y a tanto alzado de los gastos de gestión asociados a mantener la participación, gastos que en rigor nunca deberían haberse deducido íntegramente sin límite— convive con una crítica doctrinal bien documentada: al aplicarse ese 5% en **cada eslabón** de una cadena societaria con varios niveles, el efecto no es una única minoración del 5%, sino un **efecto cascada acumulativo** que erosiona el beneficio distribuido tantas veces como sociedades interpuestas existan entre la que genera la renta original y el accionista final —precisamente el mismo problema estructural de imposición en cascada que la exención pretendía evitar en 1995, reintroducido ahora, de forma deliberada y a menor escala, por razones puramente recaudatorias. [Suposición] Es un ejemplo interesante para tu exposición oral de cómo una medida de ensanchamiento de base puede, sin proponérselo, reabrir parcialmente el mismo problema estructural que el propio impuesto había resuelto una generación antes.

3.3. El grupo fiscal como unidad de tributación: la consolidación fiscal
[Probable] Si la exención del art. 21 LIS resolvía el problema de que un mismo beneficio tributara varias veces al circular en forma de dividendo entre los distintos eslabones de una cadena societaria, el régimen de consolidación fiscal resuelve un problema distinto pero de la misma familia: evitar que la mera decisión empresarial de organizar una actividad económica en varias sociedades jurídicamente distintas —en lugar de en una sola— penalice fiscalmente al grupo frente a lo que pagaría si esa misma actividad estuviera integrada en una única entidad.

[Seguro] El régimen, regulado hoy en el Capítulo VI del Título VII LIS (arts. 55 a 75), permite que un grupo fiscal —una entidad dominante y sus dependientes, participadas directa o indirectamente en al menos el 75% de su capital (70% si cotizan en un mercado regulado), y durante todo el periodo impositivo— tribute mediante una declaración única (Modelo 220), en la que la base imponible del grupo se determina como la suma algebraica de las bases imponibles individuales de cada entidad integrante. [Probable] La consecuencia inmediata, y la que da sentido económico al régimen, es que las pérdidas de una filial se compensan en el mismo ejercicio con los beneficios de otra filial del grupo, sin necesidad de esperar, como ocurriría tributando individualmente, a que cada entidad genere beneficios propios futuros contra los que aplicar sus propias bases imponibles negativas. A esto se suma un segundo mecanismo, más técnico: las eliminaciones e incorporaciones de resultados derivados de operaciones internas del grupo —por ejemplo, el margen de beneficio en una venta de existencias entre dos sociedades del mismo grupo— se neutralizan mientras el bien permanece dentro del grupo, y solo se incorporan a la base imponible consolidada cuando ese bien sale del grupo hacia un tercero independiente, de nuevo el mismo principio de diferimiento hasta la realización frente a un tercero ajeno que ya hemos visto operar, con variantes, en toda la divergencia temporal.

[Suposición, dato de interés expositivo] Este régimen tiene, además, una genealogía histórica singularmente antigua dentro del IS español, anterior incluso al propio impuesto moderno: se remonta al Real Decreto-ley 15/1977, de 25 de febrero, desarrollado por el Real Decreto 1414/1977 —es decir, nace bajo el régimen preconstitucional, meses antes incluso de la Ley 50/1977 de medidas urgentes de reforma fiscal que situamos al principio del tema—, y fue modificado sustancialmente por la Ley 18/1982, que redefinió el concepto de sociedad dominante, antes de su incorporación definitiva al cuerpo general del impuesto con la Ley 43/1995. [Probable] Conviene que sepas justificar por qué lo trato aquí, en la divergencia por atribución, y no como un beneficio fiscal del apartado 3: el régimen no concede ninguna ventaja que no se derivaría, en rigor, de medir correctamente la capacidad económica agregada de una unidad de decisión económica real —el grupo—, sino que corrige la distorsión que resultaría de tratar como compartimentos estancos a entidades que, aunque formalmente separadas, actúan y se financian como una sola. Es, en el fondo, la misma idea que sustentaba la exención del art. 21 LIS, aplicada ahora no a la circulación de un dividendo ya generado, sino a la propia determinación conjunta de la base imponible del ejercicio.


3.4. Operaciones vinculadas: el principio de plena competencia y su origen normativo
El segundo pilar de este bloque, las **operaciones vinculadas**, tiene un objeto distinto: no ya la renta que fluye entre matriz y filial en forma de dividendo, sino el **precio** al que se facturan bienes, servicios, préstamos, cesiones de uso de marcas o de propiedad intelectual, entre entidades que forman parte de un mismo grupo o que están relacionadas de cualquier otro modo tasado por la ley —el art. 18.2 LIS enumera, entre otros supuestos de vinculación, la relación entre una entidad y sus socios, entre una entidad y sus administradores (salvo por su retribución como tales), entre entidades del mismo grupo, y vínculos de parentesco hasta el tercer grado con socios o administradores—. [Seguro] El art. 18.1 LIS exige que estas operaciones se valoren, a efectos fiscales, por su **valor de mercado** —definido como el que habrían acordado partes independientes en condiciones de libre competencia (el llamado *arm's length principle*, principio de plena competencia, estándar internacional desarrollado y sistematizado por las Directrices de Precios de Transferencia de la OCDE)—, y no por el precio que las propias partes, precisamente por estar vinculadas y no enfrentarse a intereses contrapuestos, hayan podido pactar libremente entre sí.

[Seguro] Esta obligación de valorar a mercado no nace con la Ley 27/2014: existía ya, en forma más rudimentaria, desde el origen mismo del impuesto en 1978, pero fue la **Ley 36/2006, de 29 de noviembre**, de medidas para la prevención del fraude fiscal —la misma reforma, recordarás, que aparecía en el Tema 16 introduciendo el gravamen especial sobre premios de lotería y otras medidas antifraude— la que introdujo, por primera vez en el ordenamiento tributario español, la **obligación de documentar** la valoración a mercado realizada por el contribuyente, obligando a mantener esa documentación a disposición de la Administración. [Seguro] Esta obligación de documentación fue desarrollada reglamentariamente por el **Real Decreto 1793/2008**, y posteriormente simplificada —tras la crisis, con criterio de reducción de cargas administrativas para pymes— por el Real Decreto-ley 6/2010, que exoneró de la obligación de documentar a las empresas de reducida dimensión cuando el conjunto de sus operaciones vinculadas no superase los 100.000 euros anuales a valor de mercado; hoy, el propio art. 18.3 LIS mantiene un contenido documental simplificado para entidades con cifra de negocios inferior a 45 millones de euros.

### 3.5. De las operaciones vinculadas a BEPS: cuando el problema deja de ser doméstico

[Probable] Mientras la operación vinculada se produce dentro de un único territorio fiscal, el riesgo que corrige el art. 18 LIS es sobre todo de reparto interno de la carga tributaria entre socio y sociedad, o entre sociedades de un mismo grupo residentes en España. El problema se vuelve mucho más grave, y adquiere una dimensión completamente distinta, cuando las entidades vinculadas están repartidas entre distintas jurisdicciones fiscales: ahí, fijar un precio de transferencia artificialmente bajo o alto entre una filial española y una filial en un territorio de baja tributación permite **desplazar artificialmente el beneficio** hacia la jurisdicción donde tributa menos, sin que exista ninguna razón económica real —ningún cambio en dónde se genera efectivamente el valor— que justifique ese desplazamiento.

[Seguro] Este es precisamente el fenómeno que dio origen, en el seno de la OCDE y el G20, al **Proyecto BEPS** (*Base Erosion and Profit Shifting*, erosión de bases imponibles y traslado de beneficios), un paquete de 15 Planes de Acción desarrollado entre 2013 y 2015, entre los que destacan, para lo que aquí interesa, las Acciones 8 a 10 (alinear los resultados en materia de precios de transferencia con la creación de valor real, no con la mera titularidad formal de activos intangibles) y la Acción 13 (documentación de precios de transferencia e **informe país por país**, obligatorio en España para grupos multinacionales con cifra de negocios consolidada igual o superior a 750 millones de euros, que deben desglosar ante la Administración, país por país, dónde declaran beneficios, dónde pagan impuestos y dónde tienen empleados y activos reales —precisamente para detectar desajustes entre ambas magnitudes—).

[Suposición, gancho expositivo] El caso que mejor ilustra, ante cualquier auditorio, por qué este problema dejó de ser una cuestión técnica de especialistas y se convirtió en un debate político y mediático de primer orden es el de **Apple e Irlanda**. En 2016, la Comisión Europea concluyó que Irlanda había concedido a Apple, entre 1991 y 2014, resoluciones fiscales anticipadas (*tax rulings*) que permitían atribuir la práctica totalidad de los beneficios europeos de Apple a una "sede central" sin sustancia económica real —sin empleados, sin oficinas operativas— que, por una peculiaridad del sistema fiscal irlandés de la época, no tributaba en ningún territorio: ni en Irlanda, porque formalmente no residía allí a efectos fiscales, ni en ningún otro lugar. El resultado, documentado por la propia Comisión, fue que en 2014 Apple pagó, sobre sus beneficios europeos, un tipo efectivo cercano al **0,005%**. [Seguro] La Comisión calificó esta ventaja como ayuda de Estado ilegal y ordenó a Irlanda recuperar 13.000 millones de euros más intereses; tras un primer revés en 2020 (el Tribunal General de la UE anuló la decisión, considerando que la Comisión no había probado suficientemente la ventaja selectiva), el **Tribunal de Justicia de la Unión Europea, en su sentencia de 10 de septiembre de 2024 (asunto C-465/20 P)**, resolvió definitivamente el litigio, anuló la sentencia del Tribunal General y confirmó la Decisión original de la Comisión de 2016: Apple debía devolver, y de hecho ya había ingresado en una cuenta de garantía bloqueada desde 2018, los 13.000 millones de euros más unos 1.200 millones de intereses acumulados. [Probable] El caso no se resolvió, técnicamente, por la vía de las normas de precios de transferencia u operaciones vinculadas que hemos visto en este subapartado, sino por la vía del Derecho de la competencia europeo (ayudas de Estado, art. 107 TFUE) — pero el problema de fondo, la asignación artificial de beneficios a una entidad sin sustancia económica real, es exactamente el mismo que persigue corregir el estándar de plena competencia y el conjunto del proyecto BEPS, y por eso el caso funciona como cierre perfecto de este subapartado: muestra, con una cifra y un nombre que cualquier tribunal reconocerá de inmediato, qué es lo que está en juego cuando hablamos de "divergencia por atribución" en el Impuesto sobre Sociedades.

\subsection{Debates abiertos: incidencia económica y doble imposición del dividendo}

### 4.1. Por qué cierro el apartado con dos preguntas y no con una síntesis técnica más

Estos dos debates no son un añadido enciclopédico al final del apartado: son la culminación natural de la pregunta que abría el subapartado 1 —por qué gravar a la persona jurídica—. Si la respuesta fuera enteramente la teoría del beneficio, la incidencia económica del impuesto sería una cuestión secundaria (la sociedad "merece" pagar, con independencia de sobre quién recaiga finalmente la carga) y la doble imposición del dividendo no sería, en rigor, un problema (dos sujetos distintos, dos capacidades económicas distintas, dos impuestos legítimos). Si la respuesta fuera enteramente la teoría del anticipo, ambas cuestiones se vuelven centrales y urgentes: hay que saber quién soporta realmente la carga para valorar si el impuesto cumple su función redistributiva, y hay que corregir la doble imposición porque, de no hacerlo, se estaría gravando dos veces la misma renta del mismo sujeto último. España, como veremos, no ha resuelto esta tensión de forma consistente —trata de forma completamente distinta al socio-entidad y al socio-persona física—, y esa inconsistencia es, en sí misma, el mejor indicador de que el debate sigue genuinamente abierto.

### 4.2. La incidencia económica del Impuesto sobre Sociedades: ¿quién paga, en realidad?

[Seguro] La pregunta de la incidencia —quién soporta finalmente la carga económica de un impuesto, con independencia de quién esté legalmente obligado a ingresarlo— tiene, para el IS, una respuesta de referencia clásica: el modelo de equilibrio general de dos sectores de **Arnold Harberger** (1962, *The Incidence of the Corporation Income Tax*, *Journal of Political Economy*). Harberger modela una economía cerrada con dos sectores —uno corporativo, sujeto al IS, y otro no corporativo, exento— y capital que se mueve libremente entre ambos buscando la máxima rentabilidad después de impuestos. [Probable] Su conclusión, que se convirtió en el consenso dominante durante más de tres décadas, es que, aunque el impuesto se introduce solo en el sector corporativo, el capital huye de ese sector hasta que la rentabilidad después de impuestos se iguala en toda la economía —de modo que, en equilibrio, es **el capital en su conjunto**, no solo el capital invertido en sociedades, el que acaba soportando la carga del impuesto, porque el propio proceso de huida deprime la rentabilidad también en el sector no corporativo al saturarlo de capital desplazado. Es un resultado elegante y, durante mucho tiempo, tranquilizador para quien defendía el IS como instrumento redistributivo: si el capital soporta la carga, y el capital está más concentrado en los deciles altos de renta que el trabajo, el IS es progresivo.

[Seguro] Este consenso se ha visto sustancialmente cuestionado desde que el modelo se extiende a una **economía abierta**, con capital móvil no ya solo entre sectores dentro de un país sino entre países. [Probable] Si el capital puede huir no solo del sector corporativo sino del propio territorio nacional hacia jurisdicciones con menor tributación —exactamente el fenómeno de movilidad internacional del capital que fundamentaba, en el Tema 16, la lógica del modelo Dual Income Tax—, entonces gravar el capital a nivel doméstico deja de deprimir la rentabilidad del capital a escala global (el capital simplemente se reubica) y empieza a recaer, en cambio, sobre el factor que **no puede** moverse: el trabajo, vía menores salarios, o los consumidores, vía precios más altos. Trabajos posteriores en esta línea, como el de William Randolph (2006, *International Burdens of the Corporate Income Tax*, documento de trabajo del Congressional Budget Office estadounidense), estimaron con modelos de economía abierta que el trabajo podría llegar a soportar una parte muy sustancial de la carga del impuesto —en algunas estimaciones, más de la mitad—, un resultado radicalmente distinto del de Harberger y con implicaciones de política completamente opuestas: si es el trabajo quien paga, el IS deja de ser claramente progresivo y su papel como instrumento redistributivo queda en entredicho.

[Probable] Conviene que no zanjes este debate con una respuesta única ante el tribunal, porque no está zanjado en la literatura: la magnitud real de esta traslación depende de parámetros empíricos —la elasticidad de la movilidad internacional del capital, el tamaño de la economía, el grado de competencia en los mercados de bienes y de trabajo— sobre los que no existe consenso cuantitativo firme, y las estimaciones varían enormemente según el modelo y los supuestos empleados. **(OJO: esta incertidumbre sobre quién soporta realmente la carga del IS es, en el fondo, la misma tensión eficiencia-equidad en la imposición del capital que sustenta el debate del Dual Income Tax que desarrollamos en el Tema 16 —la única diferencia es que allí el debate se libraba en sede del contribuyente individual, y aquí se libra en sede de la persona jurídica—, así que si necesitas desarrollar en profundidad los fundamentos de ese debate, la referencia está en el Tema 12, no hace falta repetirlo aquí.)** Lo que sí puedes afirmar con seguridad es la implicación de política que se deriva de cualquiera de las dos posiciones: si el trabajo soporta una parte relevante de la carga del IS, entonces la reducción del tipo general del 30% al 25% que trajo la Ley 27/2014 —enmarcada, recuerda, en la misma doctrina de "ensanchar bases, bajar tipos" que ya vimos en el Tema 16— podría defenderse no solo por razones de competitividad empresarial, sino también como una medida con un componente distributivo mucho menos regresivo de lo que una lectura superficial ("bajar impuestos a las empresas") sugeriría a primera vista.

### 4.3. La doble imposición económica del dividendo: modelos comparados y la asimetría española

[Seguro] El segundo debate es más concreto y tiene una respuesta normativa mucho más precisa en el ordenamiento español, aunque —como verás— profundamente asimétrica. La literatura comparada, sistematizada ya en estudios clásicos como el de Álvaro de la Cueva González-Cotera (2000) para la Comisión Europea, distingue tres grandes modelos de integración entre el IS y el impuesto personal del socio: el **sistema clásico** (sin ninguna corrección; el dividendo tributa íntegramente en ambos niveles, como si fueran dos rentas independientes), la **integración total** (mediante sistemas de imputación con crédito fiscal pleno, que reconstruyen para el socio la ficción de que ha tributado él mismo, directamente, por el beneficio societario, descontando de su cuota el impuesto ya pagado por la sociedad), y la **integración parcial** (exención total o parcial del dividendo en sede del socio, o tipos reducidos específicos para las rentas del capital de origen societario, como hace la base del ahorro española que vimos en el Tema 16).

[Seguro] España ha ensayado, históricamente, más de uno de estos modelos, y aquí está el punto que conviene fijar con precisión: **el modelo aplicado ha sido distinto según si el socio es una persona física o una entidad**, y esa distinción se ha ido ensanchando con el tiempo en vez de reducirse. Hasta la reforma de 2006, el socio persona física corregía la doble imposición mediante un **sistema de imputación en su variante estimativa** (heredado de la Ley 43/1995 y el TRLIS de 2004): el contribuyente elevaba al íntegro el dividendo percibido —sumándole una estimación del impuesto que presumiblemente había pagado la sociedad— y aplicaba después una deducción en cuota por ese importe estimado, reconstruyendo así, de forma aproximada, el impuesto que habría pagado si hubiera obtenido esa renta directamente. La **Ley 35/2006** sustituyó este mecanismo, complejo de calcular y de justificar, por uno mucho más simple pero también mucho más burdo: una **exención de los primeros 1.500 euros anuales** de dividendos percibidos por el contribuyente, con independencia del impuesto efectivamente pagado por la sociedad. Y la **Ley 26/2014** —ya la conoces bien del Tema 16— eliminó incluso esa exención simplificada. Desde 2015, el socio persona física español no dispone de **ningún mecanismo de corrección de la doble imposición del dividendo**: el dividendo tributa íntegramente en sede de la sociedad (al 25% general) y, después, íntegramente en sede del socio, dentro de la base del ahorro del IRPF, a los tipos que ya conoces (19% a 30%). España aplica, para el accionista persona física, un **sistema clásico puro**, mitigado únicamente —y de forma indirecta, no diseñada específicamente para corregir la doble imposición— por el hecho de que la base del ahorro tributa a tipos más bajos que la base general.

[Seguro] Frente a esto, para el socio-entidad, España mantiene —como vimos en el subapartado anterior— un régimen de **exención** que hoy alcanza el 95% del dividendo (art. 21 LIS). [Probable] La asimetría es, por tanto, radical: una sociedad que cobra un dividendo de otra sociedad corrige casi íntegramente la doble imposición; una persona física que cobra exactamente el mismo dividendo no corrige nada en absoluto. Esta asimetría no es casual ni fruto del descuido: responde, en gran medida, a un factor que raramente se menciona fuera de la literatura especializada y que conviene que conozcas, porque explica por qué España —y prácticamente toda la Unión Europea— abandonó los sistemas de imputación para personas físicas en la primera década de este siglo. [Seguro] La **Sentencia del Tribunal de Justicia de las Comunidades Europeas de 7 de septiembre de 2004, asunto C-319/02 (Manninen)**, declaró contrario a la libertad de circulación de capitales (entonces art. 56 TCE) el sistema de imputación finlandés, porque el crédito fiscal por el impuesto societario ya pagado solo se reconocía cuando el dividendo procedía de una sociedad finlandesa, no cuando procedía de una sociedad de otro Estado miembro —una discriminación que penalizaba, de facto, la inversión transfronteriza dentro de la propia Unión—. [Probable] La sentencia puso a los Estados miembros ante un dilema práctico: o extendían el crédito fiscal a los dividendos de sociedades extranjeras —algo administrativamente muy complejo, porque exige conocer con precisión cuánto impuesto societario ha pagado realmente una sociedad sometida a un sistema fiscal distinto—, o abandonaban directamente el sistema de imputación en favor de mecanismos más simples de exención parcial o de tipos reducidos, que no requieren esa trazabilidad transfronteriza. La mayoría de los Estados miembros, España entre ellos, optaron por la segunda vía —lo que explica, en buena medida, por qué la simplificación operada por la Ley 35/2006 en 2007 (la exención de los primeros 1.500 euros) no fue solo una decisión de política fiscal doméstica, sino también una respuesta a una presión de Derecho comunitario que hacía inviable mantener el antiguo sistema de imputación tal y como estaba diseñado.

[Probable] El resultado final, con el que cierro este subapartado, es una fotografía que conviene que sepas describir con precisión ante el tribunal: España integra plenamente al socio-entidad (95% de exención) y no integra en absoluto al socio-persona física (sistema clásico puro), una asimetría que **no tiene fundamento en ningún principio de capacidad económica** —el dividendo es la misma renta, tributando dos veces sobre el mismo beneficio subyacente, con independencia de qué tipo de sujeto lo reciba en última instancia— y que se explica, más bien, por la combinación de dos factores de naturaleza completamente distinta: la necesidad estructural de evitar la imposición en cascada dentro de los grupos societarios (que justifica la exención del art. 21 LIS) y la simplificación administrativa forzada por la jurisprudencia comunitaria tras Manninen (que llevó a renunciar a cualquier corrección para el socio persona física, en lugar de mantener una corrección más compleja pero más precisa).


\subsection{Balance: ¿qué tipo de Impuesto sobre Sociedades tiene España?}

Cierro el apartado con la misma operación de síntesis que hicimos en el Tema 16 para el IRPF, pero aquí sobre dos ejes de caracterización distintos y, en cierto modo, independientes entre sí: el grado de corrección de la doble imposición del dividendo, y el grado de conexión entre contabilidad y fiscalidad. Ninguno de los dos admite una etiqueta única y limpia; en ambos, España ocupa una posición intermedia y, sobre todo en el primero, marcadamente asimétrica.

5.1. Primer eje: un sistema bifronte según quién sea el socio

Ya lo hemos fijado con precisión en el subapartado anterior, así que aquí solo lo sintetizo como conclusión: España no tiene un modelo de integración IS-IRPF/IS, sino dos, aplicados según la naturaleza del perceptor. Frente al socio-entidad, el sistema es de exención (hoy al 95%, art. 21 LIS) — el modelo de integración más generoso de los tres que vimos, y el que menos doble imposición residual deja. Frente al socio-persona física, el sistema es clásico puro, sin ningún mecanismo de corrección desde 2015 — el modelo menos generoso de los tres, el que en principio la literatura comparada considera menos defendible en términos de neutralidad y de capacidad económica. [Probable] Esta bifurcación no responde a ningún principio unificador de diseño tributario, sino a la confluencia de dos lógicas independientes que hemos ido desenredando a lo largo del apartado: la necesidad estructural de evitar la imposición en cascada dentro de las cadenas societarias (que justifica por sí sola la exención del art. 21 LIS, con independencia de cualquier consideración sobre el socio final) y la renuncia forzada, tras la jurisprudencia comunitaria Manninen, a mantener un sistema de imputación más preciso mecanismo para el socio persona física. El resultado, involuntario pero real, es que el sistema español de integración discrimina sistemáticamente al pequeño accionista minoritario, persona física, frente al gran accionista corporativo — exactamente la dirección de desigualdad que un sistema tributario progresivo debería evitar, no reforzar.

5.2. Segundo eje: dependencia contable con corrección extensiva

[Seguro] La literatura de contabilidad comparada, sistematizada en el trabajo clásico de Lamb, Nobes y Roberts (1998, International Variations in the Connections between Tax and Financial Reporting, Accounting and Business Research), clasifica los sistemas fiscales comparados según su grado de conexión entre normativa contable y normativa fiscal, situando en un extremo del espectro a los países de tradición anglosajona —donde contabilidad y fiscalidad son, en la práctica, dos sistemas de determinación de resultado casi completamente independientes, cada uno con su propio conjunto de normas y sin pretensión de partir el uno del otro— y en el extremo opuesto a los países de tradición continental, donde la fiscalidad toma como punto de partida obligado el resultado contable, en aplicación de principios como el alemán Maßgeblichkeitsprinzip ("principio de autoridad" o de vinculación, recogido en el § 5 de la Einkommensteuergesetz alemana), que durante más de un siglo ha exigido que los valores reconocidos en el balance mercantil se trasladen, salvo excepción expresa, al balance fiscal.

[Probable] España, según el análisis específico de Nobes, Oliveras y Puig (2004, The Changing Relationship between Tax and Financial Reporting in Spain, Universitat Pompeu Fabra), se sitúa claramente en el polo continental de este espectro, y así lo confirma el propio art. 10.3 LIS que hemos usado como benchmark de referencia de todo este apartado: el resultado contable es, por mandato legal expreso, el punto de partida obligado de la base imponible. [Probable] Pero conviene matizar esta clasificación con algo que el propio recorrido de este apartado ya te ha mostrado de forma práctica: una cosa es la dependencia formal (partir del resultado contable) y otra distinta es la magnitud real de los ajustes que se practican sobre ese punto de partida. El catálogo que hemos desarrollado en los subapartados 2 y 3 —amortizaciones, provisiones, deterioros, instrumentos financieros, compensación de BINs, exención de participaciones— muestra que, aunque el punto de arranque sea siempre contable, el número y la magnitud económica de las correcciones fiscales que se le aplican son considerables. [Probable] Es, de hecho, el mismo fenómeno que la literatura empírica ha documentado para Alemania —el país con la tradición de conexión más fuerte del entorno comparado—: estudios como el de Zinn (2012) sobre datos reales de declaraciones fiscales alemanas muestran que, pese al Maßgeblichkeitsprinzip, la renta fiscal alemana supera en torno a un 10% a la renta contable declarada a los accionistas, una divergencia nada desdeñable que ha llevado a la propia Alemania a debilitar su principio de conexión mediante la reforma BilMoG de 2009 (Bilanzrechtsmodernisierungsgesetz, Ley de Modernización del Derecho Contable). [Suposición] La lectura más precisa para España, por tanto, no es "sistema de fuerte conexión contable-fiscal" sin más matices, sino "sistema de dependencia formal con corrección extensiva" — mantiene la contabilidad como punto de partida legal y como ancla conceptual del impuesto, pero se aparta de ella en un número suficientemente amplio de partidas como para que el resultado fiscal final pueda diferir sustancialmente del resultado contable que ven los accionistas en la memoria anual.

5.3. Cierre del apartado y tránsito al apartado 3

Con estos dos ejes de caracterización —bifurcación en la corrección de la doble imposición según el tipo de socio, y dependencia contable con corrección extensiva— queda cerrado el "sistema tributario de referencia" real del Impuesto sobre Sociedades español, del mismo modo que en el Tema 16 cerramos la caracterización del IRPF como modelo semi-dual. Y la función de este cierre es exactamente la misma que allí: no es SHS lo que hemos estado midiendo en este apartado —ya lo advertimos desde el primer subapartado, porque aquí no existe un benchmark económico equivalente—, sino la distancia entre el resultado contable y la base imponible fiscalmente exigida, corregida allí donde el legislador ha decidido no confiar plenamente en el juicio contable por razones de verificabilidad, de prevención de la elusión, o de protección recaudatoria. El apartado 3 no vuelve a comparar el IS con ningún ideal teórico externo: parte de este sistema de referencia ya caracterizado —contabilidad corregida, con integración desigual según el tipo de socio— para catalogar, dentro de él, las desviaciones que sí responden a una finalidad deliberada de fomento económico: la reserva de capitalización, la libertad de amortización, las deducciones por I+D+i, por producciones cinematográficas, por creación de empleo, y el resto del catálogo que ya anticipamos al fijar la estructura completa de este tema.



\clearpage
\section{Beneficios fiscales: }
\puntosclave{
    
}

Antes de proponer la estructura, señalo dónde se rompe otra vez la analogía mecánica con el Tema 16, porque si no lo advierto ahora se me va a colar en la redacción: en el IRPF, la regresividad de las reducciones en base nacía de la progresividad de la tarifa (tramos marginales distintos). El IS, en cambio, tiene una tarifa prácticamente proporcional ($25\%$ general, con algún tipo reducido tasado), así que ese argumento no se traslada sin más. Aquí el problema distributivo real de los incentivos fiscales es otro, y mucho más interesante: no depende del tramo de renta del beneficiario, sino de si la empresa tiene cuota suficiente para absorber la deducción. Lo desarrollo en la estructura.

Antes de entrar en el catálogo, fijo el objeto exacto de este apartado, porque ya hemos aprendido, a lo largo del tema, que la frontera entre "estructural" y "beneficio fiscal" no es evidente por sí sola y que equivocarla tiene consecuencias analíticas serias. Este apartado trata exclusivamente de los mecanismos del IS que premian deliberadamente una conducta —invertir, capitalizar, contratar, rodar una película en España— y no de los mecanismos que corrigen una distorsión introducida por el propio diseño del impuesto. Con ese criterio, quedan expresamente fuera del catálogo dos regímenes que podrían parecer, a primera vista, beneficios fiscales de primer orden, y que ya hemos situado en su lugar correcto: la neutralidad fiscal en las reestructuraciones empresariales, que desarrollamos como extensión de la divergencia temporal del apartado 2, y el régimen de consolidación fiscal, que desarrollamos como extensión de la divergencia por atribución del mismo apartado. Si necesitas revisar por qué ambos son estructurales y no gasto fiscal, el criterio está ya razonado allí; no lo repito aquí. Lo que sí trato en este apartado son cuatro cuestiones: la forma jurídica de los incentivos y el problema específico de absorción que plantean en un impuesto de tarifa casi proporcional; el catálogo de incentivos por finalidad; la dispersión foral, con su singular historial de litigio europeo; y el balance de magnitud y eficacia que cierra el tema completo.

\subsection{Forma jurídica de los incentivos: El problema latente}
1.1. Por qué el argumento del Tema 16 no se traslada sin más

[Seguro] En el IRPF, el criterio decisivo para valorar el efecto distributivo de un beneficio fiscal era su ubicación en el esquema de liquidación —reducción en base frente a deducción en cuota—, porque la tarifa progresiva multiplicaba el valor de una reducción en base según el tramo marginal del contribuyente. El IS tiene una tarifa prácticamente proporcional: un tipo general único del 25%, con tipos reducidos tasados para supuestos muy concretos (empresas de nueva creación, entidades de reducida dimensión en algún ejercicio puntual, cooperativas), pero sin la escala de tramos que hacía tan relevante, en el Tema 16, distinguir dónde interviene el beneficio dentro del esquema de liquidación. Si aplicásemos aquí mecánicamente el argumento del Tema 16, concluiríamos —equivocadamente— que la forma jurídica del incentivo carece de relevancia distributiva en el IS. No es así: el problema simplemente se ha desplazado a otro punto de la liquidación.

1.2. El problema real: la absorción de cuota

[Probable] Una deducción en cuota —que es, de hecho, la forma jurídica predominante en el catálogo de incentivos del IS, frente a la reducción en base, mucho más marginal aquí— solo genera ahorro fiscal efectivo si la empresa tiene cuota íntegra positiva suficiente para absorberla. Una empresa en pérdidas, o con una cuota íntegra reducida, puede tener derecho pleno a una deducción y, sin embargo, no beneficiarse de ella en el ejercicio en que la genera, porque no hay nada contra lo que aplicarla. [Probable] Esto tiene una consecuencia distributiva que la literatura de política de innovación ha documentado extensamente y que conviene que conozcas con nombre propio: los incentivos fiscales a la I+D+i, por su propia naturaleza de deducción en cuota, tienden a beneficiar de forma sistemática a las empresas ya consolidadas y rentables —que son, con frecuencia, las que menos necesitan el incentivo para decidir invertir en innovación— y a dejar sin beneficio efectivo a las empresas jóvenes, intensivas en I+D y todavía no rentables —el perfil arquetípico de una start-up tecnológica en sus primeros años, que es exactamente el sujeto que la política de innovación querría incentivar con mayor intensidad—. Bronwyn Hall y John Van Reenen, en su revisión de referencia sobre la eficacia de los incentivos fiscales a la I+D (Hall y Van Reenen, 2000, How Effective Are Fiscal Incentives for R&D? A Review of the Evidence, Research Policy), señalan precisamente este sesgo hacia empresas rentables y consolidadas como uno de los problemas de diseño más persistentes de este tipo de incentivos en el conjunto de países de la OCDE, no solo en España.

1.3. El mecanismo técnico del límite y la respuesta legislativa: la monetización

[Seguro] El art. 39.1 LIS fija el marco general de este problema: el conjunto de deducciones para incentivar determinadas actividades no puede superar, en cada ejercicio, el 25% de la cuota íntegra minorada en las deducciones por doble imposición y las bonificaciones, límite que se eleva al 50% cuando el importe de la deducción por I+D+i del propio ejercicio excede el 10% de esa misma cuota íntegra. El exceso no se pierde —puede aplicarse en ejercicios futuros dentro de los plazos que la propia ley establece—, pero para una empresa con pérdidas recurrentes ese diferimiento puede no llegar nunca a materializarse en un ahorro fiscal real.

[Seguro] El legislador fue consciente de este problema y diseñó, específicamente para la deducción por I+D+i, una solución singular: el art. 39.2 LIS permite optar por excluir esta deducción del límite general —aplicándola incluso hasta dejar la cuota líquida en cero— o, si ni siquiera así hay cuota suficiente para absorberla, solicitar directamente su abono en efectivo a la Administración tributaria a través de la propia autoliquidación del IS —lo que la práctica profesional conoce, con acierto descriptivo, como "monetización de la deducción"—. [Seguro] Ambas vías —exceder el límite general o solicitar el abono— se aplican con una penalización del 20% sobre el importe de la deducción, y están sujetas a límites cuantitativos propios (un millón de euros anuales para innovación tecnológica, tres millones anuales para el conjunto de I+D+i) y a requisitos materiales exigentes: mantener o incrementar la plantilla media durante los veinticuatro meses siguientes, destinar un importe equivalente a la deducción aplicada o abonada a gastos de I+D+i o a inversión en inmovilizado afecto a estas actividades en ese mismo plazo, y disponer de un informe motivado que certifique la calificación de la actividad como investigación, desarrollo o innovación tecnológica. [Probable] Es, en sustancia, la traducción española del mismo problema que resuelven, con mecanismos distintos, otros países de nuestro entorno: Francia, con su Crédit d'Impôt Recherche, permite igualmente que las pequeñas y medianas empresas soliciten el reembolso inmediato del crédito no absorbido por insuficiencia de cuota, precisamente para no penalizar a las empresas jóvenes e intensivas en I+D frente a las ya consolidadas.

[Seguro] Este mecanismo, sin embargo, ha generado un problema de encaje normativo reciente y todavía no resuelto con claridad, que sirve de excelente ejemplo de "problema de implementación" para tu exposición oral: la tributación mínima del 15% que introdujo el art. 30 bis LIS en 2022 —ya la conoces del apartado 1— no aclara expresamente cómo interactúa con la posibilidad de aplicar la deducción de I+D+i sin límite hasta dejar la cuota líquida en cero. Como han señalado análisis especializados recientes (Cuatrecasas, 2022, sobre la convivencia entre ambos regímenes), la opción de solicitar el abono en efectivo por insuficiencia de cuota —el "cash back" del art. 39.2— parece funcionar como excepción al régimen general de la cuota líquida mínima, al constituir precisamente una excepción a la regla de que la cuota líquida no puede ser negativa; pero la posibilidad de aplicar la deducción sin límite hasta reducir la cuota líquida a cero plantea más dudas, porque la norma de 2022 no fue redactada teniendo en cuenta esta interacción específica. [Suposición] Es un buen ejemplo de cómo dos reformas aprobadas con ocho años de diferencia, cada una con una lógica propia y defendible por separado —incentivar la I+D+i con la máxima flexibilidad posible, y garantizar un suelo mínimo de tributación efectiva para las grandes empresas—, pueden generar fricciones de aplicación práctica que ni el legislador ni la doctrina han terminado de resolver con seguridad jurídica plena.

1.4. La lente para leer el catálogo

[Probable] Con esto fijado, la pregunta que debes aplicar a cada incentivo del catálogo que viene a continuación no es "¿en qué tramo marginal beneficia más?" —la pregunta correcta para el IRPF del Tema 16—, sino "¿puede realmente aprovecharlo la empresa a la que se dirige, dada su situación de rentabilidad y de cuota disponible?". Verás que esta pregunta separa con nitidez los incentivos bien diseñados —como la I+D+i, con su válvula de escape de la monetización— de otros, como veremos, que carecen de ese mecanismo corrector y que, por tanto, benefician de forma mucho más desigual según la madurez y la rentabilidad de la empresa que los reclama.


\subsection{Catálogo de incentivos: Objetivo de política económica}
Antes de entrar en el catálogo, conviene fijar el marco teórico que explica por qué el legislador interviene, en general, con incentivos fiscales sobre la actividad empresarial, porque sin ese marco cada figura del catálogo parece una decisión aislada y arbitraria, y no lo es. La justificación económica clásica de cualquier incentivo fiscal a una actividad privada es la existencia de un fallo de mercado: alguna razón por la que el mercado, dejado a su libre funcionamiento, produce menos de esa actividad de lo que sería socialmente óptimo. Verás que esta idea —tomada de la teoría económica del bienestar, no de la doctrina tributaria— reaparece, con variantes distintas, en casi todas las figuras del catálogo: unas veces el fallo de mercado es una externalidad positiva no apropiada por quien invierte (I+D+i); otras veces es una asimetría estructural en el propio diseño del impuesto que distorsiona una decisión de financiación neutral (capitalización empresarial); otras veces es una discriminación de origen social que el mercado laboral no corrige por sí solo (empleo de personas con discapacidad); y otras veces, más discutible doctrinalmente, es simplemente una competencia entre jurisdicciones por una actividad económica móvil (producciones audiovisuales). Organizo el catálogo precisamente en ese orden, de mayor a menor solidez doctrinal del fallo de mercado que se invoca, porque es un criterio más informativo que el orden alfabético o el orden del articulado.

### 2.1. Capitalización empresarial: dos reservas con lógica de fondo distinta
Empiezo por esta figura porque es la que descansa sobre el fallo de mercado más antiguo y mejor documentado de todo el catálogo, y porque, a diferencia de las demás, no depende de ninguna externalidad tecnológica o social: depende de una distorsión que el propio Impuesto sobre Sociedades introduce en la decisión de financiación de cualquier empresa, con independencia de a qué se dedique.

[Seguro] El problema que la reserva de capitalización trata de corregir se conoce en la literatura de fiscalidad internacional como sesgo de la deuda frente al capital (debt-equity bias): en prácticamente todos los sistemas de IS comparados, incluido el español, los intereses pagados por financiación ajena son gasto fiscalmente deducible, mientras que la retribución al capital propio —dividendos, beneficios retenidos— no lo es. [Probable] Esta asimetría, que no responde a ningún principio de capacidad económica sino a una tradición contable e histórica del propio impuesto, empuja a las empresas hacia un nivel de endeudamiento superior al que elegirían en un mundo fiscalmente neutral, con dos consecuencias que la literatura considera indeseables: mayor fragilidad financiera agregada de la economía (empresas más endeudadas son más vulnerables a shocks de tipos de interés o de demanda) y una asignación del capital distorsionada por consideraciones fiscales y no por la rentabilidad real de los proyectos de inversión. [Seguro] Ruud de Mooij, en su trabajo de referencia para el Fondo Monetario Internacional (de Mooij, 2011, Tax Biases to Debt Finance: Assessing the Problem, Finding Solutions), cuantifica este sesgo para un amplio panel de países y lo señala como uno de los problemas estructurales más relevantes —y menos visibles para el público general— de la fiscalidad societaria contemporánea, singularmente agravado tras la crisis financiera de 2008, en la que el sobreendeudamiento corporativo jugó un papel causal ampliamente documentado.

[Seguro] Dos son las grandes familias de solución que la doctrina y la práctica comparada han ensayado frente a este problema: limitar más severamente la deducción de intereses (lo que ya vimos, de pasada, que hace la Ley 27/2014 desde su origen, endureciendo la deducibilidad de gastos financieros como parte del ensanchamiento de bases de 2015), o bien —la vía que sigue España con la reserva de capitalización— reconocer una deducción "espejo" sobre el capital propio, para equiparar su tratamiento fiscal al de la deuda. Esta segunda vía tiene nombre técnico propio en la literatura: Allowance for Corporate Equity (ACE), propuesta seminal del IFS Capital Taxes Group británico (1991) y ensayada, con diseños distintos, en países como Bélgica (la llamada Notional Interest Deduction, vigente desde 2006 aunque progresivamente recortada desde 2018 por su elevado coste fiscal) o Italia. [Seguro] La Unión Europea ha llevado esta misma lógica al plano normativo comunitario con la propuesta de Directiva DEBRA (Debt-Equity Bias Reduction Allowance), presentada por la Comisión Europea el 11 de mayo de 2022, que combina una deducción sobre el incremento de fondos propios con una restricción adicional a la deducibilidad de intereses —el mismo doble mecanismo, en esencia, que el legislador español ya aplica de forma separada desde 2015—. [Probable] La doctrina española más reciente ha analizado con detalle la relación entre ambos instrumentos: Cencerrado Millán (2025, DEBRA. Estado de la cuestión y análisis crítico, Revista Española de Derecho Financiero) examina si la reserva de capitalización española satisface ya, totalmente o en parte, las exigencias de la futura Directiva europea, en un contexto en el que la propia tramitación de DEBRA se ha ralentizado sensiblemente en el Consejo por la falta de unanimidad entre Estados miembros.

Contenido y evolución material

Ya desarrollamos con detalle, en la primera versión de este subapartado, la mecánica y la escalada normativa reciente del porcentaje aplicable (10% en su origen bajo la Ley 27/2014; 15% desde 2024; entre el 20% y el 30% desde 2025 según el incremento de plantilla vinculado) y del plazo de mantenimiento (de cinco a tres años). No lo repito; recupero únicamente el dato de la reserva de nivelación como segunda figura del mismo bloque, exclusiva de empresas de reducida dimensión y orientada, no a la capitalización, sino a anticipar la compensación de pérdidas futuras dentro de un horizonte de cinco años.

Implicaciones

[Probable] La experiencia comparada con los sistemas de tipo ACE es, en el mejor de los casos, mixta, y conviene que lo sepas para no presentar la reserva de capitalización como una solución sin coste. El caso belga es el más estudiado: la deducción por intereses nocionales generó un coste recaudatorio muy superior al inicialmente previsto —en gran parte porque grandes grupos multinacionales reestructuraron artificialmente sus balances para maximizar el capital computable, sin que ello implicara ninguna inversión productiva adicional real—, lo que llevó al Gobierno belga a recortar drásticamente el mecanismo desde 2018. [Suposición] El diseño español, con su requisito de mantenimiento efectivo durante varios años y su reciente vinculación expresa al incremento de plantilla —que dificulta mucho más la ingeniería puramente contable que aprovechó el caso belga—, parece haber aprendido de ese precedente, aunque todavía no existe, hasta donde alcanza la literatura de evaluación disponible en España, un estudio de impacto público y sistemático equivalente al que ya vimos para otros incentivos en el Tema 16.

Agrupo aquí dos figuras que a menudo se presentan juntas por compartir mecánica formal —ambas reducen la base imponible mediante la dotación de una reserva indisponible— pero que responden a finalidades de política económica distintas y conviene no confundir.

[Seguro] La **reserva de capitalización** (art. 25 LIS) premia la autofinanciación empresarial: permite reducir la base imponible en un porcentaje del incremento de los fondos propios del ejercicio —es decir, del beneficio que la empresa decide no repartir como dividendo y retener como reserva—, siempre que ese incremento se mantenga durante un plazo determinado y se dote una reserva específica, indisponible durante ese plazo, por el importe de la reducción. [Seguro] Es una figura que ha experimentado una escalada normativa notable en muy pocos años: nació con un porcentaje del 10% y un plazo de mantenimiento de cinco años; la reforma con efectos desde el 1 de enero de 2024 elevó el porcentaje al 15% y redujo el plazo a tres años; y la reforma con efectos desde el 1 de enero de 2025 lo eleva de nuevo, hasta el **20%** con carácter general, con una escala adicional que premia expresamente la creación de empleo —23% si la plantilla media crece entre el 2% y el 5%, 26,5% si crece entre el 5% y el 10%, y hasta el **30%** si el incremento de plantilla supera el 10%—, vinculando así, de forma cada vez más explícita, el incentivo a la capitalización con el incentivo a la contratación. [Seguro] El requisito de mantenimiento no es meramente formal: una resolución reciente del **Tribunal Económico-Administrativo Central, de 20 de octubre de 2025**, ha confirmado que el incremento de fondos propios debe mantenerse **año a año** durante todo el plazo, y no solo verificarse al final del periodo, rechazando la interpretación más laxa que defendía el contribuyente afectado —un buen ejemplo, para tu exposición oral, de cómo un incentivo aparentemente sencillo genera litigiosidad real cuando la empresa atraviesa dificultades económicas sobrevenidas dentro del plazo de mantenimiento.

[Seguro] La **reserva de nivelación** (art. 105 LIS), en cambio, está reservada exclusivamente a las **empresas de reducida dimensión** (cifra de negocios inferior a 10 millones de euros) y responde a una lógica distinta: no premia la capitalización, sino que permite **anticipar** el efecto de la compensación de bases imponibles negativas, reduciendo la base imponible hasta un 10% (con el límite de un millón de euros), en previsión de pérdidas que la empresa pueda generar en los cinco años siguientes. Si esas pérdidas se producen dentro de ese plazo, la reserva las compensa automáticamente; si no se producen, el importe reducido se reincorpora a la base imponible al final del quinto año, con los correspondientes intereses de demora. [Probable] Aplicando la lente del subapartado 1: mientras que la reserva de capitalización beneficia, casi por definición, a empresas ya rentables (solo hay incremento de fondos propios que premiar si ha habido beneficio), la reserva de nivelación está diseñada específicamente para empresas de reducida dimensión con expectativas de volatilidad en sus resultados —un perfil de riesgo mucho más próximo al de la pyme que necesita un colchón fiscal frente a la incertidumbre, y no al de la gran empresa consolidada.

### 2.2. I+D+i: generosidad nominal y el problema ya anticipado

Introducción

Si la capitalización empresarial descansaba en una distorsión introducida por el propio impuesto, la I+D+i descansa en el fallo de mercado más citado y mejor fundamentado de toda la teoría económica de la intervención pública en la innovación, y conviene que sepas nombrarlo con precisión, porque es prácticamente inevitable que aparezca en cualquier pregunta sobre la justificación de este incentivo.

Finalidad del legislador y debate económico de fondo

[Seguro] El fundamento es la existencia de externalidades de conocimiento (knowledge spillovers): quien invierte en investigación genera, además del beneficio privado que espera obtener, un beneficio social adicional que no puede apropiarse en su totalidad —otras empresas aprenden de sus avances, se benefician de la difusión del conocimiento generado, o directamente lo imitan una vez es público—. Este argumento, formulado en su versión más influyente por Kenneth Arrow (1962, Economic Welfare and the Allocation of Resources for Invention, en el volumen The Rate and Direction of Inventive Activity, National Bureau of Economic Research), es la justificación seminal por la que la teoría económica del bienestar concluye que, sin intervención pública, el mercado invertirá menos en investigación de lo que sería socialmente óptimo: la empresa privada calcula su decisión de inversión sobre el retorno privado esperado, no sobre el retorno social total, que incluye esos efectos de difusión. [Probable] Un incentivo fiscal a la I+D es, bajo esta lógica, un intento de acercar el retorno privado percibido por la empresa al retorno social real de la actividad, corrigiendo así una infrainversión estructural que el mercado, por sí solo, nunca resolvería.

[Probable] La discusión doctrinal no está, sin embargo, en si existe la externalidad —eso es consenso prácticamente universal en la literatura desde Arrow—, sino en cómo intervenir de forma más eficaz: mediante deducción fiscal (como hace España) o mediante subvención directa competitiva (como hacen, en paralelo, el CDTI español o los programas marco europeos de innovación). Bronwyn Hall, en trabajos posteriores a la revisión ya citada en el subapartado 1, ha señalado que ambos instrumentos tienen fortalezas complementarias: la deducción fiscal es más neutral —no exige que un funcionario público seleccione ex ante qué proyectos merecen financiación, y por tanto no sufre el riesgo de "elegir ganadores" equivocadamente— pero también es menos capaz de dirigir la inversión hacia la investigación más básica y de mayor riesgo, que es precisamente donde la externalidad de conocimiento es mayor y donde el mercado privado menos invertiría sin ayuda; la subvención competitiva puede dirigirse mejor a ese tipo de investigación, pero a costa de una capacidad de selección administrativa que no siempre acierta. [Suposición] España, al mantener ambos instrumentos en paralelo —deducción fiscal generalista vía art. 35 LIS y subvenciones competitivas vía CDTI—, ha optado, de facto, por no elegir entre ambas lógicas, sino por combinarlas, una solución pragmática que algunos análisis de política de innovación consideran razonable y otros consideran una duplicidad de instrumentos con coste administrativo evitable.

Contenido y evolución material

Ya fijamos con detalle, en la versión previa de este subapartado, los porcentajes vigentes (25%/42% para I+D, 12% para innovación tecnológica, adicionales del 17% y el 8%) y el mecanismo de absorción y monetización del art. 39.2 LIS desarrollado en el subapartado 1. No lo repito.

Implicaciones

[Probable] La pregunta de eficacia que plantea Hall y Van Reenen (2000) —input additionality frente a output additionality: si el incentivo genera gasto en I+D adicional al que ya se habría producido, o si simplemente subvenciona investigación que la empresa habría acometido de todos modos— sigue sin respuesta cerrada para el caso español específico. [Suposición] Lo que sí puede afirmarse con cierta seguridad, retomando el dato ya citado del índice B de la OCDE, es que la generosidad nominal española es de las más altas del entorno comparado; la pregunta abierta, que retomaremos en el balance final del apartado, es si esa generosidad nominal se traduce en un incremento proporcional del esfuerzo real en I+D+i del tejido empresarial español —que sigue estando, en términos de porcentaje del PIB, por debajo de la media de la Unión Europea pese a contar con uno de los regímenes fiscales más generosos—, una aparente paradoja que solo un estudio de evaluación de impacto propiamente dicho podría resolver con rigor.

---------- Anterior:

[Seguro] Ya fijamos en el subapartado 1 la mecánica de absorción de esta deducción; toca ahora su contenido sustantivo. El art. 35 LIS distingue investigación y desarrollo (I+D) de innovación tecnológica (i), con porcentajes de deducción notablemente más generosos para la primera: **25%** de los gastos del periodo (elevado al **42%** sobre el exceso respecto de la media de gastos de los dos años anteriores, un diseño que premia expresamente el **crecimiento** del esfuerzo investigador de la empresa y no solo su nivel absoluto), más una deducción adicional del **17%** sobre los gastos de personal investigador cualificado dedicado en exclusiva, y del **8%** sobre inversiones en inmovilizado afecto a la actividad investigadora. La innovación tecnológica, considerada por la ley una actividad de menor riesgo y más próxima al desarrollo comercial que a la investigación pura, recibe un porcentaje único y más modesto del **12%**. [Probable] España aparece de forma recurrente, en los índices comparados de generosidad fiscal a la I+D+i que publica la OCDE (el llamado *B-index*, que mide el coste fiscal marginal de invertir un euro adicional en I+D), entre los países con el régimen nominal más generoso de todo el entorno comparado —un dato que retomaremos con más detalle en el balance final del apartado, porque generosidad nominal y eficacia real, como ya anticipamos con Hall y Van Reenen, no son necesariamente la misma cosa.

### 2.3. Producciones audiovisuales: de incentivo cultural a plató internacional

Introducción

Esta figura es, de las cuatro que desarrollo con detalle, la que descansa sobre un fundamento doctrinal más disputado, porque en realidad combina dos justificaciones de naturaleza completamente distinta que conviene separar antes de valorarlas.

Finalidad del legislador y debate económico de fondo

[Probable] La primera justificación es de política cultural: los bienes culturales —y la producción audiovisual en particular— tienen, según la tradición de la economía de la cultura sistematizada por autores como David Throsby, un valor que excede el precio que el mercado está dispuesto a pagar por ellos en el momento del consumo —lo que se conoce como valor de existencia u opción, análogo al que se invoca para justificar la protección del patrimonio histórico—, además de una función de cohesión e identidad cultural que un mercado puramente privado tendería a infraproducir, especialmente frente a la competencia de producciones extranjeras con mayores economías de escala. Es, en el fondo, el mismo tipo de argumento de bien público que ya vimos en el Tema 16 a propósito del mecenazgo (Feldstein, 1975) trasladado ahora a la producción cultural directa.

[Probable] La segunda justificación, la que explica específicamente la deducción por producciones extranjeras rodadas en España (art. 36.2 LIS) y el régimen reforzado canario, es de naturaleza radicalmente distinta: es un argumento de atracción de actividad económica móvil, del mismo tipo, en su lógica de fondo, que la competencia fiscal internacional por la localización de sedes empresariales que ya analizamos en el apartado 2 a propósito de la incidencia del IS y del proyecto BEPS. [Probable] Y aquí es donde la doctrina se vuelve más crítica: mientras que la primera justificación (bien cultural) tiene un fundamento de fallo de mercado defendible en términos de teoría del bienestar, la segunda es, en esencia, una subvención competitiva a la deslocalización de producciones, un fenómeno muy estudiado en Estados Unidos —donde la práctica totalidad de los estados compiten entre sí ofreciendo créditos fiscales cada vez más generosos para atraer rodajes ("runaway production")—, y donde la literatura de evaluación de políticas públicas, con estudios como los sucesivos informes del Tax Foundation y de varias oficinas presupuestarias estatales, ha sido sistemáticamente escéptica sobre el retorno económico neto de estos incentivos: los puestos de trabajo generados son, en gran medida, temporales y especializados, gran parte del gasto de la producción se destina a personal técnico que ya trabajaba en el sector de todos modos, y el efecto multiplicador sobre la economía local suele ser sensiblemente inferior al que las administraciones que conceden el incentivo anuncian públicamente al aprobarlo.

Contenido y evolución material

Recupero, sin repetir en detalle, lo ya desarrollado sobre los porcentajes generales y el régimen reforzado de Canarias, con el ejemplo ya dado de grandes producciones internacionales atraídas por el incentivo canario.

Implicaciones

[Suposición] Aplicando el marco crítico anterior al caso español, conviene que distingas, en tu exposición oral, entre el juicio que mereces dar a la deducción por producciones españolas (más defendible, en la medida en que responde al argumento de bien cultural) y el juicio que mereces dar a la deducción por producciones extranjeras y al régimen reforzado canario (más próximo, en su lógica económica, a una subvención competitiva cuyo balance fiscal neto —ingresos tributarios y económicos generados frente al coste fiscal directo del incentivo— rara vez se ha sometido en España a una evaluación de impacto tan rigurosa como la que sí existe, por ejemplo, para el caso estadounidense).

----------- Anterior:

[Seguro] El art. 36 LIS regula tres deducciones relacionadas con la industria audiovisual y de las artes escénicas: por inversión en producciones cinematográficas y series audiovisuales españolas (art. 36.1), por producciones extranjeras rodadas en España (art. 36.2), y por espectáculos en vivo de artes escénicas y musicales (art. 36.3). [Probable] La deducción por producciones españolas se articula, en su estructura general, sobre un porcentaje reforzado en el primer tramo de coste de la producción y reducido en el exceso —siguiendo el mismo patrón de "más generosidad en el primer euro" que vimos en el Tema 16 a propósito del mecenazgo del IRPF—, con límites máximos de deducción por producción que se han ido ampliando en sucesivas reformas recientes, en un contexto de competencia internacional cada vez más intensa por atraer rodajes. [Suposición, gancho expositivo] El caso de **Canarias** es el más llamativo de todo el catálogo del IS, porque el Régimen Económico y Fiscal (REF) del archipiélago —una singularidad territorial con fundamento constitucional propio, distinta de la dispersión foral que trataremos en el siguiente subapartado— eleva estos porcentajes hasta niveles muy superiores a los del territorio general: la propia Agencia Tributaria cifra la deducción por producciones cinematográficas en Canarias en el 54% (45% en el segundo tramo), frente a porcentajes sensiblemente más bajos en el resto del territorio, con un incentivo adicional todavía mayor para las islas de La Palma, La Gomera y El Hierro. Esta ventaja fiscal explica, en gran medida, por qué Canarias se ha convertido en plató habitual de superproducciones internacionales en la última década —desde secuencias de sagas como *Star Wars* o *Juego de Tronos* hasta numerosas producciones de plataformas de streaming—, y es un ejemplo perfecto de cómo un incentivo fiscal puede generar un ecosistema industrial completo (equipos técnicos locales, infraestructura de rodaje, formación especializada) allí donde, sin el incentivo, la actividad simplemente no se localizaría en el territorio.

### 2.4. Empleo de personas con discapacidad: Discriminación de mercado

Introducción

Cierro el catálogo desarrollado con detalle con la figura que, aunque de menor peso presupuestario que las anteriores, tiene detrás la justificación doctrinal más directamente vinculada a un principio constitucional expreso, y no solo a un argumento de eficiencia económica.

Finalidad del legislador y debate económico de fondo

[Seguro] El fundamento no es aquí, en sentido estricto, un fallo de mercado en el sentido técnico de externalidad no apropiada, sino la corrección de una discriminación estadística documentada en el mercado de trabajo: la literatura de economía laboral sobre discapacidad y empleo —con trabajos de referencia como los de Kruse y Schur en el contexto estadounidense, o los sucesivos informes de la OIT en el ámbito comparado— muestra de forma consistente que las personas con discapacidad se enfrentan a tasas de desempleo y de inactividad sistemáticamente superiores a las de la población general, en un grado que no se explica solo por diferencias reales de productividad, sino también por prejuicios de contratación y por costes de ajuste (adaptación de puestos de trabajo) que el empleador internaliza de forma asimétrica. [Probable] La deducción del art. 38 LIS se complementa, y no sustituye, con la obligación de cuota de reserva de empleo del 2% para empresas de más de 50 trabajadores que impone la legislación general sobre derechos de las personas con discapacidad (hoy, Real Decreto Legislativo 1/2013), configurando así un sistema de doble instrumento —obligación cuantitativa más incentivo fiscal complementario— que es habitual en la política comparada de empleo de colectivos vulnerables.

Contenido y evolución material

Recupero, sin repetir, los importes ya fijados (9.000/12.000 euros en origen, 11.700/15.600 euros desde noviembre de 2018).

Implicaciones

[Suposición] A diferencia de la I+D+i, no existe en España, hasta donde alcanza la literatura de evaluación disponible, un estudio de impacto específico y reciente sobre la eficacia de esta deducción en términos de creación neta de empleo para personas con discapacidad frente al escenario sin incentivo —una laguna de evaluación que contrasta con la relativa abundancia de literatura sobre la eficacia de la cuota de reserva del 2%, mucho más estudiada por su carácter obligatorio y por la existencia de sanciones asociadas a su incumplimiento.

------------- Anterior:

[Seguro] El art. 38 LIS reconoce una deducción por cada persona/año de incremento del promedio de plantilla de trabajadores con discapacidad, contratados por tiempo indefinido: 9.000 euros por cada persona/año de incremento del promedio de plantilla con discapacidad en un grado igual o superior al 33% e inferior al 65%, y 12.000 euros si el grado de discapacidad es igual o superior al 65%, importes elevados desde noviembre de 2018 a 11.700 y 15.600 euros respectivamente. [Probable] Es una deducción de importe fijo por trabajador, no proporcional a ningún gasto ni a ninguna inversión, lo que la distingue estructuralmente de casi todo el resto del catálogo —no premia un esfuerzo económico variable, sino directamente una decisión de contratación—, y aplicando de nuevo la lente del subapartado 1, es también un ejemplo de incentivo cuyo aprovechamiento efectivo depende enteramente de que la empresa contratante tenga cuota íntegra suficiente, sin que exista aquí, a diferencia de la I+D+i, ningún mecanismo equivalente de monetización para las empresas sin cuota positiva.



### 2.5. Mecenazgo: la vertiente societaria del problema del "polizón"

La justificación de fondo de esta figura no es distinta, en su núcleo teórico, de la que ya fijamos para la I+D+i, aunque el fallo de mercado que corrige tiene otro nombre en la literatura: la financiación privada de bienes públicos o cuasi-públicos —cultura, investigación básica, asistencia social— tiende a la infraprovisión estructural porque cualquier potencial donante puede beneficiarse del bien generado por las donaciones de terceros sin necesidad de contribuir él mismo, el clásico problema del **polizón** (*free-rider problem*) que Paul Samuelson formalizó en su artículo seminal de 1954 sobre la teoría pura del gasto público (*The Pure Theory of Public Expenditure*, *Review of Economics and Statistics*). El incentivo fiscal al mecenazgo es, en esa lectura, un intento de corregir esa infraprovisión abaratando artificialmente el coste de donar. Pero cuando el donante no es una persona física movida por un impulso altruista, como en el IRPF del Tema 16, sino una **sociedad mercantil**, aparece un matiz doctrinal adicional que no existía en aquel contexto y que conviene que domines: buena parte de la literatura de responsabilidad social corporativa —con trabajos ya clásicos sobre el vínculo entre filantropía empresarial y reputación de marca— sostiene que la donación corporativa no siempre responde a un impulso genuinamente altruista, sino que funciona, en gran medida, como un **gasto de naturaleza cuasi-publicitaria**, orientado a mejorar la reputación de la empresa frente a consumidores, inversores y reguladores, y que por tanto la empresa habría incurrido en ese gasto de todos modos, con o sin deducción fiscal, por su propio interés comercial. Si esta segunda lectura es correcta —y no hay consenso empírico cerrado sobre en qué proporción lo es—, entonces la deducción del art. 20 Ley 49/2002 estaría subvencionando, al menos parcialmente, un gasto que ya seria rentable en términos puramente privados, lo que reintroduce, con otro disfraz, el mismo problema de **peso muerto** que ya identificamos en el Tema 16 a propósito de Feldstein (1975) y de la previsión social: no toda donación adicional generada por el incentivo es donación que sin él no se habría producido. En cuanto a su contenido material, el art. 20 de la Ley 49/2002 reconoce una deducción del 35% de la base de la donación, elevable al 40% cuando la empresa ha donado a la misma entidad, por importe igual o superior, en cada uno de los dos ejercicios anteriores —la misma lógica de "prima a la fidelidad" que ya conocías del ámbito del IRPF—, con un límite del 10% de la base imponible del periodo (traslado del exceso a los diez ejercicios siguientes) y con la posibilidad, si el Gobierno declara la actividad donada como prioritaria en la Ley de Presupuestos, de elevar esos porcentajes hasta el 40%/45% y ese límite hasta el 15%. El **Real Decreto-ley 6/2023, de 19 de diciembre**, aprobado dentro de la ejecución del Plan de Recuperación, Transformación y Resiliencia, ha ampliado además el propio concepto de donación deducible —incorporando, desde el 1 de enero de 2024, la cesión gratuita y temporal del uso de bienes muebles o inmuebles como nueva modalidad con derecho a deducción—, y la evolución normativa más reciente apunta a una nueva elevación de los porcentajes generales hacia el entorno del 40-50%, en la línea de refuerzo del mecenazgo que ya veníamos observando desde 2021. Aplicando la lente del subapartado 1, conviene no perder de vista que esta deducción está, como casi todo el catálogo, sujeta al problema de absorción sobre cuota: solo la empresa con beneficio suficiente puede aprovecharla plenamente, lo que en la práctica concentra el mecenazgo empresarial español en el segmento de grandes empresas rentables, y explica en parte por qué la filantropía corporativa española sigue estando, en comparación con el entorno anglosajón, menos desarrollada entre la pequeña y mediana empresa.

------------ Anterior:
No desarrollo aquí de nuevo el régimen de la Ley 49/2002, que ya tratamos con detalle en el Tema 16 desde la perspectiva del IRPF; me limito a señalar su proyección en el IS, regulada en los mismos arts. 20 y 21 de esa ley: las entidades que realizan donativos a entidades sin fines lucrativos acogidas al régimen especial de mecenazgo pueden deducir en su cuota íntegra un porcentaje del importe donado —con una estructura de tramos y de "prima a la fidelidad" por donaciones recurrentes análoga, aunque con porcentajes propios, a la que ya conoces del ámbito del IRPF—, sujeta también, como el resto del catálogo de este subapartado, al límite general de absorción sobre cuota que fijamos en el subapartado 1.

### 2.6. El régimen de empresas de reducida dimensión: entre la restricción de crédito y el efecto umbral

La justificación teórica de este bloque no procede de la literatura de innovación ni de la de bienes públicos, sino de la literatura sobre **imperfecciones del mercado de crédito**: el trabajo fundacional de Joseph Stiglitz y Andrew Weiss (1981, *Credit Rationing in Markets with Imperfect Information*, *American Economic Review*) mostró que, en presencia de información asimétrica entre prestamista y prestatario, el mercado de crédito no se ajusta necesariamente vía precio —el banco no puede simplemente subir el tipo de interés hasta encontrar el equilibrio, porque hacerlo expulsaría selectivamente a los prestatarios más solventes y dejaría una cartera de riesgo peor—, lo que genera **racionamiento de crédito** incluso para proyectos con rentabilidad esperada positiva. Las empresas pequeñas, sin historial crediticio extenso ni activos suficientes para ofrecer como garantía, son las que sufren este racionamiento con mayor intensidad, lo que constituye el fundamento económico clásico —más allá de cualquier consideración de mero tamaño o de peso en el empleo agregado— para that el Estado intervenga con instrumentos que mejoren su posición financiera, entre ellos la libertad de amortización (que, al adelantar la deducción fiscal, funciona de facto como un préstamo sin interés del Estado a la empresa, mejorando su liquidez inmediata sin necesidad de acudir al mercado de crédito) y la reserva de nivelación ya desarrollada en el subapartado 2.1. El contenido material de este bloque ha experimentado, sin embargo, una transformación reciente de primer orden que conviene incorporar con el detalle que merece, porque es de rabiosa actualidad para la campaña del impuesto de estos mismos ejercicios: durante años, tras la Ley 27/2014, el régimen de reducida dimensión perdió su antiguo tipo de gravamen reducido y quedó reducido, en la práctica, a la libertad de amortización y la reserva de nivelación, mientras el tipo general se aplicaba de forma casi universal. La **Ley 7/2024, de 20 de diciembre** —la misma norma que trasladó a España el Pilar 2 de la OCDE y que ya citamos en el apartado 1—, ha revertido parcialmente esta tendencia, introduciendo una nueva escala de tipos reducidos y progresivos, distinta según el tamaño exacto de la entidad: las **microempresas** (cifra de negocios inferior a un millón de euros) tributan, para los periodos iniciados en 2025, al 21% por los primeros 50.000 euros de base imponible y al 22% por el resto, con una reducción escalonada hasta el 17%/20% en 2027 conforme a la nueva Disposición Transitoria 44ª LIS; las **entidades de reducida dimensión propiamente dichas** (entre uno y diez millones de cifra de negocios) tributan al 24% en 2025, con descensos progresivos anuales hasta el 20% en 2029; y se mantiene, al margen de ambas categorías, el tipo del 15% para entidades de nueva creación en su primer periodo con base imponible positiva y en el siguiente. Esta reintroducción de un tipo reducido específico por tramos de tamaño reabre, sin embargo, un debate doctrinal que la literatura de economía laboral y de organización industrial ha documentado con datos muy sólidos fuera del ámbito estrictamente tributario, pero que se aplica con la misma lógica a cualquier umbral fiscal por tamaño: los llamados **efectos umbral** o de *bunching*. Luis Garicano, Claire Lelarge y John Van Reenen (2016, *Firm Size Distortions and the Productivity Distribution: Evidence from France*, *Review of Economic Studies*) mostraron, con datos del universo de empresas francesas, que cuando una regulación —en su caso, laboral, pero el mecanismo es idéntico para cualquier umbral fiscal— empieza a aplicarse de forma discontinua a partir de un tamaño concreto de empresa, un número significativo de empresas elige deliberadamente **no crecer** más allá de ese umbral para conservar el régimen favorable, generando una distorsión visible en la propia distribución estadística del tamaño empresarial —una acumulación anómala de empresas justo por debajo del umbral— y una pérdida de eficiencia agregada, porque esas empresas terminan operando a una escala inferior a la que sería óptima en ausencia de la regulación. [Suposición] Aplicado al caso español, los umbrales de uno y diez millones de cifra de negocios que ahora delimitan microempresa, entidad de reducida dimensión y régimen general son, precisamente, el tipo de frontera discreta que esta literatura identifica como generadora de incentivos a la "atrofia deliberada": es razonable esperar, aunque no exista todavía un estudio de *bunching* específico y publicado para el caso español tras esta reforma de 2024, que una parte del tejido empresarial español ajuste su facturación declarada, o incluso su estrategia de crecimiento real, para permanecer justo por debajo de esos umbrales y conservar el tipo reducido —el mismo tipo de coste de eficiencia, documentado empíricamente en Francia, que cualquier diseño de política fiscal por tramos de tamaño debe sopesar frente al beneficio, también real, de aliviar la restricción de crédito de las empresas genuinamente pequeñas.

--------- Anterior:
[Probable] Cierro el catálogo con una figura transversal que atraviesa varias de las anteriores: el régimen de empresas de reducida dimensión (cifra de negocios inferior a 10 millones de euros) no es, en rigor, un incentivo único sino un conjunto de ventajas —la ya vista reserva de nivelación, exclusiva de este colectivo, y regímenes de **libertad de amortización** para determinados elementos nuevos del inmovilizado material afectos a la actividad, vinculados a menudo a la creación o mantenimiento de empleo—. Es aquí donde reaparece, con nombre propio, la figura que ya identificamos como gasto fiscal genuino al distinguirla de las tablas ordinarias de amortización en el apartado 2: la libertad de amortización no corrige ninguna estimación contable, sino que permite adelantar íntegramente la deducción del coste de un activo muy por delante de su depreciación económica real, precisamente para incentivar la inversión productiva en el segmento empresarial que, en España, representa la inmensa mayoría del tejido productivo en número de entidades.


\subsection{Dispersión territorial: Competencia fiscal interna}

A diferencia de la dispersión autonómica del IRPF que analizamos en el Tema 16 —amplísima en número de medidas pero de perfil jurídico relativamente bajo—, la dispersión foral en el IS descansa sobre un fundamento constitucional propio y arrastra un historial de litigio europeo de primer nivel, que conviene entender desde su raíz teórica antes de entrar en el episodio que mejor la ilustra. La justificación económica clásica de permitir que entidades subestatales fijen su propia fiscalidad procede de Charles Tiebout (1956, A Pure Theory of Local Expenditures, Journal of Political Economy): si contribuyentes y empresas pueden elegir libremente su jurisdicción entre varias que ofrecen combinaciones distintas de impuestos y servicios públicos, esa movilidad —"votar con los pies"— revela preferencias reales y permite una asignación más eficiente de recursos que la que lograría una autoridad fiscal centralizada y uniforme. Es el mismo argumento de fondo que ya vimos justificar la dualización de la base del ahorro en el Tema 16 y la competencia internacional por bases imponibles móviles en el apartado 2 de este tema, trasladado ahora a la competencia entre territorios de un mismo Estado. Pero tiene, desde los años ochenta, un contrapunto igualmente clásico: el modelo de George Zodrow y Peter Mieszkowski (1986, Pigou, Tiebout, Property Taxation, and the Underprovision of Local Public Goods, Journal of Urban Economics) muestra que, cuando el factor móvil es el capital y no la mera elección residencial de un consumidor, la competencia entre jurisdicciones no conduce a un óptimo de eficiencia, sino a una infraprovisión sistemática de bienes públicos: cada territorio, temeroso de perder base imponible móvil, recorta su tipo por debajo del nivel que sus propios ciudadanos preferirían, en una carrera a la baja de la que ningún territorio puede salir unilateralmente sin perder competitividad frente a los demás — el mismo problema, en clave interna, que el proyecto BEPS trata de corregir en clave internacional.

España convive con ambas lógicas mediante dos vías distintas: la cesión limitada de capacidad normativa a las Comunidades Autónomas de régimen común —marginal en el IS, como ya señalamos al fijar el marco jurídico— y una autonomía normativa plena en los territorios de Concierto Económico (País Vasco) y Convenio Económico (Navarra), amparada en la Disposición Adicional Primera de la Constitución, que ampara y respeta los derechos históricos de los territorios forales. Esta autonomía permite a las Diputaciones Forales vascas y a la Comunidad Foral de Navarra fijar tipos y deducciones propias sustancialmente distintas de las del territorio común —tipos generales forales históricamente inferiores al estatal, y un catálogo de deducciones propias, especialmente en I+D+i, más generoso que el estatal—, lo que ha generado, desde los años noventa, sucesivas reacciones de las Comunidades limítrofes —principalmente La Rioja y Castilla y León— ante los tribunales españoles y la Comisión Europea, al considerar que esa ventaja distorsiona la competencia empresarial dentro del propio mercado español.

3.2. El test de autonomía: la doctrina Azores y su aplicación al caso vasco

El marco doctrinal para resolver esta tensión lo fijó el Tribunal de Justicia en el llamado caso Azores (Sentencia de 6 de septiembre de 2006, asunto C-88/03, República Portuguesa contra Comisión), relativo a una reducción del IS en las islas Azores portuguesas. El Tribunal estableció que, cuando una medida fiscal procede de una entidad infraestatal dotada de verdadera autonomía respecto del gobierno central, el marco de referencia para apreciar su selectividad —y, por tanto, si constituye ayuda de Estado prohibida— no tiene por qué ser el territorio completo del Estado miembro, sino que puede limitarse al ámbito geográfico de la propia entidad, siempre que concurran tres condiciones acumulativas: autonomía institucional (estatuto político y administrativo propio), autonomía de procedimiento (el gobierno central no puede intervenir en la adopción de la norma territorial) y autonomía económica (la propia entidad asume las consecuencias financieras de su menor recaudación, sin compensación estatal).

Esta doctrina se aplicó directamente al caso español en la Sentencia de 11 de septiembre de 2008 (asuntos acumulados C-428/06 a C-434/06, "caso País Vasco"), donde el Tribunal —siguiendo las conclusiones de la Abogada General Juliane Kokott— concluyó que los tres Territorios Históricos vascos satisfacen las tres condiciones, gracias muy singularmente al propio mecanismo del cupo: al no recibir compensación estatal por la menor recaudación que generen sus propias decisiones normativas —el llamado "riesgo unilateral"—, las Diputaciones Forales asumen íntegramente el coste fiscal de sus políticas, cumpliendo la exigencia de autonomía económica que el caso Azores había dejado como la más exigente en la práctica comparada. La consecuencia es, en principio, favorable a la autonomía fiscal vasca en general: un tipo foral más bajo que el estatal no constituye, por sí mismo, ayuda de Estado selectiva, porque el marco de referencia correcto es el propio territorio vasco, no España en su conjunto.

3.3. El límite del test: las vacaciones fiscales vascas y la selectividad material

Esta doctrina de 2008, favorable en abstracto, no resuelve —y aquí está el matiz que distingue a un candidato que entiende el problema de uno que lo simplifica— el problema específico que habían planteado, casi una década antes, las llamadas "vacaciones fiscales vascas": un conjunto de normas forales de los años noventa que establecían, para empresas de nueva creación en los tres Territorios Históricos, una exención total del IS durante varios ejercicios, junto con un crédito fiscal del 45% sobre determinadas inversiones. La Comisión Europea, en Decisiones de 11 de julio de 2001 (entre ellas la Decisión 2002/820/CE), calificó estas medidas como ayudas de Estado incompatibles, calificación confirmada en sucesivas resoluciones: la Sentencia de 14 de diciembre de 2006 (asuntos acumulados C-485/03 a C-490/03) declaró que España había incumplido su obligación de recuperación; la Sentencia de 9 de junio de 2011 (asuntos acumulados C-465/09 P a C-470/09 P, caso Confebask) rechazó en casación los últimos recursos de las Diputaciones Forales; y ante la persistente demora en ejecutar la recuperación, el Tribunal, en Sentencia de 13 de mayo de 2014, impuso a España una multa de 30 millones de euros, repercutida internamente vía cupo sobre las propias Diputaciones.

[Probable] La razón por la que estas medidas concretas no se benefician de la doctrina favorable de 2008, aunque su litigio se prolongara más allá de esa fecha, es una distinción técnica de primer orden: el test Azores resuelve el problema de la selectividad regional —si comparar el tipo vasco con el estatal es, en sí mismo, relevante—, pero las vacaciones fiscales planteaban, además, un problema de selectividad material dentro del propio territorio vasco: la exención no beneficiaba a todas las empresas vascas por igual, sino específicamente a las de nueva creación, lo que constituye selectividad —y por tanto ayuda de Estado potencialmente ilegal si no se notifica a la Comisión conforme al art. 108.3 TFUE— con independencia de cuál sea el marco de referencia territorial correcto. Es el mismo tipo de distinción que ya trazamos en el apartado 2 a propósito de Apple-Irlanda: no basta con que la autoridad tributaria tenga competencia legítima para fijar su propia fiscalidad —el Tribunal nunca cuestionó la legitimidad institucional del Concierto Económico—; el problema surge cuando, dentro de esa competencia legítima, se concede una ventaja que favorece selectivamente a unas empresas frente a otras sin justificación en la lógica interna del propio sistema tributario.

El episodio tuvo un epílogo interno de interés: las empresas vascas obligadas a devolver las ayudas reclamaron indemnización patrimonial a las Diputaciones Forales, alegando buena fe. El Tribunal Supremo, en Sentencia 1361/2018, de 5 de septiembre, rechazó estas reclamaciones, señalando —con cita de doctrina consolidada del propio Tribunal de Justicia— que el principio de confianza legítima no ampara a los beneficiarios de ayudas declaradas ilegales, sino a sus competidores, que se encontraban en desventaja frente a quienes sí las disfrutaron: un cierre, dos décadas después del origen normativo del problema, que traza con más precisión que cualquier explicación abstracta el límite real de la autonomía fiscal foral.


\subsection{Balance: magnitud, competencia fiscal internacional y evaluación de eficacia}

### 4.1. La magnitud cuantitativa: un peso menor, pero no marginal, dentro del sistema

Cierro el apartado, y con él el tema completo, retomando el mismo ejercicio de cuantificación que hicimos en el Tema 16 para el IRPF, porque solo con cifras concretas se entiende la escala real de lo que hemos catalogado. [Seguro] Según el Presupuesto de Beneficios Fiscales del Estado, en el proyecto de Presupuestos para 2023 el conjunto de beneficios fiscales del Impuesto sobre Sociedades se cifró en 5.697 millones de euros, equivalentes al 12,6\% del total de beneficios fiscales presupuestados (45.269 millones de euros en conjunto), con un descenso del 2,8\% respecto del ejercicio anterior explicado, expresamente, por la caída de los beneficios asociados a las sociedades de inversión y por el efecto de contención que empezaba a producir la tributación mínima del 15\% para grandes empresas que vimos en el apartado 1. 

[Probable] Es una cifra sensiblemente inferior, en términos absolutos, a la que vimos para el IRPF en el Tema 16 (en torno a 11.000 millones de euros en los mismos ejercicios), lo que no debe leerse como que el IS tenga menos "gastos fiscales" en sentido relativo, sino como el reflejo de que el propio impuesto recauda menos —en torno a 28.500 millones de euros previstos para 2023, frente a los más de 113.000 millones del IRPF— y de que buena parte de sus beneficios formales, como ya aprendimos al principio de este apartado, no se contabilizan como tales en la Memoria por ser estructurales y no gasto fiscal en sentido estricto. [Seguro] La serie histórica muestra, además, una volatilidad singularmente marcada en el IS, mucho mayor que la que vimos para el IRPF: en el Presupuesto de 2013, en plena crisis, los beneficios fiscales del IS cayeron un 16,6% en un solo ejercicio, hasta 3.084,8 millones de euros, por la conjunción de tres factores que ya conoces bien de este tema —la limitación de la deducción por reinversión de beneficios extraordinarios, la supresión de ciertos regímenes de libertad de amortización, y la caída generalizada de los beneficios empresariales sobre los que se calculan buena parte de estos incentivos—. [Probable] Esta mayor volatilidad, frente a la relativa estabilidad del gasto fiscal en el IRPF, es coherente con algo que hemos visto repetirse a lo largo de todo el apartado: buena parte del catálogo del IS —capitalización, I+D+i, cine— está diseñado como porcentaje de una magnitud variable (incremento de fondos propios, gasto en innovación, inversión en producción), no como importe fijo, de modo que su coste recaudatorio fluctúa mecánicamente con el propio ciclo económico, amplificando en las fases de crisis la sensación —y en parte la realidad— de que el Estado "recauda más" simplemente porque las empresas invierten menos.

### 4.2. La competencia fiscal internacional: de la carrera de tipos a la carrera de créditos

Ya situamos en el apartado 1 la respuesta normativa más reciente a la competencia fiscal internacional sobre el IS —el Pilar 2 de la OCDE/G20 y su transposición española mediante la Ley 7/2024—; cierro ahora el argumento con la perspectiva de conjunto que ese dato aislado no permitía ver por sí solo. [Seguro] La caída de los tipos nominales del IS en el mundo desde los años ochenta es uno de los fenómenos más documentados de la fiscalidad comparada: el tipo medio de gravamen del IS en los países de la OCDE ha descendido de forma prácticamente ininterrumpida durante más de cuatro décadas, desde niveles superiores al 45% a comienzos de los años ochenta hasta situarse, en los años previos a la introducción del Pilar 2, en el entorno del 21-23% de media —una caída que la literatura de competencia fiscal, siguiendo la lógica de Zodrow y Mieszkowski que ya desarrollamos en el subapartado 3, atribuye precisamente a la movilidad creciente del capital y a la presión competitiva entre jurisdicciones por atraer bases imponibles empresariales—. [Probable] El Pilar 2 nace explícitamente para poner un suelo a esta dinámica, mediante un tipo mínimo efectivo del 15% aplicable a los grandes grupos multinacionales con independencia de dónde localicen sus beneficios. Pero conviene que cierres este bloque con una nota de cautela que la propia literatura de fiscalidad internacional más reciente ha empezado a señalar, y que constituye el debate genuinamente abierto y de mayor actualidad en esta materia: las reglas técnicas del Pilar 2 (las llamadas reglas GloBE, *Global Anti-Base Erosion*) tratan de forma distinta los incentivos fiscales según su diseño técnico, excluyendo del cómputo de la tributación efectiva mínima a los llamados **créditos fiscales reembolsables cualificados** (*qualified refundable tax credits*) —aquellos que, como la monetización de la deducción de I+D+i española que analizamos en el subapartado 1, se abonan en efectivo al contribuyente con independencia de su cuota—, mientras que sí computan como reducción de la tributación efectiva a los incentivos articulados como simple reducción de tipo o de base. [Suposición] Especialistas en fiscalidad internacional del entorno de centros como el Oxford University Centre for Business Taxation, dirigido durante años por Michael Devereux, han advertido de que esta asimetría técnica puede generar, paradójicamente, un efecto no deseado por los propios diseñadores del Pilar 2: en lugar de eliminar la competencia fiscal internacional, podría limitarse a **desplazarla** —de la tradicional carrera a la baja de tipos nominales, ya cerrada por el suelo del 15%, hacia una nueva carrera de créditos fiscales reembolsables cada vez más generosos, precisamente porque estos escapan al cómputo del tipo mínimo global—. [Probable] Si esta lectura es correcta, el incentivo español a la I+D+i que hemos analizado en detalle en el subapartado 1 —con su mecanismo de monetización ya diseñado, años antes del Pilar 2, como un crédito reembolsable— podría convertirse, de forma no del todo buscada, en un instrumento particularmente bien posicionado para competir internacionalmente por inversión en el nuevo escenario post-Pilar 2, un ángulo que conecta, de forma inesperada, el debate más técnico de todo el apartado 1 con el incentivo más concreto de todo el apartado 3.

### 4.3. La evaluación de eficacia: el caso concreto de la I+D+i, cerrando el círculo

Cierro el tema retomando, con datos concretos y no solo con la literatura teórica de Hall y Van Reenen ya citada en el subapartado 1, la pregunta de si el incentivo más generoso del catálogo español cumple realmente su función. [Seguro] La propia AIReF, dentro de la misma Fase II del Spending Review de 2020 que ya conoces del Tema 16 —el estudio que evaluó la tributación conjunta y la reducción por previsión social del IRPF—, evaluó también la deducción por I+D+i del IS, con un retraso derivado de las dificultades de acceso a los datos durante el estado de alarma de la pandemia. [Seguro] La conclusión de la AIReF es matizada, y conviene que la sepas transmitir con esa misma matización ante el tribunal: el incentivo **sí** fomenta la inversión empresarial en I+D+i —no es, por tanto, un caso comparable al fracaso que la propia AIReF había identificado para la previsión social del IRPF—, pero existe una **"amplia distancia" entre su eficacia potencial y su eficacia real**, constatando además que España sigue presentando un nivel de inversión en I+D+i por debajo de la media de los países de su entorno, y que ese nivel de inversión cayó de forma singular durante el periodo de recesión 2008-2015, justo cuando la teoría de Arrow que citamos en el subapartado 2 sugeriría que la externalidad de conocimiento —y por tanto la necesidad de sostener el incentivo— era más urgente, no menos. [Seguro] Entre sus recomendaciones concretas, la AIReF propone permitir solicitar el abono de la deducción reembolsable en el **mismo ejercicio** en que se genera el derecho —hoy el mecanismo del art. 39.2 LIS que analizamos en el subapartado 1 exige, en la práctica, esperar a la insuficiencia de cuota constatada en la autoliquidación, con el consiguiente desfase temporal para la empresa—, eliminar algunos de los requisitos formales que dificultan el acceso al mecanismo de monetización, reforzar las obligaciones de información sobre la composición real del gasto en I+D+i de las empresas beneficiarias, y, muy significativamente, **analizar el incentivo fiscal de forma conjunta con las políticas de subvención directa** gestionadas por el CDTI —la misma tensión entre deducción fiscal y subvención competitiva que planteamos, en términos puramente teóricos, al presentar el debate de Hall en el subapartado 2, y que la AIReF confirma, con datos españoles reales, como un problema de coordinación institucional no resuelto y no solo como una discusión académica abstracta.

### 4.4. Cierre del tema

Con este balance queda completo el desarrollo del Tema 17. El recorrido ha replicado, con las adaptaciones sustantivas que fuimos identificando en cada punto de ruptura respecto del Tema 16, la misma arquitectura de fondo: un apartado 1 que fija el marco jurídico-histórico y el esquema de liquidación, con el ajuste extracontable como bisagra hacia el siguiente bloque; un apartado 2 que confronta ese esquema con su propio benchmark —no un ideal económico externo como SHS, sino la norma contable corregida, y una pregunta previa sobre la propia justificación del impuesto que el IRPF nunca tuvo que plantearse— para identificar sus divergencias estructurales, cerrando con la caracterización de un IS semi-dependiente de la contabilidad y bifronte en su corrección de la doble imposición del dividendo; y un apartado 3 que, partiendo de ese sistema de referencia ya caracterizado, cataloga las desviaciones deliberadas del legislador, desde el sesgo estructural hacia la deuda hasta la competencia fiscal territorial, cerrando con el mismo ejercicio de evaluación empírica —AIReF, IEF— que ya aplicamos al IRPF en el Tema 16, confirmando que ambos temas, pese a tratar impuestos de naturaleza jurídica distinta, comparten el mismo método de análisis y, en última instancia, la misma exigencia: no basta con describir el impuesto tal y como está escrito en la ley, hay que poder valorarlo con el mismo rigor con que se le exige rendir cuentas a cualquier otra política pública.













































































\clearpage
\section{Conclusión} 


\subsection{Recapitulación}


\subsection{Extensiones y relación con otras partes del Temario}



\subsection{Opinión}

\subsubsection{Idea final (Salida o cierra)}

\clearpage
\pagenumbering{gobble}
\MyBibliography

Ley 50/1977, de 14 de noviembre, sobre Medidas Urgentes de Reforma Fiscal.
Ley 61/1978, de 27 de diciembre, del Impuesto sobre Sociedades.
Ley 19/1989, de 25 de julio, de reforma parcial y adaptación de la legislación mercantil a las Directivas de la CEE en materia de sociedades.
Ley 43/1995, de 27 de diciembre, del Impuesto sobre Sociedades.
Real Decreto Legislativo 4/2004, de 5 de marzo, Texto Refundido de la Ley del Impuesto sobre Sociedades.
Ley 27/2014, de 27 de noviembre, del Impuesto sobre Sociedades.
Ley 22/2021, de 28 de diciembre, de Presupuestos Generales del Estado para 2022 (art. 30 bis LIS, tributación mínima doméstica).
SANZ GADEA, E. (2003): "Historia del Impuesto sobre Sociedades. Perspectivas de evolución", en *Manual del Impuesto sobre Sociedades*, Instituto de Estudios Fiscales, Madrid.
Sentencia de la Audiencia Nacional de 7 de marzo de 2002 (sujeción de las confesiones religiosas al IS bajo la Ley 61/1978).
Ley 27/2014, arts. 2, 4, 7, 8, 27, 30 bis, 40 y 124 (ya citada en su conjunto; ahora en su contenido específico de marco jurídico).
Real Decreto 634/2015, de 10 de julio, por el que se aprueba el Reglamento del Impuesto sobre Sociedades.
Real Decreto 1514/2007, de 16 de noviembre, por el que se aprueba el Plan General de Contabilidad.
Ley 58/2003, General Tributaria (ya citada; ahora en su función de norma supletoria del IS, arts. 23, 35.2.c y 122).
Ley 27/2014, art. 10.3 (principio de conexión contable-fiscal) y art. 26 (compensación de bases imponibles negativas).
Ley 43/1995, de 27 de diciembre, Exposición de Motivos (función de retención en la fuente del IS respecto de inversores extranjeros).
MUSGRAVE, R. A. y MUSGRAVE, P. B.: Public Finance in Theory and Practice, McGraw-Hill (ediciones sucesivas).
Joint Committee on Taxation (2019): "Estimates of Federal Tax Expenditures for Fiscal Years 2019-2023", JCX-55-19 (tratamiento de las compensaciones de pérdidas como "normal income tax law").
GUTMAN, H. (2020): "Net Operating Loss Carrybacks Are Not Tax Subsidies", carta al editor, Tax Notes Federal, 22 de junio de 2020.
Ley 27/2014, arts. 12, 13, 15.k) y 26 (ya citada; ahora en su contenido específico sobre amortizaciones, deterioros y compensación de BINs).
Ley 16/2013, de 29 de octubre, por la que se establecen determinadas medidas en materia de fiscalidad medioambiental y se adoptan otras medidas tributarias y financieras (no deducibilidad del deterioro de participaciones).
Real Decreto-ley 3/2016, de 2 de diciembre, por el que se adoptan medidas en el ámbito tributario dirigidas a la consolidación de las finanzas públicas.
STC 11/2024, de 18 de enero (BOE de 20 de febrero de 2024), cuestión de inconstitucionalidad núm. 2577-2023.
STC 78/2020 (precedente sobre el Real Decreto-ley 2/2016 y los pagos fraccionados del IS).
Constitución Española de 1978, art. 86.1 (límites materiales del decreto-ley).
Real Decreto 1514/2007, Plan General de Contabilidad (ya citado; ahora en su Norma de Valoración 2ª sobre amortización).
Ley 27/2014, arts. 18 y 21 (ya citada; ahora en su contenido específico sobre operaciones vinculadas y exención de participaciones).
Ley 36/2006, de 29 de noviembre, de medidas para la prevención del fraude fiscal (origen de la obligación de documentación de operaciones vinculadas).
Real Decreto 1793/2008, de 3 de noviembre, de desarrollo reglamentario en materia de documentación de operaciones vinculadas.
Real Decreto-ley 6/2010, de 9 de abril, de medidas para el impulso de la recuperación económica y el empleo (simplificación documental para pymes).
Ley 11/2020, de 30 de diciembre, de Presupuestos Generales del Estado para 2021 (limitación al 95% de la exención del art. 21 LIS).
Directiva 2011/96/UE del Consejo, de 30 de noviembre de 2011 (Directiva matriz-filial).
OCDE (2013-2015): Proyecto BEPS (*Base Erosion and Profit Shifting*), Planes de Acción 8-10 y 13.
OCDE: *Transfer Pricing Guidelines for Multinational Enterprises and Tax Administrations* (ediciones sucesivas).
Sentencia del Tribunal de Justicia de la Unión Europea de 10 de septiembre de 2024, asunto C-465/20 P, *Comisión/Irlanda y otros* (caso Apple).
Decisión de la Comisión Europea de 30 de agosto de 2016, SA.38373, sobre la ayuda estatal concedida por Irlanda a Apple.
Ley 35/2006 (ya citada; ahora en su redacción original de 2006 sobre exención de 1.500 euros en dividendos, y su posterior eliminación por la Ley 26/2014, ya citada en el Tema 16).
Real Decreto Legislativo 4/2004 (TRLIS, ya citado; ahora en sus arts. 23 —antiguo sistema de imputación IRPF—, 30 y 32 —deducciones por doble imposición interna e internacional—).
HARBERGER, A. C. (1962): "The Incidence of the Corporation Income Tax", *Journal of Political Economy*, vol. 70.
RANDOLPH, W. C. (2006): "International Burdens of the Corporate Income Tax", Congressional Budget Office Working Paper 2006-09.
DE LA CUEVA GONZÁLEZ-COTERA, Á. (2000): "La eliminación de la doble imposición económica de dividendos en la Unión Europea", *Estudios Financieros*, CEF, núm. 206.
STJCE de 7 de septiembre de 2004, asunto C-319/02, *Manninen*.
LAMB, M., NOBES, C. y ROBERTS, A. (1998): "International Variations in the Connections between Tax and Financial Reporting", Accounting and Business Research, vol. 28, núm. 3.
NOBES, C., OLIVERAS, E. y PUIG, X. (2004): "The Changing Relationship between Tax and Financial Reporting in Spain", Economics Working Papers 782, Universitat Pompeu Fabra.
ZINN, B. (2012): "Tax Accounting in Germany: Empirical Evidence on the Relationship between Financial and Tax Accounting and Options for Reform", Universität Mannheim / ZEW Discussion Paper 12-051.
BLAKE, J., FORTES, H., AMAT, O. y AKERFELDT, K. (2002): "The Relationship between Tax Regulations and Financial Accounting: A Comparison of Germany, Sweden, Spain and the United Kingdom", Journal of Applied Accounting Research, vol. 4, núm. 1.
§ 5 Abs. 1 EStG (Einkommensteuergesetz alemana, Maßgeblichkeitsprinzip); Bilanzrechtsmodernisierungsgesetz (BilMoG), 2009.
Ley 27/2014, arts. 76 a 89 (Capítulo VII, régimen de neutralidad fiscal en reestructuraciones).
Directiva 2009/133/CE del Consejo, de 19 de octubre de 2009 (Directiva de Fusiones), y su antecedente, la Directiva 90/434/CEE.
Ley 29/1991, de 16 de diciembre, de adecuación de determinados conceptos impositivos a las Directivas y Reglamentos de las Comunidades Europeas (primera transposición española de la Directiva de Fusiones).
Sentencia del Tribunal de Justicia de las Comunidades Europeas de 17 de julio de 1997, asunto C-28/95, Leur-Bloem.
Ley 27/2014, arts. 55 a 75 (Capítulo VI del Título VII, régimen de consolidación fiscal).
Real Decreto-ley 15/1977, de 25 de febrero, y Real Decreto 1414/1977, de 17 de junio (origen del régimen de declaración consolidada).
Ley 18/1982, de 26 de mayo, de modificación del concepto de sociedad dominante a efectos de la tributación consolidada de los grupos de sociedades.
Ley 27/2014, arts. 29, 35, 39.1 y 39.2 (ya citada en su conjunto; ahora en su contenido específico sobre límites y monetización de deducciones).
HALL, B. H. y VAN REENEN, J. (2000): "How Effective Are Fiscal Incentives for R&D? A Review of the Evidence", Research Policy, vol. 29, núms. 4-5.
Cuatrecasas (2022): "Monetización de la deducción I+D+i en el Impuesto sobre Sociedades", Legal Flash Fiscalidad.
Ley 27/2014, arts. 25, 35, 36, 38 y 105 (ya citada; ahora en su contenido específico sobre el catálogo de incentivos).
Resolución del Tribunal Económico-Administrativo Central de 20 de octubre de 2025 (mantenimiento anual del incremento de fondos propios en la reserva de capitalización).
Ley 19/1994, de 6 de julio, de modificación del Régimen Económico y Fiscal de Canarias (marco general del REF, con desarrollo reglamentario posterior de los incentivos a la inversión).
OCDE: *B-index* de generosidad fiscal a la I+D+i, *R&D Tax Incentives* (informes comparados periódicos).
ARROW, K. J. (1962): "Economic Welfare and the Allocation of Resources for Invention", en NELSON, R. (ed.), The Rate and Direction of Inventive Activity: Economic and Social Factors, National Bureau of Economic Research / Princeton University Press.
DE MOOIJ, R. A. (2011): "Tax Biases to Debt Finance: Assessing the Problem, Finding Solutions", IMF Staff Discussion Note SDN/11/11.
IFS Capital Taxes Group (1991): Equity for Companies: A Corporation Tax for the 1990s, Institute for Fiscal Studies, Londres (formulación original del Allowance for Corporate Equity).
Comisión Europea (2022): Propuesta de Directiva DEBRA (Debt-Equity Bias Reduction Allowance), COM(2022) 216 final, 11 de mayo de 2022.
CENCERRADO MILLÁN, E. (2025): "DEBRA. Estado de la cuestión y análisis crítico", Revista Española de Derecho Financiero, núm. 205.
THROSBY, D. (2001): Economics and Culture, Cambridge University Press.
Real Decreto Legislativo 1/2013, de 29 de noviembre, Texto Refundido de la Ley General de derechos de las personas con discapacidad y de su inclusión social (cuota de reserva de empleo del 2%).
- SAMUELSON, P. A. (1954): "The Pure Theory of Public Expenditure", *Review of Economics and Statistics*, vol. 36, núm. 4.
- Ley 49/2002, art. 20 (ya citada; ahora en su contenido específico sobre la deducción del donante persona jurídica).
- Real Decreto-ley 6/2023, de 19 de diciembre, por el que se aprueban medidas urgentes para la ejecución del Plan de Recuperación, Transformación y Resiliencia en materia de servicio público de justicia, función pública, régimen local y mecenazgo.
- STIGLITZ, J. E. y WEISS, A. (1981): "Credit Rationing in Markets with Imperfect Information", *American Economic Review*, vol. 71, núm. 3.
- Ley 7/2014, art. 29 y Disposición Transitoria 44ª LIS, en la redacción dada por la Ley 7/2024, de 20 de diciembre (nueva escala de tipos reducidos para microempresas y entidades de reducida dimensión).
- GARICANO, L., LELARGE, C. y VAN REENEN, J. (2016): "Firm Size Distortions and the Productivity Distribution: Evidence from France", *Review of Economic Studies*, vol. 83, núm. 1.
Disposición Adicional Primera de la Constitución Española de 1978.
Ley 12/2002 (Concierto Económico con el País Vasco) y Ley 28/1990 (Convenio Económico con Navarra), ya citadas en el Tema 16, ahora en su proyección sobre el IS.
TIEBOUT, C. M. (1956): "A Pure Theory of Local Expenditures", Journal of Political Economy, vol. 64, núm. 5.
ZODROW, G. R. y MIESZKOWSKI, P. (1986): "Pigou, Tiebout, Property Taxation, and the Underprovision of Local Public Goods", Journal of Urban Economics, vol. 19, núm. 3.
Sentencia del Tribunal de Justicia de 6 de septiembre de 2006, asunto C-88/03, República Portuguesa/Comisión (caso Azores).
Sentencia del Tribunal de Justicia de 11 de septiembre de 2008, asuntos acumulados C-428/06 a C-434/06 (caso País Vasco).
Decisión 2002/820/CE de la Comisión, de 11 de julio de 2001, y decisiones conexas (vacaciones fiscales vascas).
Sentencia del Tribunal de Justicia de 14 de diciembre de 2006, asuntos acumulados C-485/03 a C-490/03.
Sentencia del Tribunal de Justicia de 9 de junio de 2011, asuntos acumulados C-465/09 P a C-470/09 P (Confebask).
Sentencia del Tribunal de Justicia de 13 de mayo de 2014 (multa a España por demora en la recuperación de ayudas).
Sentencia del Tribunal Supremo 1361/2018, de 5 de septiembre (Sala de lo Contencioso-Administrativo).
- Ministerio de Hacienda: Presupuesto de Beneficios Fiscales, Proyectos de Ley de Presupuestos Generales del Estado 2013, 2022 y 2023 (ya citado de forma genérica en el Tema 16; ahora con cifras específicas del IS).
- AIReF (2020): evaluación de la deducción por I+D+i, dentro del estudio de *Spending Review, Fase II: Beneficios Fiscales* (ya citado en el Tema 16 a propósito de otras figuras del IRPF).
- OCDE: series históricas de tipos nominales del Impuesto sobre Sociedades en los países miembros (*Corporate Tax Statistics*, informes periódicos).
- DEVEREUX, M. P. (Oxford University Centre for Business Taxation): trabajos sobre el diseño de las reglas GloBE del Pilar 2 y el tratamiento de los créditos fiscales reembolsables cualificados.


\MyBibItem{DAWID y GATTI}{Chapter 2 - Agent-Based Macroeconomics}{2018}{https://doi.org/10.1016/bs.hescom.2018.02.006}


% ANEXOS
\clearpage










\end{document}